% !TeX root = ../Book5.tex
% 纲要标题：# 第一编　有限群的表示与特征标
% 纲要位置：工作记录/计划纲要.md，第 3719 行。


\BookPart{有限群的表示与特征标}{b5:part:01}
\BookPartGuide
本编研究有限群在线性空间上的作用。群可以有许多互不同构的表示，一个表示也可能遗漏群的部分结构；问题是怎样找到不可约成分、比较不同表示，并从它们的数值约束反过来认识群。我们从具体矩阵开始，逐步建立特征标、诱导、双陪集和整诱导理论，最后用这些工具证明一个群论中的结构定理。

\ProseParagraphBreak
\chapref{b5:ch:01}给出表示与群代数模的对应，构造正则、置换、对偶和张量表示，证明 Maschke 定理\footnote{马施克（Heinrich Maschke，1853-1908），德国数学家，研究有限群与微分几何；其完全可约性定理是群表示论的基本结果。}和 Schur 引理\footnote{舒尔（Issai Schur，1875-1941），德籍数学家，研究群表示、矩阵理论与数论。}。本章允许一般底域，并分别说明“特征不整除群阶”对完全可约性的作用，以及“底域代数闭”对不可约表示自同态环的影响。\chapref{b5:ch:02}转到复数域，从平均投影建立正交关系，计算重数、中心幂等元和小群的特征标表。

\ProseParagraphBreak
\chapref{b5:ch:03}把子群表示诱导到整个群，给出张量和函数两种模型，并用 Frobenius 互反计算重数。诱导的维数、作用矩阵、特征标公式和传递性都有具体表达式，读者可在 \(S_3\) 的二阶、三阶子群上逐项检验。\chapref{b5:ch:04}继续分析诱导后限制的双陪集分块，再研究不可约表示限制到正规子群时的共轭轨道和重数；射影表示用来记录延拓时出现的标量障碍。

\ProseParagraphBreak
\chapref{b5:ch:05}包含两项不同的整性问题。Artin 定理\footnote{阿丁（Emil Artin，1898-1962），奥地利裔数学家，研究类域论、代数和表示论。}允许有理系数，Brauer 定理\footnote{布劳尔（Richard Brauer，1901-1977），德裔美国数学家，创立模表示论，并研究群的特征标。}进一步用初等子群消去整体分母。后者的证明先把诱导像组织成表示环的理想，再逐个素数进行检测；所需代数整数与有限环扩张的论证都在本章准备。随后证明 Burnside 的 \(p^aq^b\) 阶群可解性定理\footnote{伯恩赛德（William Burnside，1852-1927），英国数学家，研究有限群及其表示，著有《Theory of Groups of Finite Order》。}，明确区分它与轨道计数公式、不可约矩阵代数定理。

\begin{center}
\begin{tikzcd}[column sep=large,row sep=large]
\begin{gathered}\text{第1章}\\\text{表示与分解}\end{gathered}\arrow[r]
&\begin{gathered}\text{第2章}\\\text{特征标与正交}\end{gathered}\arrow[r]\arrow[d]
&\begin{gathered}\text{第3章}\\\text{诱导与互反}\end{gathered}\arrow[d]\\
&\begin{gathered}\text{第5章}\\\text{整诱导与群论}\end{gathered}
&\begin{gathered}\text{第4章}\\\text{Mackey 与 Clifford}\end{gathered}\arrow[l]
\end{tikzcd}
\end{center}

图中箭头串起表示的分解、特征标、诱导与双陪集分析；第5章把这些工具用于特征标的整数表达和群的可解性问题。

\ProseParagraphBreak
有限群表示还通向不变量理论：\chapref{b5:ch:22}的第22.1—22.4节计算有限群不变环及其 Hilbert 级数\footnote{希尔伯特（David Hilbert，1862-1943），德国数学家，证明不变量有限生成定理，并对数学基础和数论作出重要贡献。}，第22.4节的 Gröbner 消去可与第一卷附录 F 对照，第22.5节则把\chapref{b5:ch:07}的 Schur--Weyl 对偶用于矩阵不变量。回到本编，\secref{b5:sec:0504}通过表示环的逐素数检测证明整诱导，\secref{b5:sec:0505}把特征标整性用于 Burnside 可解性定理。

\ProseParagraphBreak
所需群论为子群、共轭、商群、Sylow 定理\footnote{西罗（Peter Ludwig Mejdell Sylow，1832-1918），挪威数学家，研究有限群的素数幂阶子群及其存在和共轭性质。}与有限群作用；线性代数包括迹、特征值、直和与张量的基本运算。所需表示论构造和证明在本卷展开。先计算习题中的具体例子，再读一般证明，有助于理解条件的用途；遇到特征、底域或左右作用的改变，应重新核对公式。

\ProseParagraphBreak
初读可配合 Etingof 等人的《Introduction to Representation Theory》\cite[Chapters 4--5]{EtingofEtAl2011RepresentationTheory}。本编使用普通复特征标建立 Brauer 整诱导理论。正特征的 Brauer 特征标在第8章定义，它用于研究模表示的合成因子，与整诱导定理讨论的问题不同。

% !TeX root = ../../Book5.tex
% 纲要标题：## 第1章　表示、群代数与完全可约性
% 纲要位置：工作记录/计划纲要.md，第 3721 行。


\chapter{表示、群代数与完全可约性}
\label{b5:ch:01}
\BookChapterNumber{b5:ch:01}

\begin{chapterabstract}
表示把群元素变成可逆线性变换，群代数则把同一作用写成模结构。本章从置换、循环群及张量构造认识这些对象，以平均投影证明 Maschke 定理，并给出模特征下的反例。Schur 引理随后确定不可约表示之间的同态，区分代数闭域与一般域，建立等型分解及重数的唯一性。
\end{chapterabstract}

\begin{learninggoals}
\begin{itemize}
\item 在群作用、矩阵及群代数模三种描述之间转换，正确核验同态与同构。
\item 构造置换、正则、一维、对偶、张量及外幂表示。
\item 用平均投影分解表示，指出群阶可逆性在证明中的作用。
\item 将平均法用于各次齐次多项式，计算 Reynolds 算子及简单不变环。
\item 用 Schur 引理判断交织空间，区分等型分量与所选的简单直和项。
\end{itemize}
\end{learninggoals}

本章以群、商群和有限维线性代数为基础。模、张量及外代数只使用其基本构造，所需的群作用公式均在正文中给出。具体矩阵表示及其不变子空间会使可约性与完全可约性的差别直接显现。

\ProseParagraphBreak
可参阅 Etingof 等人的《Introduction to Representation Theory》\cite[Proposition 2.3.9; Corollary 2.3.10; Section 4.1, pp. 59--61]{EtingofEtAl2011RepresentationTheory}。该书从一般代数的表示讲起，再讨论有限群；本章先讨论群作用，一般底域上分解所涉及的除代数条件另在使用处说明。

% !TeX root = ../../Book5.tex
% 纲要标题：### 1.1 表示与模
% 纲要位置：工作记录/计划纲要.md，第 3723 行。
% 纲要细目（写作范围）：
% * 群的线性表示
% * 群代数模与表示的等价
% * 子表示、商表示及直和


\section{表示与模}
\label{b5:sec:0101}

一个群可以置换顶点，也可以旋转平面。表示论把后一类作用推广到任意向量空间：每个群元素给出一个可逆线性变换，群乘法对应变换复合。以下从作用的定义讲起，再把它改写成群代数模；这两种说法记录的是同一组运算。


\subsection{群的线性表示}
\label{b5:subsec:0101:01}

\begin{definition}\label{b5:def:linear-representation}
设 \(G\) 为群，\(k\) 为域，\(V\) 为有限维 \(k\) 向量空间。若映射 \(\rho:G\to\operatorname{GL}_k(V)\) 满足
\[
\rho(gh)=\rho(g)\rho(h)\quad(g,h\in G),\qquad \rho(1)=1_V,
\]
则称 \((V,\rho)\) 为 \(G\) 的一个 \(k\)-\term[xianxing-biaoshi]{线性表示}[linear representation]，\(\dim_kV\) 称为其\term[biaoshi-cishu]{次数}[degree of a representation]。通常简写 \(gv=\rho(g)v\)。若 \(\rho\) 单射，称表示\term[zhongshi-biaoshi]{忠实}[faithful representation]。
\end{definition}
\symidx{\(\rho:G\to\operatorname{GL}_k(V)\)：群的线性表示}

选择 \(V\) 的基后得到矩阵。换基矩阵为 \(P\) 时，全部矩阵同时变成 \(P^{-1}\rho(g)P\)；不能对不同群元素各自选一个无关的换基。零维空间也给出零表示，但后面“不可约”总排除零表示。

\begin{definition}\label{b5:def:intertwiner}
设 \(k\) 为域，\(G\) 为群，\((V,\rho)\)、\((W,\sigma)\) 为 \(G\) 的有限维 \(k\) 表示。若 \(k\) 线性映射 \(T:V\to W\) 满足
\[
T\rho(g)=\sigma(g)T\qquad(g\in G),
\]
则称 \(T\) 为表示同态，也称\term[jiaozhi-suanzi]{交织算子}[intertwining operator]。这类映射组成的 \(k\) 向量空间记为 \(\operatorname{Hom}_G(V,W)\)。若其中有可逆映射，则称两表示同构。
\end{definition}
\symidx{\(\operatorname{Hom}_G(V,W)\)：表示之间的交织算子空间}

若 \(T\) 可逆，上述等式乘以 \(T^{-1}\) 就表明逆也是交织算子。表示同构因而恰好是不依赖所选基的“同时相似”。例如循环群 \(C_m=\langle r\mid r^m=1\rangle\) 的表示由一个矩阵 \(A\) 给定，唯一要求为 \(A^m=I\)。给一组群生成元指定矩阵时，同样必须核验全部生成关系。

\subsection{群代数模与表示的等价}
\label{b5:subsec:0101:02}

\begin{definition}\label{b5:def:group-algebra}
设 \(G\) 为有限群，\(k\) 为域。以形式符号 \(g\in G\) 为基的向量空间
\[
k[G]=\left\{\sum_{g\in G}a_gg\mid a_g\in k\right\}
\]
上，若将群乘法按双线性延伸，即令
\[
\left(\sum_g a_gg\right)\left(\sum_h b_hh\right)
=\sum_{g,h}a_gb_h(gh),
\]
则所得含幺结合 \(k\) 代数称为 \(G\) 的\term[qun-daishu]{群代数}[group algebra]。其单位是群单位元对应的基向量。
\end{definition}
\symidx{\(k[G],kG\)：群代数}

结合律在基向量上就是群的结合律，双线性将它推广到全部元素。群代数 \(k[G]\) 交换当且仅当 \(G\) 交换：一个方向由乘法的双线性延伸立即得到；反过来，若群代数交换，则其中的基向量满足 \(gh=hg\)，基向量的线性无关性说明这也是 \(G\) 中的等式。不要把基向量 \(g\) 与一个任意线性组合混同；虽然每个基向量都可逆，群代数却一般不是除环。

\begin{proposition}\label{b5:prop:representations-group-modules}
设 \(G\) 为有限群，\(k\) 为域。在一个有限维 \(k\) 空间 \(V\) 上，\(G\) 的线性表示与幺性左 \(k[G]\) 模结构一一对应。表示同态恰为相应的 \(k[G]\) 模同态。
\end{proposition}
\begin{proof}
由表示定义
\[
\left(\sum_ga_gg\right)v=\sum_ga_g\rho(g)v.
\]
分配律直接成立，表示的乘法等式给出 \((ab)v=a(bv)\)，单位作用为恒等，所以得到模。反过来，从模限制到基向量 \(g\)，它与 \(g^{-1}\) 的作用互逆，故得到可逆线性变换，模结合律给出群同态。两次构造在基向量上互逆，线性延伸后在全部群代数上互逆。

最后，设 \((V,\rho_V)\)、\((W,\rho_W)\) 为两个表示，\(T:V\to W\) 为 \(k\) 线性映射。表示同态条件是
\[
T\rho_V(g)=\rho_W(g)T\qquad(g\in G).
\]
将两边对群代数的基向量线性延伸，这些等式等价于
\[
T(av)=aT(v)\qquad(a\in k[G],\ v\in V),
\]
其中两边分别使用 \(V\) 与 \(W\) 上的群代数作用。这正是 \(k[G]\) 模同态的条件。
\end{proof}

因此核、像、商和直和都可以在表示或模的语言中讨论。本卷始终使用左作用；群代数张量积的右模结构会在诱导表示的定义中另行写出。

\subsection{子表示、商表示及直和}
\label{b5:subsec:0101:03}

\begin{definition}\label{b5:def:subrepresentation-semisimple}
设 \(k\) 为域，\(G\) 为群，\((V,\rho)\) 为 \(G\) 的有限维 \(k\) 表示。若线性子空间 \(U\subseteq V\) 满足 \(\rho(g)U\subseteq U\) 对所有 \(g\in G\) 成立，称 \(U\) 为\term[zi-biaoshi]{子表示}[subrepresentation]。若 \(V\ne0\) 且子表示只有零与自身，称 \(V\) \term[bukeyue-biaoshi]{不可约}[irreducible representation]。若 \(V\ne0\) 且不能写成两个非零子表示的直和，称 \(V\) \term[bukefenjie-biaoshi]{不可分解}[indecomposable representation]。若 \(V\) 可写成不可约子表示的有限直和，称其\term[wanquan-keyue-biaoshi]{完全可约}[completely reducible representation]，也称半单表示。\termidx[biaoshi-bandan]{半单表示}空直和对应零表示。
\end{definition}

由于也对 \(g^{-1}\) 成立，不变条件实际上给出 \(\rho(g)U=U\)。商作用为 \(g(v+U)=gv+U\)，改变代表所产生的差仍在 \(U\)，故良定义。直和 \(V\oplus W\) 逐分量作用；交织算子的核与像同样不变。对有限群 \(G\)，子表示恰为群代数子模，因此不可约表示对应于简单 \(k[G]\) 模，也称简单表示。不可约表示一定不可分解，反过来一般不成立。零表示完全可约，但不是不可约或不可分解表示。

\begin{proposition}\label{b5:prop:invariant-complement-projection}
设 \(k\) 为域，\(G\) 为群，\(V\) 为 \(G\) 的有限维 \(k\) 表示，\(U\subseteq V\) 为子表示。\(U\) 有不变补空间，当且仅当存在 \(G\) 线性映射 \(p:V\to U\) 满足 \(p|_U=1_U\)。此外，\(V\) 完全可约当且仅当它的每个子表示都有不变补空间。
\end{proposition}
\begin{proof}
以下通过包含映射 \(U\hookrightarrow V\)，也把 \(p:V\to U\) 视为 \(V\) 的自同态。若 \(V=U\oplus W\) 且两项不变，沿 \(W\) 的投影与 \(V\) 上的群作用交换，并满足 \(p|_U=1_U\)。反过来，\(p|_U=1_U\) 给出 \(p^2=p\)。每个 \(v\in V\) 都写成 \(v=p(v)+(v-p(v))\)，两项分别属于 \(U\) 和 \(\ker p\)，且 \(U\cap\ker p=0\)。因此 \(V=U\oplus\ker p\)，而交织算子的核不变。

若 \(V\) 的每个子表示都有不变补空间，对 \(\dim_kV\) 归纳证明完全可约性。零表示对应空直和。若 \(V\ne0\)，在非零子表示中选维数最小者 \(S\)，则 \(S\) 不可约。取不变补空间 \(W\)，使 \(V=S\oplus W\)。还须验证 \(W\) 继承归纳所需的条件：给定子表示 \(N\subseteq W\)，在 \(V\) 中取不变补 \(D\)，使 \(V=N\oplus D\)。由于 \(N\subseteq W\)，分解任意 \(w\in W\) 时，其 \(D\) 分量仍在 \(W\) 中，所以
\[
W=N\oplus(W\cap D).
\]
两项都不变。因 \(\dim_kW<\dim_kV\)，归纳假设说明 \(W\) 是不可约表示的直和，从而 \(V=S\oplus W\) 完全可约。

反过来，设 \(V=S_1\oplus\cdots\oplus S_r\) 为不可约子表示的直和，对 \(r\) 归纳证明每个子表示 \(U\) 都可补。\(r=0\) 时结论显然；任意 \(r\) 下，\(U=0\) 都以 \(V\) 为补。若 \(U\ne0\)，在 \(U\) 中选一个维数最小的非零子表示 \(T\)，它不可约。记坐标投影为 \(q_i:V\to S_i\)。至少一个 \(q_i|_T\) 非零，否则 \(T=0\)。该映射的核是 \(T\) 的子表示，像是 \(S_i\) 的子表示，故不可约性说明
\[
q_i|_T:T\xrightarrow{\sim}S_i.
\]
令 \(R=\bigoplus_{j\ne i}S_j\)。对任意 \(v\in V\)，取
\[
t=(q_i|_T)^{-1}(q_i(v)).
\]
则 \(v-t\in\ker q_i=R\)，且 \(T\cap R=\ker(q_i|_T)=0\)，所以 \(V=T\oplus R\)。若 \(v\in U\)，因 \(T\subseteq U\)，还得到 \(v-t\in U\cap R\)，故
\[
U=T\oplus(U\cap R).
\]
\(R\) 是 \(r-1\) 个不可约子表示的直和，归纳假设给出不变子空间 \(C\subseteq R\)，使 \(R=(U\cap R)\oplus C\)。于是
\[
V=T\oplus(U\cap R)\oplus C=U\oplus C,
\]
这就是 \(U\) 的不变补空间。
\end{proof}

% !TeX root = ../../Book5.tex
% 纲要标题：### 1.2 典型表示
% 纲要位置：工作记录/计划纲要.md，第 3729 行。
% 纲要细目（写作范围）：
% * 正则表示和置换表示
% * 一维表示及交换化
% * 对偶、张量积和外幂表示


\section{典型表示}
\label{b5:sec:0102}

抽象定义允许许多不同作用。最常用的表示往往直接来自一个有限集合：把集合元素改成向量空间的基，置换便成为线性变换。另一些表示则从已有表示通过对偶和张量构造得到。以下各项都写明作用公式，以便后面计算迹和不变子空间。


\subsection{正则表示和置换表示}
\label{b5:subsec:0102:01}

设有限群 \(G\) 作用于有限集合 \(X\)。在以 \(e_x\)、\(x\in X\)，为基的 \(k\) 空间 \(k[X]\) 上定义 \(ge_x=e_{gx}\)，得到\term[zhihuan-biaoshi]{置换表示}[permutation representation]。矩阵每行每列各有一个一，其余位置为零。\(X=G\) 且采用左平移时，得到\term[zhengze-biaoshi]{正则表示}[regular representation]，也就是 \(k[G]\) 在自身上的左乘作用。

\ProseParagraphBreak
正则表示忠实，因为 \(g\cdot1=g\)，不同基向量不相等。对一般 \(X\)，不动向量
\[
V^G=\bigl\{v\in k[X]\mid gv=v\text{ 对所有 }g\in G\bigr\}
\]
在同一个轨道上的系数必须相同。因此每个轨道的基向量之和给出一组不动空间的基，\(\dim_k k[X]^G\) 等于轨道数，任意特征均成立。

\begin{example}\label{b5:exm:permutation-augmentation}
设 \(S_n\) 置换 \(k^n\) 的标准基，\(n\ge2\)。令
\[
L=k(e_1+\cdots+e_n),\qquad
W=\left\{\sum_i a_ie_i\mid\sum_i a_i=0\right\}.
\]
两者都是子表示。当 \(n\ne0\) 于 \(k\) 时，\(k^n=L\oplus W\)；当 \(\operatorname{char}k\mid n\) 时，却有 \(L\subseteq W\)。
\end{example}
\begin{proof}
置换不改变系数和，故两者不变。常数向量 \(a\sum_i e_i\) 的系数和为 \(na\)。若 \(n\) 可逆，交为零且维数相加为 \(n\)，得到直和；否则整个 \(L\) 都在 \(W\) 中。
\end{proof}

\subsection{一维表示及交换化}
\label{b5:subsec:0102:02}

一维 \(k\) 表示就是群同态 \(\lambda:G\to k^\times\)，称为\term[yiwei-biaoshi]{一维表示}[one-dimensional representation]。像群交换，所以 \(\lambda(ghg^{-1}h^{-1})=1\)。记 \(G'=[G,G]\) 为所有交换子生成的子群，则 \(\lambda\) 唯一经由 \(G/G'\) 分解。反过来，\(G/G'\to k^\times\) 的同态与商映射复合就给出一维表示。这直接说明一维表示只能观察群的交换化。

\ProseParagraphBreak
在复数域上，\(C_m=\langle r\rangle\) 的一维表示为
\[
\lambda_j(r)=\exp(2\pi\mathrm i j/m),\qquad 0\le j<m.
\]
在实数域上，有限群的一维表示的值只能为 \(\pm1\)；阶为奇数的群因而只有平凡实一维表示。同一个群换底域以后，可用的一维表示会改变。

\ProseParagraphBreak
若 \(\lambda\) 为一维表示，\(\rho^\lambda(g)=\lambda(g)\rho(g)\) 仍满足表示乘法法则，称为用 \(\lambda\) 扭曲。子空间在扭曲前后不变的条件相同，因此不可约性不变。例如对称群的符号表示可将一个不可约表示变成另一个，而不改变次数。

\subsection{对偶、张量积和外幂表示}
\label{b5:subsec:0102:03}

设 \(V,W\) 是群 \(G\) 在域 \(k\) 上的有限维表示。对偶空间 \(V^*=\operatorname{Hom}_k(V,k)\) 上定义
\begin{equation}\label{b5:eq:dual-tensor-action}
(g\varphi)(v)=\varphi(g^{-1}v),\qquad
g(v\otimes w)=gv\otimes gw.
\end{equation}
前式给出\term[duiou-biaoshi]{对偶表示}[dual representation]，后式给出\term[zhangliang-biaoshi]{张量积表示}[tensor product representation]。逆元不可省略：计算
\[
g(h\varphi)(v)=\varphi(h^{-1}g^{-1}v)=\bigl((gh)\varphi\bigr)(v)
\]
才得到左表示的次序。对偶矩阵是 \(\rho(g^{-1})^{\mathsf T}\)，没有共轭；复共轭只会在另行选择 Hermitian 内积后出现。

\ProseParagraphBreak
外幂上的作用由
\[
g(v_1\wedge\cdots\wedge v_r)=gv_1\wedge\cdots\wedge gv_r
\]
定义，因为线性变换保持交替关系，故能从张量幂下降到商空间。特别地，最高外幂是一维表示 \(g\mapsto\det\rho(g)\)。在特征二，外代数仍按 \(v\wedge v=0\) 定义，不能仅用反交换公式替代交替条件。

\begin{proposition}\label{b5:prop:hom-tensor-invariants}
设 \(k\) 为域，\(G\) 为群，\(V,W\) 为 \(G\) 的有限维 \(k\) 表示。在 \(\operatorname{Hom}_k(V,W)\) 上置 \((gT)(v)=gT(g^{-1}v)\)，其中 \(g\in G\)、\(T\in\operatorname{Hom}_k(V,W)\)、\(v\in V\)，则自然映射
\[
W\otimes_kV^*\longrightarrow\operatorname{Hom}_k(V,W),
\qquad w\otimes\varphi\longmapsto\bigl(v\mapsto\varphi(v)w\bigr)
\]
是表示同构，且其不动空间恰为 \(\operatorname{Hom}_G(V,W)\)。
\end{proposition}
\begin{proof}
选 \(V\) 的基 \(v_i\) 及对偶基 \(\varphi_i\)，逆映射为 \(T\mapsto\sum_iT(v_i)\otimes\varphi_i\)，与所选基无关，因为它是给定自然映射的唯一逆。两侧作用于纯张量后均将 \(v\) 送到 \(\varphi(g^{-1}v)gw\)，故同构与作用相容。最后 \(gT=T\) 等价于 \(T(gv)=gT(v)\)，即交织条件。
\end{proof}

% !TeX root = ../../Book5.tex
% 纲要标题：### 1.3 Maschke 定理
% 纲要位置：工作记录/计划纲要.md，第 3735 行。
% 本轮补充的纲要细目：
% * 在对偶空间的对称代数的各齐次分量上应用平均法，定义不变多项式、不变环和 Reynolds 算子，并计算反射例子
% * 用低次循环群算例核验平均投影与对偶作用的方向；生成元、关系及商的系统理论留第22章
% #### 1.3.1 群阶的可逆性与平均算子
% #### 1.3.2 平均投影与 Maschke 定理
% #### 1.3.3 模特征情形的反例
% #### 1.3.4 多项式上的平均与 Reynolds 算子
% 纲要细目（写作范围）：
% * 特征不整除群阶的条件
% * 平均投影的证明
% * 模特征情形的反例


\section{Maschke 定理}
\label{b5:sec:0103}

线性空间中的子空间总有补空间，但任意选择的补未必被群保持。有限群允许对所有群元素取平均；只要群阶在底域中可逆，就能把一个普通投影改成与群作用相容的投影。这个除法条件决定了所有有限维表示是否都完全可约，而不决定某个指定表示是否完全可约。


\subsection{群阶的可逆性与平均算子}
\label{b5:subsec:0103:01}

对有限群 \(G\)，若 \(\operatorname{char}k=0\)，或 \(\operatorname{char}k=p>0\) 且 \(p\mathrel{\vcenter{\hbox{\scalebox{1.25}{\(\nmid\)}}}}|G|\)，则 \(|G|\cdot1_k\ne0\)，可以在 \(k\) 中除以 \(|G|\)。此时，在任意有限维表示 \(V\) 上定义
\[
E_V=\frac1{|G|}\sum_{g\in G}\rho(g).
\]
对 \(h\in G\)，左乘或右乘 \(h\) 只是重排求和项，故 \(\rho(h)E_V=E_V=E_V\rho(h)\)。因此像包含于 \(V^G\)，而不动向量经此映射保持不变，于是 \(E_V^2=E_V\)，像恰为 \(V^G\)。

\ProseParagraphBreak
这里的分母是真正的标量逆，不是一个形式记号。若特征整除群阶，同一分子 \(\sum_g\rho(g)\) 仍有意义，但对每个 \(v\in V^G\)，都有 \(\sum_g\rho(g)v=|G|v=0\)。求和算子因而在不动空间上恒为零，一般不能用它构造到不动空间的投影。

\subsection{平均投影与 Maschke 定理}
\label{b5:subsec:0103:02}

\begin{theorem}[Maschke]\label{b5:thm:maschke}
设 \(G\) 为有限群，\(k\) 为域，且 \(|G|\ne0\) 于 \(k\)。\(G\) 的每个有限维 \(k\) 表示 \(V\) 的任意子表示 \(U\) 都有不变补空间，因而每个这样的表示完全可约。
\end{theorem}
\begin{proof}
取一个普通线性投影 \(p:V\to U\)，使 \(p|_U=1_U\)，并通过包含 \(U\hookrightarrow V\) 把它看作 \(V\) 的自同态。定义
\[
\overline p=\frac1{|G|}\sum_{g\in G}\rho(g)p\rho(g^{-1}).
\]
每一项的像都在 \(U\)，因为 \(U\) 不变。对 \(u\in U\)，先有 \(\rho(g^{-1})u\in U\)，再经 \(p\) 不变，所以每一项将 \(u\) 送回 \(u\)，得到 \(\overline p|_U=1_U\)。对任意 \(h\)，将求和指标 \(g\) 改成 \(hg\)，得
\(\rho(h)\overline p\rho(h^{-1})=\overline p\)。
于是 \(\overline p\) 是 \(G\) 线性投影，\cref{b5:prop:invariant-complement-projection}给出 \(V=U\oplus\ker\overline p\)，并推出完全可约性。
\end{proof}

这就是\term[Maschke-dingli]{Maschke 定理}[Maschke's theorem]。平均只修正投影，没有假定已经知道不可约表示有哪些。因此它先提供分解的存在性，分类和重数仍需其他工具。

\ProseParagraphBreak
例如设 \(\operatorname{char}k\ne2\)，令 \(C_2=\langle s\rangle\) 在 \(k^2\) 上交换两个坐标。取子表示 \(U=k(1,1)\)，以及到 \(U\) 的普通投影
\[
\rho(s)=\begin{pmatrix}0&1\\1&0\end{pmatrix},\qquad
p=\begin{pmatrix}1&0\\1&0\end{pmatrix}.
\]
\(p\) 的像是 \(U\)，且在 \(U\) 上为恒等，但 \(p\rho(s)\ne\rho(s)p\)。共轭平均把它改成
\[
\overline p=\dfrac12\bigl(p+\rho(s)p\rho(s)\bigr)
=\dfrac12\begin{pmatrix}1&1\\1&1\end{pmatrix}.
\]
于是 \(\ker\overline p=k(1,-1)\)，得到了不变直和分解
\(k^2=k(1,1)\oplus k(1,-1)\)。原投影的核 \(k(0,1)\) 不被 \(s\) 保持；平均后的核则被保持。

\begin{proposition}\label{b5:prop:unitary-average}
设 \(G\) 为有限群，\(\rho:G\to\operatorname{GL}_{\mathbb C}(V)\) 为有限维复表示。\(V\) 上存在被 \(G\) 保持的正定 Hermitian 内积；相对于其标准正交基，每个 \(\rho(g)\) 都是酉矩阵。
\end{proposition}
\begin{proof}
任选正定 Hermitian 内积 \((\, ,\,)_0\)，约定第一变量共轭线性、第二变量线性。置
\[
(v,w)=\frac1{|G|}\sum_g\bigl(\rho(g)v,\rho(g)w\bigr)_0.
\]
令 \(w=v\ne0\)，则求和中的每一项均为正实数，所以正定；两变量的线性规则不变，重排求和给出不变性。取标准正交基即得酉矩阵。若 \(U\) 不变，其正交补也不变，因为 \((u,gw)=(g^{-1}u,w)=0\)。这另给出复数情形的不变补空间。
\end{proof}

\subsection{模特征情形的反例}
\label{b5:subsec:0103:03}

若 \(k\) 的特征为 \(p\)，在 \(C_p=\langle r\rangle\) 上取
\[
\rho(r)=\begin{pmatrix}1&1\\0&1\end{pmatrix}.
\]
写为 \(I+N\)，其中 \(N^2=0\)。二项式展开给出 \((I+N)^p=I\)，所以这是表示。直线 \(ke_1\) 不变。任何补直线可由 \(ae_1+e_2\) 生成；若它不变，\(r\) 在这条一维空间上的标量的 \(p\) 次方为一，在特征 \(p\) 下只能是一。但是
\[
\rho(r)(ae_1+e_2)=(a+1)e_1+e_2\ne ae_1+e_2,
\]
故不存在不变补。这个表示可约但不可分解。

\begin{proposition}\label{b5:prop:maschke-converse}
设 \(G\) 为有限群，\(k\) 为域。若 \(\operatorname{char}k\,\bigm\vert\,|G|\)，则正则表示 \(k[G]\) 不完全可约。因此 \(G\) 的全部有限维 \(k\) 表示完全可约，当且仅当群阶在 \(k\) 中可逆。
\end{proposition}
\begin{proof}
增广映射 \(\varepsilon:k[G]\to k\)、\(\sum a_gg\mapsto\sum a_g\)，在目标采用平凡作用时为满的表示同态。若它有 \(G\) 线性截面 \(s:k\to k[G]\)，则 \(1\in k\) 在截面下的像 \(v=s(1)\) 必为不动向量。正则作用迫使 \(v=c\sum_g g\)，所以 \(\varepsilon(v)=c|G|=0\)，不能等于一。于是其核没有不变补，正则表示不完全可约。反方向正是\cref{b5:thm:maschke}。
\end{proof}

\subsection{多项式上的平均与 Reynolds 算子}
\label{b5:subsec:0103:reynolds}

线性作用还会作用于多项式。设 \(V\) 是有限维 \(k\) 表示，取对偶空间的一组基 \(x_1,\ldots,x_n\)，写
\[
k[V]\coloneqq\operatorname{Sym}(V^*)=k[x_1,\ldots,x_n].
\]
对偶作用 \(g\ell=\ell\circ\rho(g^{-1})\) 唯一延伸为这个多项式环的代数自同构。在向量处求值时，公式为
\[
(g f)(v)=f\bigl(\rho(g^{-1})v\bigr).
\]
这里先按线性替换定义形式多项式上的作用；尤其当 \(k\) 有限时，不能把两个多项式在全部 \(k\) 点取值相同当成多项式相等。逆元则保证 \(g(hf)=(gh)f\)。

\begin{definition}\label{b5:def:polynomial-invariant-ring}
设 \(k\) 为域，\(G\) 为群，\(\rho:G\to\operatorname{GL}_k(V)\) 为有限维表示。在 \(k[V]=\operatorname{Sym}(V^*)\) 上，将对偶作用 \(g\ell=\ell\circ\rho(g^{-1})\) 延伸为代数自同构，其中 \(\ell\in V^*\)、\(g\in G\)。满足 \(gf=f\) 对所有 \(g\in G\) 成立的多项式 \(f\in k[V]\)，称为\term[bubian-duoxiangshi]{不变多项式}[invariant polynomial]。它们组成的 \(k\) 子代数
\[
k[V]^G=\bigl\{\,f\in k[V]\mid gf=f\text{ 对所有 }g\in G\,\bigr\}
\]
称为该作用的\term[bubian-huan]{不变环}[invariant ring]。若 \(G\) 有限且 \(|G|\ne0\) 于 \(k\)，则映射
\[
\mathcal R(f)=\frac1{|G|}\sum_{g\in G}gf
\]
称为\term[Reynolds-suanzi]{Reynolds 算子}[Reynolds operator]。\footnote{雷诺（Osborne Reynolds，1842-1912），爱尔兰数学家、物理学家，以流体流动中的雷诺数和湍流研究著称。}
\end{definition}

线性替换保持各个齐次分量，故 \(k[V]_d\) 是有限维表示。把本节开头的平均投影分别用于每个 \(k[V]_d\)，可知 \(\mathcal R(k[V]_d)\subseteq k[V]_d\)，即 \(\mathcal R\) 保持分次；它的像恰为 \(k[V]^G\)，且 \(\mathcal R^2=\mathcal R\)。平均可能消去最高次齐次部分，所以当 \(\mathcal R(f)\ne0\) 时，只能断言 \(\deg\mathcal R(f)\le\deg f\)。对于 \(a\in k[V]^G\)，逐项使用 \(g(af)=a\,g(f)\)，还有
\[
\mathcal R(af)=a\mathcal R(f).
\]
所以平均可以同乘以不变多项式交换，但一般不保持任意乘积。

\begin{example}\label{b5:exm:reynolds-reflection}
设 \(\operatorname{char}k\ne2\)，令 \(C_2\) 在 \(k^2\) 上作用为 \((u,v)\mapsto(-u,v)\)。在对偶坐标 \(x,y\) 中，
\[
\mathcal R(f)=\frac{f(x,y)+f(-x,y)}2.
\]
它删掉所有含奇数次 \(x\) 的单项式，所以不变环为 \(k[x^2,y]\)。这个生成结论来自逐个单项式的计算，不仅来自平均公式。例如 \(\mathcal R(x)=0\)，却有 \(\mathcal R(x^2)=x^2\)，因此 \(\mathcal R\) 不是环同态。又如 \(\mathcal R(x^3+y^2)=y^2\)，说明保持分次不等于保持每个多项式的次数。
\end{example}

这里已能把不变多项式按次数作为表示来研究。至于怎样找到有限个生成元、确定它们之间的全部关系，以及能否由不变量区分两个轨道，留待\chapref{b5:ch:22}讨论。

% !TeX root = ../../Book5.tex
% 纲要标题：### 1.4 Schur 引理与分解
% 纲要位置：工作记录/计划纲要.md，第 3753 行。
% 纲要细目（写作范围）：
% * 不可约表示的自同态环
% * 代数闭域和一般域的区别
% * 等型分解、重数与正则表示；一般域上的自同态除代数与代数闭域情形的区别
% ---


\section{Schur 引理与分解}
\label{b5:sec:0104}

完全可约性让我们能够拆开表示，还须知道拆出的成分是否唯一。交织算子给出这个问题的答案：两个不可约成分之间的非零映射必可逆，而一个不可约表示的自同态环一般是除代数。只有在合适的底域上，才能再把它缩成标量。


\subsection{不可约表示的自同态环}
\label{b5:subsec:0104:01}

\begin{lemma}[Schur]\label{b5:lem:schur}
设 \(G\) 为群，\(k\) 为域，\(V,W\) 为非零不可约有限维 \(k\) 表示。每个非零 \(T\in\operatorname{Hom}_G(V,W)\) 都是同构。因此不相同构时 \(\operatorname{Hom}_G(V,W)=0\)，而 \(\operatorname{End}_G(V)\) 是一个有限维 \(k\) 除代数。
\end{lemma}
\symidx{\(\operatorname{End}_G(V)\)：表示的自同态代数}
\begin{proof}
\(\ker T\) 与 \(\operatorname{im}T\) 都是子表示。非零性及不可约性使它们分别等于零和 \(W\)，所以 \(T\) 双射；其逆满足 \(T^{-1}(gw)=gT^{-1}(w)\)，仍为表示同态。对 \(V=W\) 就得到每个非零自同态可逆；它作为 \(\operatorname{End}_k(V)\) 的线性子空间有限维。
\end{proof}

这称为\term[Schur-yinli]{Schur 引理}[Schur's lemma]。其第一部分不要求群有限，也不需要完全可约性；它只比较已经不可约的表示。

\subsection{代数闭域与一般域}
\label{b5:subsec:0104:02}

\begin{corollary}\label{b5:cor:schur-algebraically-closed}
设 \(k\) 代数闭，\(G\) 为群，\(V\) 为非零不可约有限维 \(k\) 表示。则
\[
\operatorname{End}_G(V)=k\,1_V.
\]
特别地，有限交换群在代数闭域上的不可约表示都是一维的。
\end{corollary}
\begin{proof}
对 \(T\in\operatorname{End}_G(V)\)，代数闭性保证存在特征值 \(\lambda\in k\)。\(T-\lambda I\) 不可逆，由 Schur 引理只能为零。若群交换，每个 \(\rho(g)\) 都是交织算子，因而都是标量，于是每条直线不变。不可约性迫使维数为一。
\end{proof}

底域条件不能省去。例如实空间 \(\mathbb R^2\) 上，令 \(C_3\) 的生成元作用为旋转 \(120^\circ\) 的矩阵。它没有实特征向量，所以没有不变实直线，表示不可约。与旋转交换的实矩阵却为
\[
\begin{pmatrix}a&-b\\b&a\end{pmatrix},\qquad a,b\in\mathbb R,
\]
构成与 \(\mathbb C\) 同构的实除代数，而非只有实标量。复化后旋转出现两个不同特征值，表示分裂成两个一维复表示。

\ProseParagraphBreak
即使底域代数闭，结论“自同态只有标量”也不能无条件反推不可约。取 \(k=\overline{\mathbb F}_3\)，令 \(r=(123)\)、\(s=(12)\) 在 \(k^2\) 上分别作用为
\[
R=\begin{pmatrix}1&1\\0&1\end{pmatrix},\qquad
S=\begin{pmatrix}-1&0\\0&1\end{pmatrix}.
\]
写 \(R=I+N\)，则 \(N^2=0\)，所以在特征 \(3\) 下
\[
R^3=I,\qquad S^2=I,\qquad
SRS=I-N=R^{-1}.
\]
这些正是 \(S_3\) 的生成关系，因而给出一个表示。直线 \(ke_1\) 是非零真子表示，故该表示可约。另一方面，与 \(S\) 交换的矩阵必须为对角矩阵，因为 \(-1\ne1\)；再与 \(R\) 交换，就迫使两个对角元相等。因此
\[
\operatorname{End}_{S_3}(k^2)=kI.
\]
若一个非零的完全可约表示能分成两个非零子表示，相应投影就是一个非标量自同态。因此，非零的完全可约表示若只有标量自同态，就必不可约。

\subsection{等型分解、重数与正则表示}
\label{b5:subsec:0104:03}

\begin{proposition}\label{b5:prop:isotypic-decomposition}
设 \(G\) 为有限群，\(k\) 为域，\(V\) 为完全可约的有限维左 \(k[G]\) 模。设
\[
V\simeq\bigoplus_{i=1}^r S_i^{\oplus m_i}
\]
为其分解，其中 \(S_i\) 为两两不同构的不可约有限维左 \(k[G]\) 模，\(m_i\) 为非负整数。记
\[
D_i=\operatorname{End}_G(S_i),\qquad
M_i=\operatorname{Hom}_G(S_i,V).
\]
称 \(M_i\) 为类型 \(S_i\) 的\term[chongshu-kongjian]{重数空间}[multiplicity space]，并给它平凡 \(G\) 作用。则
\[
\dim_k M_i=m_i\dim_k D_i.
\]
若 \(V\ne0\)，则
\[
\operatorname{End}_G(V)\simeq\prod_{\substack{1\le i\le r\\m_i>0}}M_{m_i}(D_i),
\]
即乘积中略去零重数项。若 \(V=0\)，则各 \(m_i=0\)，且 \(\operatorname{End}_G(V)\) 只有零映射。
相应的\term[dengxing-fenliang]{等型分量}[isotypic component] \(V_i\)，即 \(V\) 中所有同构于 \(S_i\) 的简单子表示之和，是内在确定的；其内部拆成 \(m_i\) 个简单直和项通常不唯一。
若 \(k\) 代数闭，则 \(D_i=k\,1_{S_i}\)，所以 \(m_i=\dim_kM_i\)。此时求值映射给出等变同构
\[
\varepsilon_i:S_i\otimes_kM_i\longrightarrow V_i,
\qquad v\otimes f\longmapsto f(v),
\]
其中 \(g(v\otimes f)=gv\otimes f\)。
\end{proposition}
\symidx{\(V_i,M_i\)：完全可约表示的等型分量与重数空间}
\begin{proof}
把同态分别投到各个直和坐标，Schur 引理使不同类型间的块为零，相同类型间的块为 \(D_i\) 中的元素。因此 \(M_i\) 与 \(D_i^{\oplus m_i}\) 作为 \(k\) 向量空间同构，得到维数公式。自同态在所选分解下由同型块内的矩阵给出；两个自同态复合时，其第 \((a,b)\) 项是 \(\sum_c d_{ac}\circ d'_{cb}\)。按自同态复合定义 \(D_i\) 的乘法，这就是 \(M_{m_i}(D_i)\) 的通常矩阵乘法，故得到自同态代数的乘积公式。若 \(V=0\)，所有结论直接成立。

任意同构于 \(S_i\) 的子表示向其他类型的坐标投影都为零，所以它全在所选分解的 \(S_i^{\oplus m_i}\) 中；反过来，这些坐标子表示已生成整个块。故所有同型简单子表示之和恰为该块，不依赖所选分解；当 \(m_i=0\) 时这个和为零。

若 \(k\) 代数闭，Schur 引理给出 \(D_i=k\,1_{S_i}\)。求值公式对 \(v,f\) 双线性，且像包含于 \(V_i\)，所以诱导所述张量映射。由 \(f(gv)=gf(v)\)，它是等变的。在所选分解下，\(M_i\) 的一组基由到各个 \(S_i\) 坐标的嵌入组成，求值映射就是 \(m_i\) 份 \(S_i\) 的恒等识别，故为同构。
\end{proof}

在代数闭域上，选择重数空间 \(M_i\) 的一组基 \(f_1,\ldots,f_{m_i}\)，就由求值同构得到
\[
V_i=\bigoplus_{j=1}^{m_i}\varepsilon_i\bigl(S_i\otimes_k kf_j\bigr).
\]
不同的基可能给出不同的简单直和项；但把基向量分别乘以非零标量，或只改变它们的排列，并不改变这些子空间组成的无序集合。当 \(m_i=1\) 时，唯一的简单直和项就是 \(V_i\) 本身。

\ProseParagraphBreak
例如，设 \(k\) 为代数闭域，\(S\) 为群 \(G\) 的一个非零不可约 \(k\) 表示。在 \(S\oplus S\) 中，
\[
U_1=S\oplus0,\qquad U_2=0\oplus S,\qquad
\Delta=\{\,(v,v):v\in S\,\}
\]
都是同构于 \(S\) 的简单子表示，而
\[
S\oplus S=U_1\oplus U_2=U_1\oplus\Delta.
\]
后一分解可由 \((v,w)=(v-w,0)+(w,w)\) 直接验证，且 \(U_1\cap\Delta=0\)。两种分解中的具体子空间不同，但等型分量都是整个 \(S\oplus S\)，重数都是 \(2\)。在重数空间中，这对应于将基 \(f_1(v)=(v,0)\)、\(f_2(v)=(0,v)\) 改为 \(f_1,f_1+f_2\)。

\ProseParagraphBreak
一般域上的重数空间则还有右 \(D_i\) 作用 \(f\cdot d=f\circ d\)，其右 \(D_i\) 维数为 \(m_i\)；它的 \(k\) 维数还要乘以 \(\dim_kD_i\)，不能直接当作重数。这里出现的矩阵块和除代数，也正是半单代数理论中的基本结构。

\begin{lemma}\label{b5:lem:irreducibles-in-regular}
设 \(G\) 为有限群，\(k\) 为域，\(|G|\ne0\) 于 \(k\)。每个不可约有限维 \(k\) 表示都同构于正则表示的某个直和项，因此不可约表示只有有限多个同构类型。
\end{lemma}
\begin{proof}
对不可约 \(S\) 选 \(0\ne v\in S\)。映射 \(k[G]\to S\)、\(a\mapsto av\)，像是非零子表示，故满射。由 Maschke 定理，其核有不变补，所以 \(S\) 是正则表示的直和项。把有限维正则表示分成有限个不可约项；任意不可约 \(S\) 向这些项的至少一个投影非零，Schur 引理使它同构于该项。因而该有限列表已经穷尽所有类型。
\end{proof}


\begin{chaptersummary}
群代数把有限个群作用线性组合起来，表示与左模因而给出等价描述。置换和正则表示从群作用直接构造，对偶和张量则从已有表示产生新表示；每种构造都须保留正确的作用次序。

\ProseParagraphBreak
群阶可逆时，平均将普通投影改成不变投影，所有有限维表示完全可约。模特征中的增广映射不分裂，说明分解结论不能脱离特征条件。不可约与不可分解在这里也必须区分。

\ProseParagraphBreak
把平均逐次用于对称代数，就得到保持次数的 Reynolds 算子。它是到不变多项式的投影，并对不变子环线性；反射的例子说明这种投影一般不是环同态。生成元与关系的问题因此需要在平均之外继续研究。

\ProseParagraphBreak
Schur 引理控制不可约之间的同态，代数闭性进一步把自同态除代数化成标量。等型分量及其重数由表示本身确定，而同型成分内部的直和分解依赖选择。下一章将用迹与正交关系把这些重数变成可计算的数值。
\end{chaptersummary}

\BookExercises

\begin{exercise}\label{b5:ex:01-cyclic-diagonalization}
设 \(k\) 中 \(m\) 可逆，且 \(T^m-1\) 在 \(k\) 上分裂。证明 \(C_m\) 的每个有限维 \(k\) 表示都能选基使全部作用矩阵同时对角，并用生成元的特征空间写出分解。
\end{exercise}
\begin{solution}
生成元矩阵 \(A\) 满足 \(A^m=I\)。多项式 \(T^m-1\) 与导数 \(mT^{m-1}\) 互素，故分裂为互异一次因子。最小多项式也无重根，所以 \(V=\bigoplus_{\lambda^m=1}\ker(A-\lambda I)\)。群元素均为生成元的幂，在每个该特征空间上以 \(\lambda^j\) 作用；对各空间分别选基便得同时对角化。
\end{solution}

\begin{exercise}\label{b5:ex:01-change-basis}
在 \(\mathbb C^2\) 上令 \(C_2\) 生成元分别作用为 \(\operatorname{diag}(1,-1)\) 和 \(\begin{pmatrix}0&1\\1&0\end{pmatrix}\)。写出两表示之间的同构及全部交织算子。
\end{exercise}
\begin{solution}
取 \(P=\begin{pmatrix}1&1\\1&-1\end{pmatrix}\)，有 \(BP=PA\)，且 \(\det P=-2\ne0\)，所以 \(P\) 给同构。任意交织 \(T\) 的第一列须属于 \(B\) 的一特征空间，第二列属于负一特征空间，故
\[
T=\begin{pmatrix}a&b\\a&-b\end{pmatrix}.
\]
它可逆恰好要求 \(ab\ne0\)。
\end{solution}

\begin{exercise}\label{b5:ex:01-regular-hom}
设 \(G\) 为有限群，\(k\) 为域，\(V\) 为有限维表示。证明评价 \(T\mapsto T(1)\) 给出 \(\operatorname{Hom}_G\bigl(k[G],V\bigr)\simeq V\)。不假设群阶可逆。
\end{exercise}
\begin{solution}
交织条件给出 \(T(g)=gT(1)\)，再由线性性得到 \(T(a)=a\cdot T(1)\)，所以该评价单射。给定 \(v\in V\)，公式 \(a\mapsto av\) 定义左模同态，评价为 \(v\)，故满射。两构造都是线性的，且对 \(V\) 的表示同态相容。
\end{solution}

\begin{exercise}\label{b5:ex:01-faithful-linear}
证明有限群有忠实一维复表示，当且仅当它是循环群。
\end{exercise}
\begin{solution}
忠实一维表示将群嵌入 \(\mathbb C^\times\)。设群阶为 \(m\)，每个像的 \(m\) 次方为一，故像是由 \(\exp(2\pi\mathrm i/m)\) 生成的循环群的子群，仍循环。反向将 \(C_m\) 的生成元送到这个本原单位根，得到忠实表示。
\end{solution}

\begin{exercise}\label{b5:ex:01-permutation-hom-orbits}
设有限群 \(G\) 作用于有限集合 \(X,Y\)。证明 \(\dim_k\operatorname{Hom}_G\bigl(k[X],k[Y]\bigr)\) 等于对角作用下 \(Y\times X\) 的轨道数，任意域均成立。计算 \(S_n\) 在自然 \(n\) 点置换表示上的自同态维数，其中 \(n\ge2\)。
\end{exercise}
\begin{solution}
矩阵条目记为 \(a_{y,x}\)，交织条件恰为 \(a_{gy,gx}=a_{y,x}\)。每个轨道上可自由选一个常数，各轨道指示矩阵线性无关且生成全部交织算子，所以维数为轨道数。\(S_n\) 在有序点对上有对角与非对角两个轨道，故维数二；矩阵可写为 \(aI+bJ\)，其中 \(J\) 的全部条目为一。
\end{solution}

\begin{exercise}\label{b5:ex:01-modular-simple}
在 \(\mathbb F_2^3\) 中，设 \(S_3\) 置换三个坐标。证明和为零的二维子空间不可约，虽然底域特征整除 \(|S_3|\)。
\end{exercise}
\begin{solution}
该子空间的非零向量为 \((1,1,0),(1,0,1),(0,1,1)\)，群在三者上传递。一个不变的一维子空间只有一个非零向量，若存在就要求该向量被全群固定，矛盾。因此此二维表示不可约。Maschke 条件失效只是说并非全部表示完全可约，不是说不可约表示不再存在。
\end{solution}

\begin{exercise}\label{b5:ex:01-invariants-exact}
设有限群 \(G\) 的阶在 \(k\) 中可逆。证明对短正合表示列 \(0\to U\to V\to W\to0\)，取不动空间后仍短正合。
\end{exercise}
\begin{solution}
核及单射可以直接在原列中检验。对 \(w\in W^G\) 选任一提升 \(v\in V\)，平均 \(v_0=|G|^{-1}\sum_g gv\) 不动，且映到 \(|G|^{-1}\sum_g gw=w\)。因此右端满射。这个提升证明明确使用了群阶的逆元。
\end{solution}

\begin{exercise}\label{b5:ex:01-double-dual}
证明有限维表示的自然映射 \(V\to V^{**}\)、\(v\mapsto\bigl(\varphi\mapsto\varphi(v)\bigr)\)，是表示同构。说明对偶作用中的逆元如何保证等变性。
\end{exercise}
\begin{solution}
有限维线性代数使评价映射双射。对 \(g\in G\)，双对偶中 \(g\) 作用于评价泛函的值为 \((g^{-1}\varphi)(v)=\varphi(gv)\)，正好是 \(gv\) 的评价。若在第一次对偶中漏掉逆元，复合的作用次序就不再是左群作用，此计算也无法保持给定的等变关系。
\end{solution}

\begin{exercise}\label{b5:ex:01-coevaluation}
设 \(V\) 为 \(d\) 维表示。证明 \(\sum_i v_i\otimes\varphi_i\in V\otimes V^*\) 与所选对偶基无关且为不动向量；评价 \(v\otimes\varphi\mapsto\varphi(v)\) 将它送到什么？当 \(\operatorname{char}k\mid d\) 时，这一复合为什么不能给出平凡表示的分裂？
\end{exercise}
\begin{solution}
在\cref{b5:prop:hom-tensor-invariants}的自然识别下，这个张量对应恒等算子，因此与基无关且不动。评价等于算子的迹，所以值为 \(d\cdot1_k\)。若该值为零，则从 \(k\) 送入这个不动向量、再评价回 \(k\) 的复合为零，不能是恒等映射。这里仅否定这对具体映射构成分裂，不排除使用其他映射得到补空间。
\end{solution}

\begin{exercise}\label{b5:ex:01-exterior-dual}
对 \(d\ge1\) 的 \(d\) 维 \(k\) 表示 \(V\)，证明自然同构
\[
\Lambda^{d-1}V\simeq V^*\otimes\Lambda^dV
\]
与群作用相容。
\end{exercise}
\begin{solution}
将 \(u\in\Lambda^{d-1}V\) 送到 \(v\mapsto v\wedge u\) 的线性映射，得到 \(\Lambda^{d-1}V\to\operatorname{Hom}_k(V,\Lambda^dV)\)。在一组基下，删去第 \(i\) 个基向量的楔积映为带符号的第 \(i\) 个对偶基乘最高楔积，因此双射，任意特征均成立。群作用保持楔积，故映射等变；再用有限维 Hom 与张量的自然识别即得所述同构。
\end{solution}

\begin{exercise}\label{b5:ex:01-cyclic-modular-regular}
设 \(k\) 特征为 \(p\)，\(G=C_p\)。证明 \(k[G]\simeq k[t]/(t^p)\)，并求正则表示的全部子表示。它是否不可分解？
\end{exercise}
\begin{solution}
把 \(t\) 送到 \(r-1\)，由 \((r-1)^p=r^p-1=0\) 得满同态；两侧维数 \(p\)，故同构。子表示就是理想。任一非零理想中取最低非零次数最小的元素 \(t^j u\)，其中 \(u\) 常数项非零，因有限等比求逆而可逆，所以理想含 \(t^j\)，而其他元素都为它的倍数。故全部子表示为 \((t^j)\)、\(0\le j\le p\)。这些子模成一条链，任意两个非零子模有非零交，因此正则表示不可分解；当 \(p>1\) 时它却不是不可约。
\end{solution}

\begin{exercise}\label{b5:ex:01-augmentation-idempotent}
设 \(|G|\) 在 \(k\) 中可逆，\(e=|G|^{-1}\sum_g g\)。证明 \(e\) 为中心幂等元，\(\ker\varepsilon=k[G](1-e)\)，并写出群代数的直和分解。
\end{exercise}
\begin{solution}
左右乘群元素都固定群元素之和，且 \(e^2=e\)，所以 \(e\) 中心。对任意 \(a\)，有 \(ae=\varepsilon(a)e\)。若 \(\varepsilon(a)=0\)，则 \(a=a(1-e)\)；反向 \(\varepsilon(1-e)=0\)，得核等式。因此
\[
k[G]=ke\oplus k[G](1-e).
\]
第一项为平凡表示，第二项为增广理想；正交幂等元同时给出代数直积分解。
\end{solution}

\begin{exercise}\label{b5:ex:01-finite-field-schur}
在 \(\mathbb F_3^2\) 上令 \(C_4\) 生成元作用为 \(A=\begin{pmatrix}0&-1\\1&0\end{pmatrix}\)。证明此表示不可约，并计算自同态环。
\end{exercise}
\begin{solution}
\(A^2=-I,A^4=I\)。多项式 \(T^2+1\) 在 \(\mathbb F_3\) 无根，所以 \(A\) 没有不变直线，表示不可约。解矩阵交换方程得到全部自同态为 \(aI+bA\)。于是
\[
\operatorname{End}_{C_4}(V)\simeq\mathbb F_3[T]/(T^2+1),
\]
是九元域。这个除代数比标量域大，正好显示 Schur 标量结论需要代数闭性。
\end{solution}

\begin{exercise}\label{b5:ex:01-isotypic-diagonal}
设 \(S\) 为有限群 \(G\) 的不可约复表示。分类 \(S\oplus S\) 的不可约子表示，说明为什么两个坐标子空间并非唯一的简单直和项。
\end{exercise}
\begin{solution}
任意不可约子表示至少有一个非零坐标投影，故由 Schur 引理同构于 \(S\)。其嵌入由一对标量 \((a,b)\ne(0,0)\) 给出，即 \(v\mapsto(av,bv)\)。两对标量给同一个像，当且仅当相差一个共同非零倍数，所以这些子表示由射影直线 \(\mathbb P^1(\mathbb C)\) 参数化。两条不同参数对应的子表示交为零，维数相加为 \(2\dim S\)，从而给出许多不同的直和分解。
\end{solution}

\begin{exercise}\label{b5:ex:01-center-faithful}
设有限群 \(G\) 有忠实不可约复表示。证明 \(Z(G)\) 循环。
\end{exercise}
\begin{solution}
中心元素的作用与全群作用交换，故由 Schur 引理都是标量。忠实性使中心到这些非零标量的同态单射。因此中心同构于 \(\mathbb C^\times\) 的有限子群，由\cref{b5:ex:01-faithful-linear}的单位根论证为循环群。
\end{solution}

\begin{exercise}\label{b5:ex:01-invariant-hermitian}
设 \(V\) 为不可约复表示。证明两种不变正定 Hermitian 内积仅相差一个正实数倍。
\end{exercise}
\begin{solution}
固定第一种内积，存在唯一线性算子 \(A\)，使第二种写成 \((v,Aw)_1\)。两内积均不变，故 \(A\) 与群作用交换；Schur 引理使 \(A=cI\)。在任意非零向量上比较两个正实范数平方，得到 \(c>0\)。因此两内积相差正实数倍。
\end{solution}

\begin{exercise}\label{b5:ex:01-tensor-weights}
设 \(C_5=\langle r\rangle\)，\(\zeta\) 为本原五次单位根。若 \(V\) 的生成元矩阵为 \(\operatorname{diag}(\zeta,\zeta^{-1})\)，分解 \(V\otimes V\)、\(\Lambda^2V\) 及 \(V^*\)。
\end{exercise}
\begin{solution}
取特征基 \(v_+,v_-\)，张量基的权分别为 \(\zeta^2,1,1,\zeta^{-2}\)。因此 \(V\otimes V=\lambda_2\oplus1^{\oplus2}\oplus\lambda_{-2}\)。最高楔积的权为一，所以 \(\Lambda^2V\) 平凡。对偶的两个权互换，故 \(V^*\simeq V\)。这些分解来自具体权空间，不以维数吻合代替同构。
\end{solution}

\begin{exercise}\label{b5:ex:01-polynomial-average}
设 \(\zeta\in\mathbb C\) 为本原三次单位根，\(G=C_3=\langle s\rangle\) 在 \(V=\mathbb C^2\) 上的生成元矩阵为 \(\operatorname{diag}(\zeta,\zeta^{-1})\)。以 \(x,y\) 为对偶坐标，对整数 \(a,b\ge0\) 计算 \(\mathcal R(x^ay^b)\)，并求全次数 \(0\le d\le4\) 的不变齐次多项式空间的基。另将全部二次齐次多项式组成的表示分解为一维表示。
\end{exercise}
\begin{solution}
对偶作用为 \(s(x)=\zeta^{-1}x\)、\(s(y)=\zeta y\)，所以
\[
\mathcal R(x^ay^b)=
\frac{1+\zeta^{b-a}+\zeta^{2(b-a)}}3\,x^ay^b
=
\begin{cases}
x^ay^b,&a\equiv b\pmod3,\\
0,&a\not\equiv b\pmod3.
\end{cases}
\]
后一情形用到三次单位根之和 \(1+\zeta+\zeta^2=0\)。单项式是各齐次部分的基，且均为特征向量，因而按同余条件选出的单项式给出不动空间的全部基。次数依次为零至四时，可以取
\[
\{1\},\qquad \varnothing,\qquad \{xy\},\qquad
\{x^3,y^3\},\qquad \{x^2y^2\}.
\]
全部二次多项式以 \(x^2,xy,y^2\) 为基，生成元在这三个向量上的特征值依次为 \(\zeta,1,\zeta^2\)。因此它是三个不同一维表示的直和，不动项恰为 \(\mathbb Cxy\)。这些有限次数的计算尚未单独证明整个不变环的生成定理。
\end{solution}

\begin{exercise}\label{b5:ex:01-cyclic-unipotent}
设 \(k\) 特征为 \(p\)，\(1\le d\le p\)，\(N\) 为一个 \(d\) 阶幂零 Jordan 块。证明 \(r\mapsto I+N\) 给出 \(C_p\) 的不可分解表示，并求其不动空间维数。
\end{exercise}
\begin{solution}
\((I+N)^p=I+N^p=I\)，所以群关系成立。不动空间是 \(\ker N\)，维数一。若表示是两个非零子表示的直和，\(N\) 在每一项上幂零，故每项都有非零核，两核直和会使 \(\dim\ker N\ge2\)，矛盾。因此不可分解；\(d=1\) 时还不可约，\(d>1\) 时有真子表示 \(\ker N\)。
\end{solution}

% !TeX root = ../../Book5.tex
% 纲要标题：## 第2章　特征标与正交关系
% 纲要位置：工作记录/计划纲要.md，第 3761 行。


\chapter{特征标与正交关系}
\label{b5:ch:02}
\BookChapterNumber{b5:ch:02}

\begin{chapterabstract}
本章用迹构造复特征标，证明正交关系与不可约特征标的完备性，并从具体表示计算若干小群的特征标表及其信息边界。
\end{chapterabstract}

\begin{learninggoals}
\begin{itemize}
\item 计算直和、张量、对偶和外平方的特征标，始终核对工作域。
\item 用内积计算 Hom 空间维数与分解重数，理解列正交所需的完备性。
\item 从特征标构造中心幂等元，并解释群代数的矩阵块。
\item 实际建立小群的特征标表，区分数值约束、存在性和群同构问题。
\end{itemize}
\end{learninggoals}

一个复表示给每个群元素指定一张矩阵，但同时换基会改变这些矩阵的外观。迹却不随换基改变，因此可以先问：只知道每个群元素的迹，究竟能恢复多少表示信息？对交换坐标的二元群，恒等元与交换元的迹分别是 \(2\) 和 \(0\)；若把这个表示分成平凡表示与符号表示，两者的迹相加正好给出这两个数。一般群中，这种「由迹读出分解」需要正交关系保证，还须确认已经找到全部不可约表示，才能正确填写特征标表。

\ProseParagraphBreak
本章只讨论有限维复表示。上一章的 Maschke 定理和 Schur 引理是主要工具；所需内积及共轭约定在本章固定。表格应从具体表示和正交计算得出，而不能把一组满足部分数值关系的行直接宣布为已经存在的表示。

\ProseParagraphBreak
Etingof 等人的《Introduction to Representation Theory》\cite[Sections 4.5--4.8, pp. 66--73]{EtingofEtAl2011RepresentationTheory}提供正交关系、酉表示、矩阵系数与小群特征标表的集中叙述。本章内积采用第一变量共轭线性；不同文献把共轭放在另一变量时，Hom 公式中的次序要相应转换。

% !TeX root = ../../Book5.tex
% 纲要标题：### 2.1 特征标
% 纲要位置：工作记录/计划纲要.md，第 3763 行。
% 纲要细目（写作范围）：
% * 迹、类函数与表示同构
% * 特征标在直和、张量和对偶下的行为
% * 本章默认复表示


\section{特征标}
\label{b5:sec:0201}

一个表示给每个群元素一个矩阵，但矩阵条目随换基而改变。迹在相似变换下不变，又与直和和张量积相容，因此可以把整组矩阵压缩成一个群上的函数。本章固定有限群的有限维复表示；正特征下，不能仅凭这些迹确定表示的同构类型。


\subsection{迹、类函数与表示同构}
\label{b5:subsec:0201:01}

\begin{definition}\label{b5:def:character}
设 \(G\) 为有限群，\(\rho:G\to\operatorname{GL}_{\mathbb C}(V)\) 为有限维复表示。函数
\[
\chi_V:G\longrightarrow\mathbb C,\qquad
\chi_V(g)=\operatorname{tr}\rho(g)
\]
称为 \(V\) 的\term[tezhengbiao]{特征标}[character]。非零不可约表示的特征标称为不可约特征标。
\end{definition}

有 \(\chi_V(1)=\dim V\)。因为 \(\rho(hgh^{-1})=\rho(h)\rho(g)\rho(h)^{-1}\)，迹的相似不变性使特征标在每个共轭类上取常值。群上在共轭类上常值的复函数称为\term[lei-hanshu]{类函数}[class function]；它们组成空间 \(\operatorname{CF}(G)\)，维数等于共轭类个数。

\ProseParagraphBreak
同构表示的特征标相同，但逆命题此时还没有得到。只有先证明不可约特征标的正交性，才能从特征标中读出不可约分解的重数，继而恢复表示的同构类型。不能把这项结论误当成迹的定义的一部分。

\subsection{直和、张量与对偶的特征标}
\label{b5:subsec:0201:02}

\begin{proposition}\label{b5:prop:character-operations}
设 \(G\) 为有限群，\(V,W\) 为 \(G\) 的有限维复表示，\(g\in G\)。则
\[
\chi_{V\oplus W}=\chi_V+\chi_W,\qquad
\chi_{V\otimes W}=\chi_V\chi_W,\qquad
\chi_{V^*}(g)=\chi_V(g^{-1})=\overline{\chi_V(g)}.
\]
若 \(G\) 作用于有限集合 \(X\)，则对每个 \(g\in G\)，相应复置换表示 \(\mathbb C[X]\) 的特征标为 \(g\) 在 \(X\) 中的不动点数。
\end{proposition}
\begin{proof}
直和矩阵分块对角，迹相加；张量矩阵的对角元为两矩阵对角元的两两乘积，故迹相乘。对偶作用矩阵为逆矩阵的转置，所以其迹为 \(\chi_V(g^{-1})\)。若 \(g\) 的阶为 \(m\)，\(\rho(g)\) 满足 \(T^m-1\)，这个复多项式无重根，故矩阵可对角化，特征值均为单位根；它们的逆正是复共轭，给出最后等式。置换矩阵的对角元仅在基元素被固定时为一，其余为零，所以迹等于不动点数。
\end{proof}

例如 \(S_3\) 在三个字母上的置换特征标，在单位元、换位、三轮换上依次取 \(3,1,0\)。其不动直线的特征标恒为一，和为零的二维补表示因而有值 \(2,0,-1\)。下一节将从这些数值确认该二维表示不可约。

\ProseParagraphBreak
外平方与对称平方也能从原特征标计算：
\begin{equation}\label{b5:eq:symmetric-exterior-character}
\chi_{\Lambda^2V}(g)=\frac{\chi_V(g)^2-\chi_V(g^2)}2,\qquad
\chi_{\operatorname{Sym}^2V}(g)=\frac{\chi_V(g)^2+\chi_V(g^2)}2.
\end{equation}
确实，若 \(\rho(g)\) 的特征值为 \(\lambda_i\)，外平方的特征值为 \(i<j\) 的 \(\lambda_i\lambda_j\)，对称平方还包含各 \(\lambda_i^2\)。展开 \((\sum_i\lambda_i)^2\) 即得两式。这些具体公式使用了复数域中二可逆的条件。

\subsection{复表示的基本约定}
\label{b5:subsec:0201:03}

类函数空间采用 Hermitian 内积
\begin{equation}\label{b5:eq:character-inner-product}
\langle f,h\rangle_G=\frac1{|G|}\sum_{g\in G}\overline{f(g)}h(g).
\end{equation}
它在第一变量共轭线性、第二变量线性，与本书向量空间内积的约定一致。若共轭类代表为 \(g_1,\ldots,g_r\)，类大小为 \(c_j\)，则
\[
\langle f,h\rangle_G=\frac1{|G|}\sum_{j=1}^r c_j\overline{f(g_j)}h(g_j).
\]
\symidx{\(\langle f,h\rangle_G\)：有限群类函数的内积}

类函数空间可以包含完全不是特征标的函数；例如一个任意负常数不能是非零表示在单位元处的值。后面的正交基定理会表明，类函数成为实际特征标的条件是不可约基下的系数全部为非负整数。允许整数而不要求非负，则得到虚特征标。

\ProseParagraphBreak
若在特征 \(p\) 下直接取 \(k\) 值迹，平凡表示的 \(p\) 重直和与零表示都给出零函数，却不同构。特征整除群阶时，表示还可能不完全可约。这两种困难不能靠继续套用本章复数内积消除；模表示与 Brauer 特征标在第8章讨论。

% !TeX root = ../../Book5.tex
% 纲要标题：### 2.2 正交关系
% 纲要位置：工作记录/计划纲要.md，第 3769 行。
% 纲要细目（写作范围）：
% * 矩阵系数正交关系
% * 特征标内积与行正交
% * 不可约特征标构成类函数空间的基与列正交


\section{正交关系}
\label{b5:sec:0202}

正交关系来自对线性映射取平均。平均后得到交织算子，Schur 引理迫使它在不同不可约表示之间为零，在相同表示上为标量。选取矩阵单位就把这个结论变成矩阵条目及特征标的求和公式。


\subsection{矩阵系数正交关系}
\label{b5:subsec:0202:01}

选定表示 \(\rho\) 的基后，每个矩阵条目给出函数 \(g\mapsto\rho(g)_{ia}\)，称为 \(\rho\) 的矩阵系数。它属于全函数空间 \(\mathbb C^G\)，以下使用\eqref{b5:eq:character-inner-product}在该空间上定义的内积。

\begin{theorem}[矩阵系数正交]\label{b5:thm:matrix-coefficient-orthogonality}
设 \(G\) 为有限群，\(\rho\) 与 \(\sigma\) 为 \(G\) 的不可约有限维复表示，维数分别为 \(d_\rho,d_\sigma\)。在各自的不变 Hermitian 内积下选定正交归一基，将两表示写成酉矩阵。若它们不同构，则对 \(1\le i,a\le d_\sigma\)、\(1\le j,b\le d_\rho\) 有
\[
\frac1{|G|}\sum_g\sigma(g)_{ia}\overline{\rho(g)_{jb}}=0.
\]
若 \(\rho=\sigma\) 是同一个 \(d\) 维不可约表示，并对两者采用同一组正交归一基，则对 \(1\le i,a,j,b\le d\) 有
\[
\frac1{|G|}\sum_g\rho(g)_{ia}\overline{\rho(g)_{jb}}
=\frac1d\delta_{ij}\delta_{ab}.
\]
\end{theorem}
\begin{proof}
对 \(T:V_\rho\to V_\sigma\) 定义
\[
P(T)=\frac1{|G|}\sum_g\sigma(g)T\rho(g^{-1}).
\]
重排求和表明 \(P(T)\) 是从 \(V_\rho\) 到 \(V_\sigma\) 的交织算子。两表示不同构时，由 Schur 引理得 \(P(T)=0\)。当 \(\rho=\sigma\) 且采用同一组基时，\(P(T)\) 为标量矩阵 \(cI\)。取迹给出 \(cd=\operatorname{tr}T\)，故 \(c=\operatorname{tr}T/d\)。

取矩阵单位 \(T=E_{ab}:V_\rho\to V_\sigma\)，即把 \(V_\rho\) 的第 \(b\) 个基向量送到 \(V_\sigma\) 的第 \(a\) 个基向量，其余基向量送到零。平均式的 \(i,j\) 条目为
\[
P(E_{ab})_{ij}
=\frac1{|G|}\sum_g\sigma(g)_{ia}\rho(g^{-1})_{bj}
=\frac1{|G|}\sum_g\sigma(g)_{ia}\overline{\rho(g)_{jb}},
\]
最后一步使用 \(\rho\) 的酉性。不同构时该条目为零；同一表示、同一组基时，\(\operatorname{tr}E_{ab}=\delta_{ab}\)，该条目为 \(\delta_{ab}\delta_{ij}/d\)，得到两式。
\end{proof}

这还说明同一不可约表示的 \(d^2\) 个矩阵系数线性无关：它们互相正交，且每个范数平方为 \(1/d>0\)。矩阵系数一般不是类函数，只有取对角元之和才得到特征标。

\subsection{特征标内积与行正交}
\label{b5:subsec:0202:02}

\begin{theorem}\label{b5:thm:character-hom-inner-product}
设 \(G\) 为有限群，\(V,W\) 为有限维复表示。则
\[
\langle\chi_V,\chi_W\rangle_G
=\dim_{\mathbb C}\operatorname{Hom}_G(V,W).
\]
特别地，不可约特征标两两正交且范数为一。
\end{theorem}
\begin{proof}
在 \(\operatorname{Hom}_{\mathbb C}(V,W)\) 上，对共轭作用 \(T\mapsto\rho_W(g)T\rho_V(g^{-1})\) 取平均。平均算子 \(P\) 为投影，像恰为不动空间 \(\operatorname{Hom}_G(V,W)\)，所以 \(\operatorname{tr}P\) 等于该空间维数。由\cref{b5:prop:hom-tensor-invariants}及张量迹公式，该作用在 \(g\) 处的迹为
\(\chi_W(g)\chi_V(g^{-1})=\overline{\chi_V(g)}\chi_W(g)\)。
对 \(g\) 求平均即得等式。不可约时，再用\cref{b5:lem:schur}和\cref{b5:cor:schur-algebraically-closed}，同构时 Hom 维数为一，不同构时为零。
\end{proof}

\begin{corollary}\label{b5:cor:character-irreducibility}
设 \(G\) 为有限群，\(V\) 为 \(G\) 的有限维复表示，\(S_i\) 为两两不同构的不可约复表示、\(m_i\) 为非负整数。若 \(V\) 的不可约分解为 \(\bigoplus_i S_i^{\oplus m_i}\)，则
\[
\langle\chi_{S_i},\chi_V\rangle_G=m_i,\qquad
\langle\chi_V,\chi_V\rangle_G=\sum_i m_i^2.
\]
因此非零 \(V\) 不可约当且仅当后一数为一；两个复表示同构当且仅当特征标相同。
\end{corollary}
\begin{proof}
Maschke 定理保证存在分解；迹在直和上相加，所以 \(\chi_V=\sum_i m_i\chi_{S_i}\)。逐项用刚证的正交关系即可得到公式。重数唯一确定分解的同构类型，而 \(\sum_i m_i^2=1\) 恰好表示只有一个重数为一，其他为零。
\end{proof}

这些是特征标表的行正交关系。列正交还需要知道不可约特征标已经填满类函数空间，不能仅从一组正交向量直接假定它们是基。

\subsection{不可约特征标与类函数空间的基}
\label{b5:subsec:0202:03}

\begin{theorem}\label{b5:thm:irreducible-character-basis}
设 \(G\) 为有限群。其全部不可约复特征标组成复类函数空间 \(\operatorname{CF}(G)\) 关于特征标内积 \(\langle\ ,\ \rangle_G\) 的标准正交基。因而不可约复表示的同构类型数等于共轭类数。
\end{theorem}
\begin{proof}
由\cref{b5:lem:irreducibles-in-regular}，不可约类型只有有限个；正交性已经给出线性无关。若类函数 \(f\) 与全部不可约特征标正交，令
\[
z_f=\sum_{g\in G}f(g^{-1})g\in\mathbb C[G].
\]
因系数为类函数，\(z_f\) 位于群代数中心。它在任一不可约 \(S_i\) 上由 Schur 引理作用为 \(a_iI\)，而
\[
(\dim S_i)a_i=\sum_g f(g^{-1})\chi_i(g)
=|G|\langle\chi_i,f\rangle_G=0.
\]
故在每个不可约上作用为零。正则表示是这些不可约的直和，所以 \(z_f\) 在正则表示上也为零；作用于单位基向量得到 \(z_f=0\)，于是所有系数 \(f(g^{-1})=0\)。所以正交补为零，给定正交组确为整个类函数空间的基。
\end{proof}

\begin{theorem}[列正交]\label{b5:thm:row-column-orthogonality}
设 \(G\) 为有限群，全部不可约复特征标为 \(\chi_1,\ldots,\chi_r\)。对 \(g,h\in G\)，有
\[
\sum_{i=1}^r\chi_i(g)\overline{\chi_i(h)}
=
\begin{cases}
|C_G(g)|,&g,h\text{ 共轭},\\
0,&g,h\text{ 不共轭}.
\end{cases}
\]
\end{theorem}
\begin{proof}
取各共轭类代表 \(g_1,\ldots,g_r\)，令类大小 \(c_j=|G|/|C_G(g_j)|\)。上一基定理保证不可约特征标数等于共轭类数，所以
\[
U_{ij}=\sqrt{\frac{c_j}{|G|}}\,\chi_i(g_j)
\]
定义一个方阵。行正交给出 \(U\overline U^{\mathsf T}=I\)。方阵的右逆也为左逆，故 \(\overline U^{\mathsf T}U=I\)。写出其第 \(j,k\) 个条目，得到
\[
\delta_{jk}
=\frac{\sqrt{c_jc_k}}{|G|}
\sum_{i=1}^r\overline{\chi_i(g_j)}\chi_i(g_k).
\]
当 \(j=k\) 时，这给出
\[
\sum_{i=1}^r|\chi_i(g_j)|^2
=\frac{|G|}{c_j}=|C_G(g_j)|.
\]
当 \(j\ne k\) 时，前面的实系数为正，故求和为零；对该零等式取复共轭，得到
\[
\sum_{i=1}^r\chi_i(g_j)\overline{\chi_i(g_k)}=0.
\]
特征标取值与中心化子阶都在共轭类上不变，因此这两式推广到任意 \(g,h\) 就是所述结论。
\end{proof}

行正交计算表示重数，列正交计算中心化子阶。单位元的一列特别给出
\[
\sum_i\chi_i(1)^2=|G|.
\]
这一平方和也可以从正则表示的具体分解得到，下一节将保留那个更有结构的信息。

% !TeX root = ../../Book5.tex
% 纲要标题：### 2.3 群代数的分解
% 纲要位置：工作记录/计划纲要.md，第 3775 行。
% 纲要细目（写作范围）：
% * 正则特征标与不可约表示维数的平方和
% * 复群代数的矩阵块分解
% * 中心幂等元及块的半单情形


\section{群代数的分解}
\label{b5:sec:0203}

正则表示包含全部不可约类型，并且它的迹很容易计算。两者结合，把不可约次数与群阶联系起来。随后再用中心幂等元从群代数中直接取出各个等型分量，说明特征标不仅记录数字，也能构造投影。


\subsection{正则特征标}
\label{b5:subsec:0203:01}

\begin{theorem}\label{b5:thm:regular-decomposition}
设 \(G\) 为有限群，\(\chi_1,\ldots,\chi_r\) 为全部不可约复特征标，次数为 \(d_i=\chi_i(1)\)，相应表示为 \(S_i\)。正则特征标满足
\[
\chi_{\rm reg}(g)=
\begin{cases}|G|,&g=1,\\0,&g\ne1,\end{cases}
\qquad
\mathbb C[G]\simeq\bigoplus_iS_i^{\oplus d_i}.
\]
特别地，\(\sum_i d_i^2=|G|\)。
\end{theorem}
\begin{proof}
左平移 \(x\mapsto gx\) 只有在 \(g=1\) 时有不动点，所以置换特征标给出第一式。由\cref{b5:cor:character-irreducibility}，\(S_i\) 在正则表示中的重数为
\[
\langle\chi_i,\chi_{\rm reg}\rangle_G
=\overline{\chi_i(1)}=d_i.
\]
于是得到分解。比较两侧维数即得平方和。
\end{proof}

如果已经构造了一组两两不同构的不可约表示，且它们的次数平方和等于群阶，那么列表已经完备：任何遗漏的不可约还会给正则分解增加一个正的平方。反过来，平方和吻合本身不能证明候选表示不可约，也不能排除列表重复；这两项需要先核验。

\subsection{复群代数的矩阵块分解}
\label{b5:subsec:0203:02}

\begin{theorem}\label{b5:thm:complex-group-algebra}
设 \(G\) 为有限群，\(S_1,\ldots,S_r\) 为全部不同构的不可约有限维复表示，\(d_i=\dim_{\mathbb C}S_i\)，记 \(\rho_i\) 为 \(S_i\) 的作用及其在 \(\mathbb C[G]\) 上的线性延伸。由作用定义的复代数同态
\[
\Phi:\mathbb C[G]\longrightarrow\prod_{i=1}^r
\operatorname{End}_{\mathbb C}(S_i),\qquad
a\longmapsto\bigl(\rho_i(a)\bigr)_i
\]
为同构。因此 \(\mathbb C[G]\simeq\prod_i M_{d_i}(\mathbb C)\)。
\end{theorem}
\begin{proof}
各作用保持乘法和单位，故是代数同态。若 \(a\) 的全部像为零，它零化每个不可约，因而由正则分解也零化 \(\mathbb C[G]\)；作用于一得到 \(a=0\)。所以 \(\Phi\) 单射。两侧复维数分别为 \(|G|\) 和 \(\sum_i d_i^2\)，由\cref{b5:thm:regular-decomposition}相等，有限维线性代数使它满射。
\end{proof}

不选 \(S_i\) 的基时，右端写成自同态代数，作用同态本身有明确意义；写成具体矩阵代数则需要选基。一般底域下可能出现非平凡除代数，正特征整除群阶时还会失去半单性，所以这个复群代数公式不能原样用于所有域。

\ProseParagraphBreak
这里的 \(\operatorname{End}_{\mathbb C}(S_i)\) 与 Schur 引理中的 \(\operatorname{End}_G(S_i)\) 是不同的代数：前者容许所有线性算子，后者只容许与群作用交换的算子。例如，\(S_3\) 的二维标准表示 \(S\) 是三点置换空间中坐标和为零的子空间；它的特征标在三个共轭类上为 \(2,0,-1\)，范数平方为一，故不可约。于是 \(\operatorname{End}_G(S)=\mathbb C I\)，而刚证的同构使群代数在 \(S\) 上的像为全部 \(M_2(\mathbb C)\)。这个四维矩阵块作为左模也不是一个简单模：每个矩阵唯一分成只有第一列非零与只有第二列非零的两个矩阵，左乘保持这两部分，并在每一列上按 \(S\) 作用。因此它作为左 \(G\)-模为 \(S\oplus S\)，对应正则分解中的重数二。

\subsection{中心幂等元及块的半单情形}
\label{b5:subsec:0203:03}

\begin{theorem}\label{b5:thm:central-character-idempotents}
设 \(G\) 为有限群，\(S_1,\ldots,S_r\) 为全部不同构的不可约有限维复表示，其特征标为 \(\chi_i\)，\(d_i=\chi_i(1)\)。令
\begin{equation}\label{b5:eq:central-idempotent}
e_i=\frac{d_i}{|G|}\sum_{g\in G}\chi_i(g^{-1})g.
\end{equation}
则 \(e_i\in\mathbb C[G]\) 是两两正交的中心幂等元，\(\sum_i e_i=1\)，并且在 \(G\) 的任意有限维复表示 \(V\) 中，\(e_iV\) 正是类型 \(S_i\) 的等型分量。理想 \(\mathbb C[G]e_i\) 同构于 \(M_{d_i}(\mathbb C)\)。
\end{theorem}
\begin{proof}
因系数在共轭类上常值，\(e_i\) 位于中心。在不可约 \(S_j\) 上，其作用为标量，取迹得到
\[
\operatorname{tr}\rho_j(e_i)
=d_i\langle\chi_i,\chi_j\rangle_G=d_i\delta_{ij}.
\]
除以 \(d_j\) 可知它在 \(S_i\) 上为一，在其他 \(S_j\) 上为零。借\cref{b5:thm:complex-group-algebra}的单射性，作用上的平方、相互乘积与和的等式在群代数中也成立。对一般 \(V\) 用不可约分解，得到等型投影。最后 \(\Phi(e_i)\) 仅在第 \(i\) 个矩阵因子为单位，故相应理想恰是那个因子。
\end{proof}

不能再写成两个非零正交中心幂等元之和的非零中心幂等元，称为本原中心幂等元；由上面的矩阵乘积分解，各 \(e_i\) 恰为全部这类元素。它们对应的直和因子称为群代数的\term[qun-daishu-kuai]{块}[block of a group algebra]。在当前复半单情形，每个块只有一个不可约类型；模特征情形的块可能包含多个简单类型，不能把当前特点写成块的一般定义。

\ProseParagraphBreak
中心还可直接用共轭类和描述。对每个共轭类 \(C\)，令 \(z_C=\sum_{g\in C}g\)。一个群代数元素与全部群元素交换，当且仅当其系数在各共轭类上相同，所以这些 \(z_C\) 构成中心的一组基。由矩阵乘积分解，本原中心幂等元 \(\{e_i\}\) 构成的基与共轭类和基 \(\{z_C\}\) 给出中心的两种描述；第五章将进一步考察类和乘法的整数系数。

% !TeX root = ../../Book5.tex
% 纲要标题：### 2.4 特征标表
% 纲要位置：工作记录/计划纲要.md，第 3781 行。
% 纲要细目（写作范围）：
% * 循环群、二面体群和四元数群
% * S_3、S_4 与 A_4 的计算
% * 由表能读出什么、不能读出什么
% ---


\section{特征标表}
\label{b5:sec:0204}

计算特征标表时，先确定共轭类，再从具体表示获得候选行，用内积确认不可约性，最后用次数平方和检查是否遗漏。本节逐个完成这些步骤。所有表示均在复数域上，表中同时列出共轭类大小，避免把普通欧氏内积误作特征标内积。


\subsection{循环群、二面体群与四元数群}
\label{b5:subsec:0204:01}

对 \(C_m=\langle r\rangle\)，群交换，所有不可约表示一维。令 \(\zeta=\exp(2\pi\mathrm i/m)\)，则
\[
\chi_j(r^a)=\zeta^{ja},\qquad 0\le j,a<m.
\]
这些行互异，已经有 \(m\) 个，次数平方和为 \(m\)，所以列表完整。正交关系也可由等比求和 \(\sum_{a=0}^{m-1}\zeta^{(i-j)a}=m\delta_{ij}\) 直接看到。

\ProseParagraphBreak
令 \(D_{2n}=\langle r,s\mid r^n=s^2=1,\ srs=r^{-1}\rangle\)，其中 \(n\ge3\)，下标为群阶。取 \(\zeta=\exp(2\pi\mathrm i/n)\)。对 \(j\in\mathbb Z/n\mathbb Z\)，定义
\[
\rho_j(r)=\begin{pmatrix}\zeta^j&0\\0&\zeta^{-j}\end{pmatrix},
\qquad
\rho_j(s)=\begin{pmatrix}0&1\\1&0\end{pmatrix}.
\]
矩阵满足群关系。若 \(2j\not\equiv0\pmod n\)，两个旋转特征值不同，而反射交换对应特征直线，所以表示不可约。其特征标为
\[
\chi_j(r^a)=\zeta^{ja}+\zeta^{-ja}
=2\cos(2\pi ja/n),\qquad \chi_j(r^as)=0.
\]
交换两个基向量给出 \(\rho_j\simeq\rho_{-j}\)。反过来，若 \(\rho_j\simeq\rho_\ell\)，则旋转矩阵的特征值集合相同，故 \(\ell\equiv\pm j\pmod n\)；除此之外没有重复。若 \(2j\equiv0\pmod n\)，旋转矩阵为标量矩阵，两条直线 \(\mathbb C(1,1)\)、\(\mathbb C(1,-1)\) 均不变，故 \(\rho_j\) 可约。奇数 \(n\) 时只有 \(j=0\) 属于这种情形，偶数 \(n\) 时则有 \(j=0,n/2\)。一维表示由 \(r\mapsto a,s\mapsto b\) 给出，关系要求 \(a^n=1,a^2=b^2=1\)。因此当 \(n\) 奇数时有两个一维表示和 \((n-1)/2\) 个二维不可约表示；当 \(n\) 偶数时有四个一维表示和 \((n-2)/2\) 个二维不可约表示。两种情况下次数平方和均为 \(2n\)，所以穷尽不可约类型。

\ProseParagraphBreak
旋转的共轭类为 \(\{r^a,r^{-a}\}\)，相同元素时只算一个。反射的指数在共轭下相差偶数；因此 \(n\) 奇数时反射成一个类，\(n\) 偶数时分成偶指数与奇指数两类。这给出将上述公式排成完整特征标表所需的全部列。

\ProseParagraphBreak
四元数群 \(Q_8=\{\pm1,\pm\mathbf i,\pm\mathbf j,\pm\mathbf k\}\) 的四个一维表示来自商 \(Q_8/\{\pm1\}\simeq C_2\times C_2\)。另一个表示由
\[
\rho(\mathbf i)=\begin{pmatrix}\mathrm i&0\\0&-\mathrm i\end{pmatrix},
\qquad
\rho(\mathbf j)=\begin{pmatrix}0&1\\-1&0\end{pmatrix}
\]
给出。它们平方均为 \(-I\)，且反交换，满足四元数关系；前一个的特征直线被后一个交换，故不可约。次数平方和 \(4+4=8\)，列表完整。\Cref{b5:tab:d8-q8}将四元数群与正方形对称群的特征标并排比较。表中的一维特征标按生成元取值编号：对 \(a,b\in\{1,-1\}\)，在 \(D_8\) 中令 \(\lambda_{a,b}(r)=a\)、\(\lambda_{a,b}(s)=b\)；在 \(Q_8\) 中则选生成元 \(\mathbf i,\mathbf j\)，规定 \(\mathbf k=\mathbf i\mathbf j\)，令 \(\lambda_{a,b}(\mathbf i)=a\)、\(\lambda_{a,b}(\mathbf j)=b\)。两群中这些表示都把第二列的中心元素送到 \(1\)，按各自表列的代表顺序，取值均为 \((1,1,a,b,ab)\)。

\begin{table}[htbp]
\centering
\caption{八阶二面体群与四元数群的共同特征标数值}
\label{b5:tab:d8-q8}
\begin{tabular}{lrrrrr}
\toprule
\(D_8\) 代表 & \(1\)&\(r^2\)&\(r\)&\(s\)&\(rs\)\\
\(Q_8\) 代表 & \(1\)&\(-1\)&\(\mathbf i\)&\(\mathbf j\)&\(\mathbf k\)\\
类大小 & \(1\)&\(1\)&\(2\)&\(2\)&\(2\)\\
\midrule
\(\lambda_{1,1}\)&1&1&1&1&1\\
\(\lambda_{1,-1}\)&1&1&1&-1&-1\\
\(\lambda_{-1,1}\)&1&1&-1&1&-1\\
\(\lambda_{-1,-1}\)&1&1&-1&-1&1\\
\(\chi\)&2&-2&0&0&0\\
\bottomrule
\end{tabular}
\end{table}

\subsection{\texorpdfstring{\(S_3\)、\(S_4\) 与 \(A_4\) 的特征标表}{S3、S4 与 A4 的特征标表}}
\label{b5:subsec:0204:02}

对 \(S_3\)，共轭类由轮换型给出，大小为 \(1,3,2\)。平凡与符号表示给出两行；三点置换表示减去平凡直和项，得到次数二的表示，特征标为 \(2,0,-1\)。其范数平方为 \((4+2)/6=1\)，故不可约，且平方和 \(1+1+4=6\) 已穷尽，得到\cref{b5:tab:s3}。

\begin{table}[htbp]
\centering
\caption{\(S_3\) 的特征标表}\label{b5:tab:s3}
\begin{tabular}{lrrr}
\toprule
代表&\(1\)&\((12)\)&\((123)\)\\
类大小&1&3&2\\
\midrule
\(1\)&1&1&1\\
\(\mathrm{sgn}\)&1&-1&1\\
\(\chi\)&2&0&-1\\
\bottomrule
\end{tabular}
\end{table}

对 \(S_4\)，五种轮换型的类大小为 \(1,6,3,8,6\)。四点置换表示减去平凡项，得到标准表示 \(W\)，其特征标为 \(3,1,-1,0,-1\)。直接按类大小计算范数平方为一。与符号表示张量得到另一不同的次数三表示。

\ProseParagraphBreak
集合 \(\{1,2,3,4\}\) 分成两个无序二元组共有三种方式，记为 \(A=12|34\)、\(B=13|24\)、\(C=14|23\)。换位 \((12)\) 固定 \(A\) 并交换 \(B,C\)；双换位 \((12)(34)\) 固定全部三种配对；三轮换 \((123)\) 将 \(A,C,B\) 依次循环；四轮换 \((1234)\) 交换 \(A,C\) 并固定 \(B\)。像群含有一个换位和一个三轮换，故 \(S_4\) 在这三种配对上的作用给出满同态 \(S_4\to S_3\)，核为
\[
\bigl\{1,(12)(34),(13)(24),(14)(23)\bigr\}.
\]
拉回复次数二标准表示仍不可约：一个拉回不变子空间也是整个像群 \(S_3\) 的不变子空间。在表列的五种轮换型上，固定配对数依次为 \(3,1,3,0,1\)。配对置换空间中的常向量构成一维平凡直和项，去掉它便将各迹减去一，得到这一行的值 \(2,0,2,-1,0\)。加上平凡、符号两行，平方和 \(1+1+9+9+4=24\)，得到全部不可约，见\cref{b5:tab:s4}。

\begin{table}[htbp]
\centering
\caption{\(S_4\) 的特征标表}\label{b5:tab:s4}
\begin{tabular}{lrrrrr}
\toprule
代表&\(1\)&\((12)\)&\((12)(34)\)&\((123)\)&\((1234)\)\\
类大小&1&6&3&8&6\\
\midrule
\(1\)&1&1&1&1&1\\
\(\mathrm{sgn}\)&1&-1&1&1&-1\\
\(\chi_W\)&3&1&-1&0&-1\\
\(\chi_W\mathrm{sgn}\)&3&-1&-1&0&1\\
\(\psi\)&2&0&2&-1&0\\
\bottomrule
\end{tabular}
\end{table}

对 \(A_4\)，三个双换位组成一个类。三轮换在 \(A_4\) 中的中心化子仍为它生成的三阶群，所以每类大小四，八个三轮换分成两个类。一个三轮换与其逆在 \(A_4\) 中不共轭：在 \(S_4\) 中把 \((123)\) 变为逆的元素都在奇置换 \((23)\) 乘其三阶中心化子的陪集中。商 \(A_4/V_4\simeq C_3\) 给出三个一维特征标。\(S_4\) 标准表示限制后取值 \(3,-1,0,0\)，范数平方为 \((9+3)/12=1\)，故仍不可约。平方和为 \(3+9=12\)，见\cref{b5:tab:a4}，其中 \(\omega=\exp(2\pi\mathrm i/3)\)。

\begin{table}[htbp]
\centering
\caption{\(A_4\) 的特征标表}\label{b5:tab:a4}
\begin{tabular}{lrrrr}
\toprule
代表&\(1\)&\((12)(34)\)&\((123)\)&\((132)\)\\
类大小&1&3&4&4\\
\midrule
\(1\)&1&1&1&1\\
\(\lambda\)&1&1&\(\omega\)&\(\omega^2\)\\
\(\lambda^2\)&1&1&\(\omega^2\)&\(\omega\)\\
\(\theta\)&3&-1&0&0\\
\bottomrule
\end{tabular}
\end{table}

\subsection{特征标表的信息与局限}
\label{b5:subsec:0204:03}

从次数列及其平方和可以恢复群阶；列的范数给出中心化子阶，从而给出共轭类大小。一维表示的数目等于交换化群 \(G^{\mathrm{ab}}=G/[G,G]\) 的阶：复群代数 \(\mathbb C[G^{\mathrm{ab}}]\) 的全部不可约表示都一维，再用正则分解即可得到这一计数。特征标表也决定任意两个已知表示的张量积分解，只需逐点相乘后与各不可约行取内积。

\begin{proposition}\label{b5:prop:character-kernel}
设 \(G\) 为有限群，\(\rho:G\to\operatorname{GL}_{\mathbb C}(V)\) 为 \(d\) 维复表示，其特征标为 \(\chi\)。则
\[
\ker\rho=\bigl\{g\in G\mid\chi(g)=d\bigr\}.
\]
\end{proposition}
\begin{proof}
一个方向直接取迹。反过来，\(\rho(g)\) 可对角化，特征值 \(\lambda_1,\ldots,\lambda_d\) 均在单位圆上。若其和为正实数 \(d\)，则
\(\sum_i\operatorname{Re}\lambda_i=d\)，而每项至多为一，所以每项实部都为一，即所有 \(\lambda_i=1\)。可对角化性便使 \(\rho(g)=I\)。零表示的核为全群，该公式也成立。
\end{proof}

但普通特征标表不能确定群到同构。\Cref{b5:tab:d8-q8}中的两群有相同表，\(D_8\) 却有五个阶二元素，而 \(Q_8\) 只有一个。一个表的列给出共轭类上的函数值，不自动包含各类取平方后落到哪一列的全部信息。若要计算\eqref{b5:eq:symmetric-exterior-character}中的平方类，就须再有这些群结构数据。


\begin{chaptersummary}
迹与群的共轭作用相容，所以特征标落在类函数空间中。平均交织算子给出正交性，中心群代数元素的作用再证明完备性。不可约特征标因此构成一组可用于计算的基，分解重数就是内积，列的范数则给出中心化子阶。

\ProseParagraphBreak
正则表示中每个不可约的重数等于它的次数，这既给出平方和，也给出复群代数的矩阵分解。中心幂等元把数值信息重新变成作用在任意表示上的投影；同一个群可以用共轭类和或用不可约类型组织它的中心。

\ProseParagraphBreak
制作特征标表时，需要计算共轭类大小，证明候选表示不可约，并核对它们是否穷尽全部不可约类型。特征标表能判断复表示同构，却不能单独确定群同构；下一章将从子群构造表示，为得到新的特征标提供比直接猜表更系统的方法。
\end{chaptersummary}

\BookExercises

\begin{exercise}\label{b5:ex:02-trace-bound}
设有限群的 \(d>0\) 维复表示特征标为 \(\chi\)。证明 \(|\chi(g)|\le d\)，且等号成立当且仅当 \(\rho(g)\) 为标量。注意这与 \(\chi(g)=d\) 的核判据不同。
\end{exercise}
\begin{solution}
特征值 \(\lambda_i\) 均在单位圆上，三角不等式给出界。若和的模为 \(d\)，令其辐角为 \(\theta\)，则 \(\sum_i\operatorname{Re}(e^{-\mathrm i\theta}\lambda_i)=d\)，各项至多一，故各项均为一。于是全部特征值相同；\(\rho(g)\) 可对角化，故为标量。反向标量的模必为一，所以等号成立。只有这个标量是一时，才有 \(\chi(g)=d\)。
\end{solution}

\begin{exercise}\label{b5:ex:02-s3-tensor}
令 \(V\) 为 \(S_3\) 的二维不可约复表示。分解 \(V\otimes V\)、\(\Lambda^2V\) 和 \(\operatorname{Sym}^2V\)。
\end{exercise}
\begin{solution}
在单位元、换位、三轮换三列，\(\chi_V=(2,0,-1)\)，所以张量特征标为 \((4,0,1)\)。与三条不可约行取内积，重数均为一，故
\[
V\otimes V\simeq1\oplus\mathrm{sgn}\oplus V.
\]
外平方为行列式表示。二维标准表示的行列式等于三点置换表示的行列式，因为删去的平凡直和项行列式为一，所以为符号表示。于是 \(\operatorname{Sym}^2V\simeq1\oplus V\)。
\end{solution}

\begin{exercise}\label{b5:ex:02-s4-tensor}
令 \(W\) 为 \(S_4\) 的三维标准复表示，\(\psi\) 为经 \(S_4/V_4\simeq S_3\) 拉回的二维不可约表示。分解 \(W\otimes W\)。
\end{exercise}
\begin{solution}
按正文五列顺序，张量特征标为 \((9,1,1,0,1)\)。它等于 \(1+\chi_W+\chi_W\mathrm{sgn}+\psi\)，逐列相加可核验。因此
\[
W\otimes W\simeq1\oplus W\oplus(W\otimes\mathrm{sgn})\oplus\psi.
\]
各项次数相加为九；同构由特征标判据给出，不只是维数比较。
\end{solution}

\begin{exercise}\label{b5:ex:02-a4-tensor}
令 \(\theta\) 为 \(A_4\) 的三维不可约特征标，\(\lambda\) 为非平凡一维特征标。求 \(\theta\lambda\) 和 \(\theta^2\) 的不可约展开。
\end{exercise}
\begin{solution}
\(\theta\) 在两个三轮换类为零，另两列 \(\lambda=1\)，所以 \(\theta\lambda=\theta\)。由 \(\omega+\omega^2=-1\)，逐列有
\[
\theta^2=1+\lambda+\lambda^2+2\theta.
\]
在单位元与双换位处右端分别为九与一，在三轮换处为零，恰为左端。系数为实际分解重数。
\end{solution}

\begin{exercise}\label{b5:ex:02-s4-restriction}
将 \(S_4\) 的五个不可约复表示限制到 \(A_4\)，求其分解。
\end{exercise}
\begin{solution}
平凡和符号表示都限制为平凡表示。两个三维表示在偶置换上取值相同，都限制为 \(\theta\)。二维表示的限制在 \(A_4\) 四列为 \((2,2,-1,-1)\)，正好是 \(\lambda+\lambda^2\)。这些等式结合特征标同构判据给出全部分解。限制可以使不同的不可约表示合并，也可以使一个不可约表示分裂。
\end{solution}

\begin{exercise}\label{b5:ex:02-s3-central-projectors}
在 \(\mathbb C[S_3]\) 中，记 \(c=(123)\)。写出平凡、符号和二维表示的三个中心幂等元，说明二维投影为什么不含换位项。
\end{exercise}
\begin{solution}
令 \(T\) 为三个换位的群代数和。由\eqref{b5:eq:central-idempotent}，
\[
e_1=\tfrac16(1+T+c+c^2),\quad
e_{\rm sgn}=\tfrac16(1-T+c+c^2),\quad
e_V=\tfrac13(2-c-c^2).
\]
它们和为一。二维特征标在换位处为零，所以投影中不出现那些基向量。把 \(e_V\) 作用于三点置换表示时，像为系数和为零的平面，核为常数直线。
\end{solution}

\begin{exercise}\label{b5:ex:02-cyclic-fourier}
对 \(C_m=\langle r\rangle\) 及本原 \(m\) 次单位根 \(\zeta\)，证明
\[
e_j=\frac1m\sum_{a=0}^{m-1}\zeta^{-ja}r^a
\]
满足 \(re_j=\zeta^j e_j\)、\(e_je_\ell=\delta_{j\ell}e_j\)，并求 \(r^a\) 关于这些元素的展开。
\end{exercise}
\begin{solution}
将求和指标平移一位，得到第一式。继而 \(e_je_\ell=m^{-1}\sum_a\zeta^{-ja}\zeta^{\ell a}e_\ell=\delta_{j\ell}e_\ell\)，其中用有限等比求和。又 \(\sum_je_j=1\)，乘以 \(r^a\) 得
\[
r^a=\sum_{j=0}^{m-1}\zeta^{ja}e_j.
\]
这些公式是群代数基与不可约分量坐标之间的有限 Fourier 变换\footnote{傅里叶（Jean Baptiste Joseph Fourier，1768-1830），法国数学家、物理学家，研究热传导并发展三角级数方法。}。
\end{solution}

\begin{exercise}\label{b5:ex:02-matrix-units}
在不可约酉矩阵表示 \(\rho_i\) 中定义
\[
u^{(i)}_{ab}=\frac{d_i}{|G|}\sum_g\rho_i(g^{-1})_{ba}g.
\]
证明它在第 \(i\) 个矩阵块中的像是矩阵单位 \(E_{ab}\)，在其他块为零，并求这些元素的乘法。
\end{exercise}
\begin{solution}
第 \(j\) 个像的 \(u,v\) 条目为 \(d_i|G|^{-1}\sum_g\overline{\rho_i(g)_{ab}}\rho_j(g)_{uv}\)。由矩阵系数正交，结果为 \(\delta_{ij}\delta_{ua}\delta_{vb}\)。群代数作用同态是单射，故矩阵单位的乘法在原代数中成立：
\[
u^{(i)}_{ab}u^{(j)}_{cd}=\delta_{ij}\delta_{bc}u^{(i)}_{ad},
\qquad e_i=\sum_a u^{(i)}_{aa}.
\]
\end{solution}

\begin{exercise}\label{b5:ex:02-kernels-intersection}
证明有限群全部不可约复表示的核之交为单位子群。由此说明任意有限群都有忠实完全可约复表示。
\end{exercise}
\begin{solution}
若 \(g\) 落在全部核中，它在全部不可约上恒等，因而在正则表示上恒等。正则作用忠实，所以 \(g=1\)。有限个不可约表示各取一份的直和，其核恰为这些核之交，故给出忠实完全可约表示；也可以直接取正则表示。
\end{solution}

\begin{exercise}\label{b5:ex:02-center-from-table}
设 \(\chi_i\) 遍历有限群 \(G\) 的不可约复特征标。证明 \(g\in Z(G)\) 当且仅当 \(|\chi_i(g)|=\chi_i(1)\) 对每个 \(i\) 成立。
\end{exercise}
\begin{solution}
中心元素在每个不可约上为标量，故满足模长条件。反向由\cref{b5:ex:02-trace-bound}，模长等式使各 \(\rho_i(g)\) 为标量矩阵。它们与所有 \(\rho_i(h)\) 交换；借群代数矩阵分解的单射性，得到 \(gh=hg\) 对全部 \(h\) 成立，所以 \(g\) 在中心。
\end{solution}

\begin{exercise}\label{b5:ex:02-same-table-symmetric-square}
对 \(D_8\) 与 \(Q_8\) 的二维不可约复表示，分别求对称平方中平凡表示的重数，解释为什么两群普通特征标表相同却仍得到不同答案。
\end{exercise}
\begin{solution}
按\cref{b5:tab:d8-q8}的列顺序，原特征标均为 \((2,-2,0,0,0)\)。在 \(D_8\) 中，旋转类平方为 \(r^2\)，两反射类平方为一；因此对称平方特征标为 \((3,3,-1,1,1)\)，平均为一。在 \(Q_8\) 中，三个非中心类平方均为 \(-1\)，故对称平方特征标为 \((3,3,-1,-1,-1)\)，平均为零。平方类映射不同，而它不是这张普通数值表自身包含的数据。
\end{solution}

\begin{exercise}\label{b5:ex:02-d10-tensors}
对 \(D_{10}\)，记反射符号一维表示为 \(\varepsilon\)，二维不可约为 \(V_1,V_2\)，其旋转特征标分别为 \(\zeta^a+\zeta^{-a}\)、\(\zeta^{2a}+\zeta^{-2a}\)。分解 \(V_1^{\otimes2}\) 和 \(V_1\otimes V_2\)。
\end{exercise}
\begin{solution}
在旋转上，第一乘积为 \(\zeta^{2a}+\zeta^{-2a}+2\)，在反射上为零；常数旋转值二、反射值零恰是 \(1+\varepsilon\)，故 \(V_1^{\otimes2}\simeq V_2\oplus1\oplus\varepsilon\)。第二乘积在旋转上为 \(\zeta^{3a}+\zeta^{-3a}+\zeta^a+\zeta^{-a}\)，因模五有 \(3\equiv-2\)，等于 \(\chi_2+\chi_1\)，反射上也相等。因此 \(V_1\otimes V_2\simeq V_1\oplus V_2\)。
\end{solution}

\begin{exercise}\label{b5:ex:02-s4-two-subsets}
求 \(S_4\) 作用于四元素集合全部二元子集的置换表示的不可约分解。
\end{exercise}
\begin{solution}
固定二元子集的个数在五个共轭类上分别为 \(6,2,2,0,0\)：换位和双换位都固定其对应的两块二元集合，三轮换和四轮换没有固定二元子集。该行等于 \(1+\chi_W+\psi\)，所以置换表示为 \(1\oplus W\oplus V_\psi\)。维数为六，三个重数均为一。
\end{solution}

\begin{exercise}\label{b5:ex:02-integral-values-not-character}
在 \(S_3\) 的三列上给类函数 \(f=(1,0,1)\)。证明它处处取非负整数值，却不是实际复表示的特征标，也不是虚特征标。
\end{exercise}
\begin{solution}
计算得 \(\langle1,f\rangle=\langle\mathrm{sgn},f\rangle=1/2\)，而与二维行的内积为零。因此 \(f=\dfrac12(1+\mathrm{sgn})\)。实际特征标的不可约系数须为非负整数，虚特征标须为整数；两个条件均不满足。函数取值的整数性不能代替基下系数的整数性。
\end{solution}

\begin{exercise}\label{b5:ex:02-product-groups}
设 \(G,H\) 为有限群。证明 \(G\times H\) 的全部不可约复表示为外张量 \(V_i\boxtimes W_j\)，其中 \((g,h)\) 作用为 \(\rho_i(g)\otimes\sigma_j(h)\)。
\end{exercise}
\begin{solution}
外张量特征标为 \(\chi_i(g)\psi_j(h)\)。内积的有限双和分离成两因子，故这些行两两正交且范数一，从而不可约并互不同构。其次
\[
\sum_{i,j}(\dim V_i\dim W_j)^2
=\left(\sum_i(\dim V_i)^2\right)
\left(\sum_j(\dim W_j)^2\right)=|G|\cdot|H|.
\]
平方和判据说明没有遗漏。
\end{solution}

\begin{exercise}\label{b5:ex:02-nonfaithful-irreds}
写出 \(C_2\times C_2\) 的全部不可约复表示，证明没有一个忠实，但其中两个不同的非平凡表示的直和忠实。
\end{exercise}
\begin{solution}
设生成元为 \(a,b\)。四个不可约由 \((a,b)\mapsto(\epsilon,\delta)\)、\(\epsilon,\delta\in\{\pm1\}\) 给出。平凡表示的核为全群，三个非平凡表示的核分别为三个二阶子群，均不忠实。任意两个不同的二阶子群交为单位群，所以对应两个非平凡表示的直和核为一，因而忠实。
\end{solution}

\begin{exercise}\label{b5:ex:02-q8-two-dimensional-block}
在 \(\mathbb C[Q_8]\) 中用 \(z\) 表示群元素 \(-1\) 对应的基向量。证明 \(e=(1-z)/2\) 是二维不可约表示对应的中心幂等元，并写出该块的一组基。
\end{exercise}
\begin{solution}
四个一维表示在 \(z\) 处均为一，二维表示在 \(z\) 处为 \(-I\)。所以 \(e\) 在前四者上为零，在二维表示上为一，正是所求投影。块由 \(e,e\mathbf i,e\mathbf j,e\mathbf k\) 生成，它们展开后分别支撑在互不相交的四对 \(\{g,zg\}\) 上，故线性无关。块维数四，与 \(M_2(\mathbb C)\) 的维数相同；注意公式中的标量负号与基向量 \(z\) 是不同对象。
\end{solution}

\begin{exercise}\label{b5:ex:02-ordered-pairs-s3}
求 \(S_3\) 在互不相同的有序字母对上的置换表示分解，并以不动空间维数核对轨道数。
\end{exercise}
\begin{solution}
六个有序对上作用传递，固定某一有序对的群元素只能同时固定两个字母，从而为单位。因此这是正则作用。也可逐项计算特征标为 \((6,0,0)\)，分解为 \(1+\mathrm{sgn}+2\chi_V\)。平凡表示重数为一，不动空间一维，和只有一个轨道相符。
\end{solution}

% !TeX root = ../../Book5.tex
% 纲要标题：## 第3章　诱导表示与 Frobenius 互反律
% 纲要位置：工作记录/计划纲要.md，第 3789 行。


\chapter{诱导表示与 Frobenius 互反律}
\label{b5:ch:03}
\BookChapterNumber{b5:ch:03}

\begin{chapterabstract}
本章由群代数上的张量构造诱导表示，证明 Frobenius 互反律和诱导特征标公式，并研究传递性、张量恒等式及诱导与余诱导的关系。
\end{chapterabstract}

\begin{learninggoals}
\begin{itemize}
\item 说明诱导张量积中的左右模结构与平衡关系，选取陪集代表写出作用矩阵。
\item 写出 Frobenius 互反同构及其逆，并用它计算诱导、限制中的不可约重数。
\item 证明诱导特征标公式，利用固定点与轨道数计算置换特征标。
\item 从循环子群诱导构造二面体群的不可约表示，辨认可约和重复同构的情形。
\item 对传递性、张量恒等式与诱导—余诱导同构检查良定义、等变性和自然性，说明有限指数在哪里起作用。
\end{itemize}
\end{learninggoals}

设 \(G\) 为有限群，\(H\le G\)，而 \(W\) 是域 \(k\) 上的 \(H\)-表示。想让 \(G\) 中的元素也作用在 \(W\) 所生成的空间上，可以先写形式符号 \(g\otimes w\)，再要求 \(gh\otimes w=g\otimes hw\)。这条关系把同一左陪集中的重复写法识别起来；选定左陪集代表后，所得空间便是若干份 \(W\) 的直和。群元素怎样在这些份之间移动，可由陪集分解直接算出。更重要的是，从这个空间到一个 \(G\)-表示的等变映射，只需在 \(1\otimes W\) 上给出适合的 \(H\)-等变映射。这正是诱导与限制之间互反关系的来源。

\ProseParagraphBreak
本章使用群代数模、线性映射、张量积及前两章的复特征标理论。群代数上的张量积关系和伴随公式将在本章具体推导。代数构造先在任意域上的群代数模之间进行，其中有限维模对应本卷所称的表示；涉及特征标内积、重数与不可约分解时再明确转到复数域。

\ProseParagraphBreak
陪集代表上的数据便于计算矩阵和迹，从 \(1\otimes W\) 出发的映射则便于证明泛性质与自然同构。Etingof 等人的《Introduction to Representation Theory》讨论诱导、诱导特征标与互反律\cite[Sections 5.8--5.10, pp. 105--109]{EtingofEtAl2011RepresentationTheory}，其函数模型与这里张量模型的关系在\secref{b5:sec:0304}说明。

% !TeX root = ../../Book5.tex
% 纲要标题：### 3.1 限制与诱导
% 纲要位置：工作记录/计划纲要.md，第 3791 行。
% 纲要细目（写作范围）：
% * 群代数上的扩张标量
% * 陪集模型
% * 诱导表示的维数与显式矩阵


\section{限制与诱导}
\label{b5:sec:0301}

一个群的作用限制到子群上，不需要增加任何数据。反向的任务却不能仅靠给原空间补写几个矩阵来完成：子群的表示未必能在同一空间上延拓。诱导的办法是允许空间增大，把子群已给定的作用作为关系保留下来，再让整个群自由地作用。本节先把这些关系写成商空间，随后选择陪集代表得到可计算的模型。以下构造适用于任意域，暂不要求群阶在域中可逆。

\subsection{群代数上的扩张标量}
\label{b5:subsec:0301:01}

\begin{definition}\label{b5:def:restriction-induction}
设 \(k\) 为域，\(G\) 为有限群，\(H\le G\)。若 \(V\) 为任意左 \(k[G]\)-模，把作用仅限于 \(H\)，所得左 \(k[H]\)-模称为 \(V\) 的限制，记为 \(\operatorname{Res}_H^G V\)。若 \(W\) 为任意左 \(k[H]\)-模，把 \(k[G]\) 以左右乘法看成 \(\bigl(k[G],k[H]\bigr)\)-双模，则
\[
\operatorname{Ind}_H^G W\coloneqq k[G]\otimes_{k[H]}W,
\qquad g(a\otimes w)=ga\otimes w
\]
给出一个左 \(k[G]\)-模，称为 \(W\) 的诱导模。当 \(V\)、\(W\) 有限维时，这两个构造分别给出\term[xianzhibiaoshi]{限制表示}[restricted representation]与\term[youdaobiaoshi]{诱导表示}[induced representation]。
\end{definition}
\symidx{\(\operatorname{Res}_H^G\)：从 \(G\) 到 \(H\) 的限制}
\symidx{\(\operatorname{Ind}_H^G\)：从 \(H\) 到 \(G\) 的诱导}

这里的张量积可以直接构造如下。先取域上的张量积 \(k[G]\otimes_k W\)，再模去所有
\[
 ah\otimes w-a\otimes hw
 \qquad\bigl(a\in k[G],\ h\in H,\ w\in W\bigr)
\]
生成的线性子空间。群元素生成 \(k[H]\)，所以商空间中的关系也包含 \(ab\otimes w=a\otimes bw\)，其中 \(b\in k[H]\)。左乘 \(g\) 把上述关系送到同样形式的关系，因而作用良定义；\(g_1\bigl(g_2(a\otimes w)\bigr)=(g_1g_2)a\otimes w\)，且单位元作用为恒等。

\ProseParagraphBreak
这一商构造还说明怎样从诱导空间定义线性映射。一个 \(k\)-双线性映射 \(B:k[G]\times W\rightarrow U\) 能唯一写成 \(B(a,w)=F(a\otimes w)\)，当且仅当
\[
 B(ab,w)=B(a,bw)\qquad\bigl(b\in k[H]\bigr).
\]
必要性来自张量关系；充分性来自 \(B\) 所诱导的域上张量积映射在被模去的子空间上为零。后面的同构都将实际检查这条平衡关系。

\ProseParagraphBreak
若 \(u:W\rightarrow W'\) 为 \(k[H]\)-模同态，则 \(a\otimes w\mapsto a\otimes u(w)\) 保持关系，并保持两端的 \(G\) 作用。于是诱导同时给出模与模同态的对应，并保持恒等映射与复合。限制则把一个 \(k[G]\)-模同态原样看成 \(k[H]\)-模同态。这就是后文把两者称为函子的具体含义。

\subsection{陪集模型}
\label{b5:subsec:0301:02}

选择左陪集 \(G/H\) 的一组代表 \(T\)。每个 \(g\in G\) 唯一写成 \(th\)，其中 \(t\in T\)、\(h\in H\)。所以 \(k[G]=\bigoplus_{t\in T}t k[H]\) 是右 \(k[H]\)-模的直和分解；这里使用左陪集，正是因为张量积中群代数位于右模一侧。

\begin{proposition}\label{b5:prop:induced-coset-model}
设 \(k\) 为域，\(G\) 为有限群，\(H\le G\)，\(W\) 为任意左 \(k[H]\)-模，\(T\) 为左陪集 \(G/H\) 的代表集。映射
\[
 \bigoplus_{t\in T}W\longrightarrow\operatorname{Ind}_H^G W,
 \qquad(w_t)_{t\in T}\longmapsto\sum_{t\in T}t\otimes w_t
\]
是 \(k\)-线性同构。在此模型中，若 \(gt=t'h\)，其中 \(t'\in T\)、\(h\in H\)，则 \(g\) 把第 \(t\) 个分量中的 \(w\) 送到第 \(t'\) 个分量中的 \(hw\)。
\end{proposition}
\begin{proof}
把 \(a\in k[G]\) 唯一写为 \(a=\sum_{t\in T}t b_t\)，\(b_t\in k[H]\)。给 \((a,w)\) 指定向量族 \((b_tw)_t\)，得到一个双线性映射。把 \(a\) 换成 \(ab\) 时，各系数变为 \(b_tb\)，所以这与把 \(w\) 换成 \(bw\) 的结果相同。它因而通过张量积，给出所述映射的逆。作用公式来自
\(g(t\otimes w)=gt\otimes w=t'h\otimes w=t'\otimes hw\)。
\end{proof}

代表元的改变只改变坐标，不改变诱导模。例如把 \(t\) 换成同一陪集中的 \(th_t\)，则旧坐标 \(w_t\) 在新代表下应写成 \(h_t^{-1}w_t\)，因为
\(t\otimes w_t=th_t\otimes h_t^{-1}w_t\)。这条变换也解释了为何不能把各陪集上的空间看成带有预先指定基的副本。

\begin{corollary}\label{b5:cor:induction-exact}
设 \(G\) 为有限群，\(H\le G\)，\(k\) 为任意域。诱导函子 \(\operatorname{Ind}_H^G:k[H]\text{-}\mathrm{Mod}\to k[G]\text{-}\mathrm{Mod}\) 与限制函子 \(\operatorname{Res}_H^G:k[G]\text{-}\mathrm{Mod}\to k[H]\text{-}\mathrm{Mod}\) 都保持短正合列，其中各模范畴均取全体左模。诱导还保持直和，并且对每个左 \(k[H]\)-模 \(W\)，\(\operatorname{Ind}_H^G W=0\) 当且仅当 \(W=0\)。
\end{corollary}
\begin{proof}
限制不改变底向量空间和底映射，所以核、像不变。选定 \(T\) 后，诱导一个同态就是在每个 \(t\) 分量上应用该同态；核与像均逐分量计算。因此短正合性保持，直和也可逐分量交换。代表集非空，故诱导空间为零恰好要求每个副本 \(W\) 为零。
\end{proof}

\subsection{诱导表示的维数与显式矩阵}
\label{b5:subsec:0301:03}

若 \(W\) 有限维，由\cref{b5:prop:induced-coset-model}立即得到
\begin{equation}\label{b5:eq:induction-dimension}
 \dim_k\operatorname{Ind}_H^G W=[G:H]\dim_k W.
\end{equation}
更具体地，设 \(T=\{t_1,\ldots,t_m\}\)，\(w_1,\ldots,w_d\) 是 \(W\) 的基。按 \(t_i\otimes w_j\) 排列诱导空间的基。对固定 \(g\)，写
\[
 gt_i=t_{\sigma_g(i)}h_i(g),\qquad h_i(g)\in H.
\]
若 \(\rho\) 为 \(W\) 的矩阵表示，则 \(g\) 的矩阵在第 \(i\) 个块列中仅有一个非零块，位于第 \(\sigma_g(i)\) 个块行，等于 \(\rho\bigl(h_i(g)\bigr)\)。这给出了有限次群乘法和子群矩阵计算构成的具体程序。

\begin{example}\label{b5:exm:s3-cyclic-induction}
在 \(G=S_3=\langle r,s\mid r^3=s^2=1,\ srs=r^{-1}\rangle\) 中，取 \(H=\langle r\rangle\)。设 \(k\) 为域，\(\lambda\in k^\times\) 满足 \(\lambda^3=1\)，并令 \(W=k w\)、\(rw=\lambda w\)。选择左陪集代表 \(1,s\)，诱导表示在 \(1\otimes w,s\otimes w\) 下的矩阵为
\[
 R=\begin{pmatrix}\lambda&0\\0&\lambda^{-1}\end{pmatrix},
 \qquad
 S=\begin{pmatrix}0&1\\1&0\end{pmatrix}.
\]
确实，\(rs=sr^{-1}\)，而 \(s\) 交换两个代表。矩阵满足 \(R^3=S^2=I\) 及 \(SRS=R^{-1}\)，与群关系一致。若 \(k=\mathbb C\)、\(\lambda\ne1\)，则 \(R\) 有两个不同特征值。任何一维不变子空间都须为其坐标特征线，但 \(S\) 交换两条线，故该二维表示不可约。取 \(\lambda=1\) 时则得到 \(S_3/H\) 上的置换表示；在复数域上，它分解为平凡表示与符号表示。
\end{example}

最后注意，“诱导”与“延拓”回答不同问题。延拓要求在原空间上定义整个 \(G\) 的作用；诱导的空间则按\eqref{b5:eq:induction-dimension}扩大。上例中 \(\lambda\ne1\) 的一维 \(H\)-表示不能在一维空间中延拓：若能延拓，标量之间交换，关系 \(srs=r^{-1}\) 将迫使 \(\lambda=\lambda^{-1}\)，与 \(\lambda\) 为三次本原单位根矛盾。

% !TeX root = ../../Book5.tex
% 纲要标题：### 3.2 Frobenius 互反律
% 纲要位置：工作记录/计划纲要.md，第 3797 行。
% 纲要细目（写作范围）：
% * Hom 空间伴随
% * 特征标内积形式
% * 重数计算


\section{Frobenius 互反律}
\label{b5:sec:0302}

诱导空间虽然比原空间大，但从它出发的等变映射不需要逐个指定陪集上的取值。只要知道 \(1\otimes W\) 上的取值，整个映射就由 \(G\) 等变性确定。本节把这件事写成一个自然同构；转到复数域后，它又把诱导表示的分解重数变为子群上的计算。

\subsection{Hom 空间伴随}
\label{b5:subsec:0302:01}

\begin{theorem}[Frobenius 互反律]\label{b5:thm:frobenius-reciprocity}
设 \(k\) 为域，\(G\) 为有限群，\(H\le G\)，\(W\) 为任意左 \(k[H]\)-模，\(V\) 为任意左 \(k[G]\)-模。存在对 \(W,V\) 自然的 \(k\)-线性同构
\begin{equation}\label{b5:eq:frobenius-hom}
 \operatorname{Hom}_{k[G]}(\operatorname{Ind}_H^G W,V)
 \cong
 \operatorname{Hom}_{k[H]}(W,\operatorname{Res}_H^G V).
\end{equation}
它把 \(F\) 送到 \(w\mapsto F(1\otimes w)\)；逆映射把 \(u\) 送到 \(a\otimes w\mapsto a\cdot u(w)\)。
\end{theorem}
\begin{proof}
若 \(F\) 为 \(k[G]\)-模同态，则对 \(h\in H\)，
\[
 F(1\otimes hw)=F(h\otimes w)=hF(1\otimes w),
\]
所以其限制确为 \(k[H]\)-模同态。反之，若 \(u\) 为 \(k[H]\)-模同态，则双线性映射 \((a,w)\mapsto a\cdot u(w)\) 满足
\[
 ab\cdot u(w)=a\cdot u(bw)\qquad\bigl(b\in k[H]\bigr),
\]
故通过张量积；所得 \(F_u\) 满足 \(F_u(ga\otimes w)=gF_u(a\otimes w)\)，因而保持 \(G\) 作用。

从 \(u\) 出发再限制，得到 \(F_u(1\otimes w)=u(w)\)。从 \(F\) 出发再延伸，则
\(a\cdot F(1\otimes w)=F(a\otimes w)\)，恢复 \(F\)。因此两映射互逆。记正向同构为 \(\Phi\)。若 \(\alpha:W'\rightarrow W\) 为 \(k[H]\)-模同态、\(\beta:V\rightarrow V'\) 为 \(k[G]\)-模同态，则
\[
 \Phi\left(\beta\circ F\circ\operatorname{Ind}\alpha\right)
 =\beta\circ\Phi(F)\circ\alpha.
\]
确实，两边在 \(w'\) 处的取值都是 \(\beta F\bigl(1\otimes\alpha(w')\bigr)\)。这说明同构分别与 \(W\) 变量中的前复合及 \(V\) 变量中的后复合相容，即具有两变量自然性。
\end{proof}

上述结果称为\term[Frobenius-hufanlv]{Frobenius 互反律}[Frobenius reciprocity]。用伴随的语言，诱导是限制的左伴随；\eqref{b5:eq:frobenius-hom}及其两条具体公式给出了这一关系的含义。

\ProseParagraphBreak
其中有两种值得记住的映射。取 \(V=\operatorname{Ind}_H^G W\)，恒等映射对应于 \(k[H]\)-模同态 \(w\mapsto1\otimes w\)。取 \(W=\operatorname{Res}_H^G V\)，恒等映射则对应于 \(k[G]\)-模同态
\[
 k[G]\otimes_{k[H]}V\longrightarrow V,
 \qquad g\otimes v\longmapsto gv.
\]
后一映射总是满射，因为每个 \(v\) 都是 \(1\otimes v\) 的像。若 \(V\) 是非零有限维表示且 \(H\ne G\)，两边维数不同，它便不是同构。

\subsection{特征标内积形式}
\label{b5:subsec:0302:02}

现在把基域改为 \(\mathbb C\)，并取有限维表示。沿用\cref{b5:thm:character-hom-inner-product}的约定：类函数内积对第一变量共轭线性，即
\[
 \langle\alpha,\beta\rangle_G
 =\frac1{|G|}\sum_{g\in G}\overline{\alpha(g)}\,\beta(g),
 \qquad
 \langle\chi_U,\chi_V\rangle_G
 =\dim_{\mathbb C}\operatorname{Hom}_{\mathbb C[G]}(U,V).
\]
限制表示的特征标是函数的限制；把 \(\operatorname{Ind}_H^G W\) 的特征标简记为 \(\operatorname{Ind}_H^G\chi_W\)。下一节将给出只用 \(\chi_W\) 的函数值计算它的公式。

\begin{corollary}\label{b5:cor:frobenius-character}
设 \(G\) 为有限群，\(H\le G\)，\(W\) 为有限维复 \(H\)-表示，\(V\) 为有限维复 \(G\)-表示。则
\begin{equation}\label{b5:eq:frobenius-character}
 \langle\operatorname{Ind}_H^G\chi_W,\chi_V\rangle_G
 =\langle\chi_W,\operatorname{Res}_H^G\chi_V\rangle_H.
\end{equation}
同时有交换两个变量后的等式
\(\langle\chi_V,\operatorname{Ind}_H^G\chi_W\rangle_G
=\langle\operatorname{Res}_H^G\chi_V,\chi_W\rangle_H\)。
\end{corollary}
\begin{proof}
\cref{b5:thm:character-hom-inner-product}把第一条等式两边分别化为\eqref{b5:eq:frobenius-hom}两侧的 Hom 空间维数，而\cref{b5:thm:frobenius-reciprocity}给出其同构。第二条等式由第一条取复共轭得到；不能在未说明内积约定时直接调换参数。
\end{proof}

\subsection{重数计算}
\label{b5:subsec:0302:03}

若 \(U\) 为不可约复 \(G\)-表示，则由 Maschke 定理及 Schur 引理，它在 \(\operatorname{Ind}_H^G W\) 中的重数等于
\begin{equation}\label{b5:eq:induction-multiplicity}
 [\operatorname{Ind}_H^G W:U]
 =\langle\chi_W,\operatorname{Res}_H^G\chi_U\rangle_H.
\end{equation}
方括号在此表示直和分解中的重数。若 \(W\) 也不可约，右边就是 \(W\) 在 \(\operatorname{Res}_H^G U\) 中的重数。诱导与限制因此不是互逆操作，却共享同一张重数表的两个方向。

\begin{example}\label{b5:exm:s3-transposition-induction}
取 \(G=S_3\)、\(H=\bigl\langle(12)\bigr\rangle\)，在复数域上计算。设 \(\mathbf1\)、\(\varepsilon\) 分别为 \(G\) 的平凡表示与符号表示，\(U\) 为自然置换表示 \(\mathbb C^3\) 中坐标和为零的二维子空间。其特征标在单位元、换位、三轮换上的值依次为 \(2,0,-1\)：置换矩阵的迹等于固定坐标数，再减去常向量子空间上的迹 \(1\) 即得。于是
\[
 \langle\chi_U,\chi_U\rangle_G
 =\frac{4+3\cdot0+2\cdot1}{6}=1,
\]
故由\cref{b5:cor:character-irreducibility}，\(U\) 不可约。又
\(1^2+1^2+2^2=6\)，由\cref{b5:thm:regular-decomposition}，这三个表示穷尽 \(S_3\) 的不可约表示。

\ProseParagraphBreak
记 \(H\) 的非平凡一维表示为 \(\eta\)。限制到 \(H\) 后，\(\mathbf1\)、\(\varepsilon\)、\(U\) 分别成为 \(\mathbf1\)、\(\eta\)、\(\mathbf1\oplus\eta\)。对 \(U\)，这是因为换位的平方为 \(1\)，在复数域上可对角化，且迹为 \(0\)，两个特征值只能各为 \(1,-1\) 一次。\eqref{b5:eq:induction-multiplicity}于是给出
\[
 \operatorname{Ind}_H^G\mathbf1\cong\mathbf1\oplus U,
 \qquad
 \operatorname{Ind}_H^G\eta\cong\varepsilon\oplus U.
\]
两边维数均为 \(3\)，也与\eqref{b5:eq:induction-dimension}一致。
\end{example}

还有一个常用的特殊情形。若 \(W=\mathbf1_H\)，则
\[
 [\operatorname{Ind}_H^G\mathbf1_H:U]
 =\dim_{\mathbb C} U^H,
 \qquad
 U^H=\{\,u\in U:hu=u\text{ 对所有 }h\in H\,\}.
\]
因为从平凡表示到 \(\operatorname{Res}_H^G U\) 的同态完全由一个 \(H\)-不变向量决定。这一公式把置换表示中的重数与子群固定向量联系起来，并不要求先列出整张特征标表。

% !TeX root = ../../Book5.tex
% 纲要标题：### 3.3 诱导特征标
% 纲要位置：工作记录/计划纲要.md，第 3803 行。
% 纲要细目（写作范围）：
% * 陪集求迹与诱导特征标公式；任意域上的迹函数
% * 置换特征标
% * 从子群构造不可约表示


\section{诱导特征标}
\label{b5:sec:0303}

陪集模型不仅能写出矩阵，也能直接求迹。一个群元素把不同陪集上的分量互相交换时，这些分量对迹没有贡献；只有固定陪集才留下对角块。以下涉及特征标内积与不可约分解时均在复数域上工作。

\subsection{诱导特征标公式}
\label{b5:subsec:0303:01}

\begin{theorem}[诱导特征标公式]\label{b5:thm:induced-character}
设 \(G\) 为有限群，\(H\le G\)，\(W\) 为有限维复 \(H\)-表示，特征标为 \(\psi\)。对左陪集 \(G/H\) 的任一代表集 \(T\)，\(\operatorname{Ind}_H^G W\) 的特征标满足
\begin{align}
 (\operatorname{Ind}_H^G\psi)(g)
 &=\sum_{\substack{t\in T\\t^{-1}gt\in H}}\psi(t^{-1}gt)
 \label{b5:eq:induced-character-cosets}\\
 &=\frac1{|H|}\sum_{\substack{x\in G\\x^{-1}gx\in H}}\psi(x^{-1}gx),
 \qquad g\in G.
 \label{b5:eq:induced-character-all}
\end{align}
\end{theorem}
\begin{proof}
由\cref{b5:prop:induced-coset-model}，\(g\) 把第 \(t\) 个分量送到第 \(t'\) 个分量，其中 \(gt=t'h\)。若 \(t'\ne t\)，该块不在对角线上，对迹贡献为零。若 \(t'=t\)，则 \(h=t^{-1}gt\in H\)，其对角块就是 \(h\) 在 \(W\) 上的作用，迹为 \(\psi(t^{-1}gt)\)。逐块相加得到第一式。

将 \(t\) 换成同一左陪集中的 \(th\) 后，共轭元素变成 \(h^{-1}t^{-1}gth\)。条件是否属于 \(H\) 不变，而 \(\psi\) 是 \(H\) 上的类函数，所以函数值不变。对整个群求和时，每个左陪集恰有 \(|H|\) 个元素，除以 \(|H|\) 即得第二式。
\end{proof}

若 \(k\) 为任意域，\((W,\rho)\) 为有限维 \(k\) 上的 \(H\)-表示，记 \(\psi_k(h)=\operatorname{tr}_k\rho(h)\in k\)，则同一逐块求迹论证给出
\[
 \operatorname{tr}_k\bigl(g\mid\operatorname{Ind}_H^G W\bigr)
 =\sum_{\substack{t\in T\\t^{-1}gt\in H}}\psi_k(t^{-1}gt).
\]
此时求和取值于 \(k\)；若 \(|H|\) 在 \(k\) 中可逆，也可将它改写为\eqref{b5:eq:induced-character-all}的 \(k\)-值版本。若 \(\operatorname{char}k\bigm\vert |H|\)，则不能使用其中的除法。复特征标的标准形式可对照 Etingof 等人的 Theorem 5.9.1 及 Remark 5.9.2\cite[Section 5.9, pp. 106--107]{EtingofEtAl2011RepresentationTheory}；该书使用右陪集函数模型，把代表元取逆便得到这里的方向。

\ProseParagraphBreak
\eqref{b5:eq:induced-character-all}也适用于任意复类函数 \(\alpha:H\rightarrow\mathbb C\)，并定义一个线性映射 \(\operatorname{Ind}_H^G\) 到 \(G\) 的类函数空间。把 \(g\) 换成 \(ygy^{-1}\)，再把求和变量 \(x\) 换成 \(yx\)，可见所得函数对共轭不变。对任意 \(G\) 上的类函数 \(\beta\)，有限求和换序给出
\begin{align*}
 \langle\operatorname{Ind}_H^G\alpha,\beta\rangle_G
 &=\frac1{|G|\cdot|H|}\sum_{x\in G}\sum_{h\in H}
       \overline{\alpha(h)}\,\beta(xhx^{-1})\\
 &=\frac1{|H|}\sum_{h\in H}\overline{\alpha(h)}\,\beta(h)
 =\langle\alpha,\operatorname{Res}_H^G\beta\rangle_H.
\end{align*}
这证明了\eqref{b5:eq:frobenius-character}对所有复类函数仍成立，不必先把它们写成特征标的线性组合。

\ProseParagraphBreak
公式还给出两个快速检查：若 \(g\) 不共轭于 \(H\) 中任何元素，则诱导特征标在 \(g\) 处为零；在 \(g=1\) 处，它的值为 \([G:H]\psi(1)\)，恰是诱导空间的维数。

\subsection{置换特征标}
\label{b5:subsec:0303:02}

\begin{proposition}\label{b5:prop:permutation-character}
设有限群 \(G\) 作用在有限集 \(X\) 上，\(\mathbb C[X]\) 以 \(X\) 为基并由 \(G\) 置换这组基。其\term[zhihuantezhengbiao]{置换特征标}[permutation character] \(\pi_X\) 满足
\[
 \pi_X(g)=|X^g|,
 \qquad X^g=\{\,x\in X:gx=x\,\},\qquad g\in G.
\]
若 \(H\le G\) 且 \(X=G/H\) 带左乘作用，则 \(\mathbb C[X]\cong\operatorname{Ind}_H^G\mathbf1_H\)，其中 \(\mathbf1_H\) 为平凡复 \(H\)-表示。此外，记 \(\mathbf1_G\) 为平凡复 \(G\)-表示的特征标，则
\[
 \langle\mathbf1_G,\pi_X\rangle_G=|G\backslash X|,
 \qquad
 \langle\pi_X,\pi_X\rangle_G=|G\backslash(X\times X)|,
\]
右边分别为轨道数，\(X\times X\) 上使用对角作用。
\end{proposition}
\begin{proof}
置换矩阵的对角元在对应基向量固定时为 \(1\)，否则为 \(0\)，故迹为固定点数。对 \(X=G/H\)，映射 \(g\otimes1\mapsto e_{gH}\) 保持 \(gh\otimes1=g\otimes1\) 的关系，且把\cref{b5:prop:induced-coset-model}中的基双射到陪集基，因而是 \(G\)-同构。

\(\mathbb C[X]\) 的不变向量在同一轨道上有相同系数，所以不变空间以各轨道的基向量之和为基，维数等于轨道数。\cref{b5:thm:character-hom-inner-product}将其维数写成第一内积。又 \(\pi_X(g)\) 为实数，且 \((X\times X)^g=X^g\times X^g\)，故
\[
 \langle\pi_X,\pi_X\rangle_G
 =\frac1{|G|}\sum_{g\in G}|X^g|^2
 =\langle\mathbf1_G,\pi_{X\times X}\rangle_G.
\]
把已经证明的轨道数公式用于 \(X\times X\) 即得第二式。
\end{proof}

\begin{corollary}\label{b5:cor:two-transitive-augmentation}
设有限群 \(G\) 传递地作用在有限集 \(X\) 上，且 \(|X|\ge2\)。复置换表示中的子空间
\[
 U_X=\left\{\,\sum_{x\in X}a_xe_x:\sum_{x\in X}a_x=0\,\right\}
\]
不可约，当且仅当 \(G\) 在不同点的有序对集合上也传递，即原作用为\term[erchongchuandi]{二重传递}[doubly transitive]作用。
\end{corollary}
\begin{proof}
复数域中 \(|X|\ne0\)，所以常向量子空间与 \(U_X\) 直和为 \(\mathbb C[X]\)。传递性使平凡表示的重数为 \(1\)，故
\[
 \langle\pi_X-\mathbf1_G,\pi_X-\mathbf1_G\rangle_G
 =|G\backslash(X\times X)|-1.
\]
由\cref{b5:cor:character-irreducibility}，非零表示 \(U_X\) 不可约恰好要求右边为 \(1\)。对角线本身是一个轨道，因此总轨道数为 \(2\)，恰好等价于所有不同点的有序对构成另一个轨道。
\end{proof}

例如，当 \(n\ge2\) 时，\(S_n\) 在 \(n\) 个字母上的自然作用是二重传递的，因而坐标和为零的 \((n-1)\)-维复表示不可约。这给出一族无需逐个计算特征标表的不可约表示。

\subsection{从子群构造不可约表示}
\label{b5:subsec:0303:03}

诱导不保证保持不可约性；\cref{b5:exm:s3-transposition-induction}中的两个一维表示都诱导成可约表示。实际构造时可以先求诱导特征标，再利用范数为 \(1\) 的判据。若诱导后限制的分解容易看出，\eqref{b5:eq:frobenius-character}还会进一步缩短计算。

\begin{example}\label{b5:exm:dihedral-induced-family}
设 \(n\ge3\)，取阶为 \(2n\) 的二面体群
\[
 G=D_{2n}=\langle r,s\mid r^n=s^2=1,\ srs=r^{-1}\rangle,
 \qquad H=\langle r\rangle.
\]
固定 \(n\) 次本原单位根 \(\zeta\in\mathbb C\)。对 \(j\in\mathbb Z/n\mathbb Z\)，令 \(\lambda_j(r)=\zeta^j\)，并记 \(V_j=\operatorname{Ind}_H^G\lambda_j\)。陪集代表 \(1,s\) 给出
\[
 r\longmapsto\begin{pmatrix}\zeta^j&0\\0&\zeta^{-j}\end{pmatrix},
 \qquad
 s\longmapsto\begin{pmatrix}0&1\\1&0\end{pmatrix}.
\]
因此 \(\operatorname{Res}_H^G V_j\cong\lambda_j\oplus\lambda_{-j}\)。循环群的这两个一维特征标相同，当且仅当 \(2j=0\) 于 \(\mathbb Z/n\mathbb Z\)。由 Frobenius 互反律及一维特征标的正交性，
\[
 \langle\chi_{V_j},\chi_{V_j}\rangle_G
 =\langle\lambda_j,\lambda_j+\lambda_{-j}\rangle_H
 =\begin{cases}1,&2j\ne0,\\2,&2j=0.\end{cases}
\]
所以 \(2j\ne0\) 时 \(V_j\) 不可约。交换两个基向量给出 \(V_j\cong V_{-j}\)；反之，限制的不可约分解唯一性表明 \(V_i\cong V_j\) 只能在 \(i=j\) 或 \(i=-j\) 时成立。

\ProseParagraphBreak
若 \(2j=0\)，则 \(r\) 为标量矩阵，\(s\) 的两个特征向量给出两种一维表示：\(r\mapsto\zeta^j\)，\(s\mapsto1\) 或 \(-1\)。当 \(n\) 为奇数时，这样得到两种一维表示及 \((n-1)/2\) 种二维不可约表示；当 \(n\) 为偶数时，得到四种一维表示及 \((n-2)/2\) 种二维不可约表示。两种情形的维数平方和分别为
\[
 2+4\frac{n-1}{2}=2n,
 \qquad
 4+4\frac{n-2}{2}=2n.
\]
由\cref{b5:thm:regular-decomposition}，已经列尽全部不可约复表示。这个例子说明哪些一维表示诱导后不可约，哪些给出同构的诱导表示，以及哪些诱导表示会分解。
\end{example}

一般子群的诱导后限制未必像循环正规子群这样容易分解。下一章将用双陪集把它拆成若干较小的诱导表示，从而把本节的范数检验发展为系统的不可约性判据。

% !TeX root = ../../Book5.tex
% 纲要标题：### 3.4 函子关系
% 纲要位置：工作记录/计划纲要.md，第 3809 行。
% 纲要细目（写作范围）：
% * 诱导的传递性
% * 张量恒等式
% * 有限群情形诱导与余诱导的关系；无限指数时的陪集有限支集
% ---


\section{函子关系}
\label{b5:sec:0304}

诱导常常要和另一项构造接连使用：沿子群链诱导两次，把一个表示先张量再诱导，或研究映入诱导空间的同态。只比较维数不足以处理这些问题。本节分别写出相容于群作用的映射与逆映射；这些模同构均在任意域上成立。

\subsection{诱导的传递性}
\label{b5:subsec:0304:01}

\begin{theorem}[诱导的传递性]\label{b5:thm:induction-transitivity}
设 \(k\) 为域，\(K\le H\le G\) 为有限群的子群链，\(U\) 为左 \(k[K]\)-模。存在自然的左 \(k[G]\)-模同构
\[
 \operatorname{Ind}_H^G\left(\operatorname{Ind}_K^H U\right)
 \cong \operatorname{Ind}_K^G U,
 \qquad
 a\otimes(b\otimes u)\longmapsto ab\otimes u
 \quad(a\in k[G],\ b\in k[H],\ u\in U).
\]
\end{theorem}
\begin{proof}
映射对 \(a,b,u\) 多线性。内层的关系 \(bc\otimes u=b\otimes cu\)，\(c\in k[K]\)，在目标中成为 \(abc\otimes u=ab\otimes cu\)。外层的关系 \(ah\otimes(b\otimes u)=a\otimes(hb\otimes u)\)，\(h\in k[H]\)，在目标中两边都等于 \(ahb\otimes u\)。所以映射良定义，且保持两端的 \(G\) 左乘作用。

逆映射为 \(c\otimes u\mapsto c\otimes(1\otimes u)\)。对 \(d\in k[K]\)，有
\[
 cd\otimes(1\otimes u)
 =c\otimes(d\otimes u)
 =c\otimes(1\otimes du),
\]
故它也保持张量关系。一个方向的复合显然为恒等；另一个方向使用
\(ab\otimes(1\otimes u)=a\otimes(b\otimes u)\)，同样恢复原元。对 \(K\)-同态 \(U\rightarrow U'\)，两条映射仅把最后一个分量送到其像，因此同构自然。
\end{proof}

选择 \(G/H\) 的左陪集代表 \(T\) 及 \(H/K\) 的左陪集代表 \(S\)，则 \(\{ts:t\in T,s\in S\}\) 是 \(G/K\) 的左陪集代表集。上述同构把 \(t\otimes(s\otimes u)\) 送到 \(ts\otimes u\)，因而也明确说明了两层陪集坐标与一层坐标如何对应。陪集数满足 \([G:H][H:K]=[G:K]\)，而这一同构还保留了全部群作用。

\subsection{张量恒等式}
\label{b5:subsec:0304:02}

若 \(V_1,V_2\) 为同一群的群代数模，域上张量积使用对角作用 \(g(v_1\otimes v_2)=gv_1\otimes gv_2\)。下面需要把这个作用与诱导张量积中的平衡关系分开。

\begin{theorem}[张量恒等式]\label{b5:thm:tensor-identity}
设 \(k\) 为域，\(G\) 为有限群，\(H\le G\)，\(W\) 为左 \(k[H]\)-模，\(V\) 为左 \(k[G]\)-模。以下同构对两个模自然：
\begin{equation}\label{b5:eq:tensor-identity}
 \operatorname{Ind}_H^G\left(W\otimes_k\operatorname{Res}_H^G V\right)
 \cong (\operatorname{Ind}_H^G W)\otimes_k V.
\end{equation}
它在 \(g\in G\) 的纯张量上由
\[
 g\otimes(w\otimes v)\longmapsto(g\otimes w)\otimes gv
\]
给出；右边使用 \(G\) 的对角作用。
\end{theorem}
\begin{proof}
群元素是 \(k[G]\) 的基，所以给定公式先按 \(g\) 线性延伸，具体为
\[
\left(\sum_g c_gg\right)\otimes(w\otimes v)
\longmapsto\sum_g c_g(g\otimes w)\otimes gv.
\]
不能把一般 \(a\in k[G]\) 直接放进右边的两个位置写成 \((a\otimes w)\otimes av\)：那样会产生两个求和指标和系数的乘积，不再是关于 \(a\) 的线性延伸。这里逐个群元素同时作用于两个因子，再对这些基元素的像求和。对 \(h\in H\)，映射在 \(gh\otimes(w\otimes v)\) 处的值为
\[
 (gh\otimes w)\otimes ghv=(g\otimes hw)\otimes g(hv),
\]
这正是在 \(g\otimes(hw\otimes hv)\) 处的值，故保持 \(k[H]\) 上的平衡关系。对 \(x\in G\)，将 \(g\) 换成 \(xg\) 后的像为 \((xg\otimes w)\otimes xgv\)，等于原像受到 \(x\) 的对角作用，故它是 \(G\)-同态。

逆映射在群元素上为
\[
 (g\otimes w)\otimes v\longmapsto g\otimes(w\otimes g^{-1}v).
\]
其良定义需要检查 \(gh\otimes w=g\otimes hw\)。从前者出发，像为
\[
 gh\otimes(w\otimes h^{-1}g^{-1}v)
 =g\otimes(hw\otimes g^{-1}v),
\]
与从后者出发相同。两种复合分别消去 \(g^{-1}g\) 与 \(gg^{-1}\)，故为恒等。\(H\)-同态与 \(G\)-同态分别保持相应的群作用，所以逐个代入这两条公式也验证了两变量自然性。
\end{proof}

这里的名称为\term[zhanglianghengdengshi]{张量恒等式}[tensor identity]，也常按其用途称投影公式。复数域的特征标取值满足张量积对应逐点相乘，故\eqref{b5:eq:tensor-identity}给出
\[
 \operatorname{Ind}_H^G\left(\psi\cdot\operatorname{Res}_H^G\chi\right)
 =(\operatorname{Ind}_H^G\psi)\cdot\chi.
\]
这个等式对任意复类函数 \(\psi,\chi\) 也成立：在\eqref{b5:eq:induced-character-all}中，\(\chi(x^{-1}gx)=\chi(g)\)，可直接将该因子提出求和。

\ProseParagraphBreak
例如令 \(H=\{1\}\)、\(W=k\)，可得 \(k[G]\otimes_k V\) 上的对角作用同构于 \(\dim_k V\) 个正则表示的直和（\(V\) 有限维时）。这不是半单性结论；\cref{b5:thm:tensor-identity}中的显式映射说明，即使 \(\operatorname{char}k\) 整除 \(|G|\)，该同构仍然成立。

\subsection{有限群情形诱导与余诱导的关系}
\label{b5:subsec:0304:03}

Frobenius 互反律将 \(\operatorname{Hom}_{k[G]}(\operatorname{Ind}_H^G W,V)\) 识别为 \(\operatorname{Hom}_{k[H]}(W,\operatorname{Res}_H^G V)\)。若给定的却是 \(\operatorname{Res}_H^G V\rightarrow W\)，就要寻找一个 \(G\)-模，使映入它的同态对应于这些 \(H\)-同态。下面的余诱导满足
\[
 \operatorname{Hom}_{k[G]}(V,\operatorname{Coind}_H^G W)
 \cong\operatorname{Hom}_{k[H]}(\operatorname{Res}_H^G V,W).
\]

\begin{definition}\label{b5:def:coinduction}
设 \(k\) 为域，\(G\) 为有限群，\(H\le G\)，\(W\) 为左 \(k[H]\)-模。把 \(k[G]\) 以左乘看成左 \(k[H]\)-模，定义
\[
 \operatorname{Coind}_H^G W
 \coloneqq\operatorname{Hom}_{k[H]}\bigl(k[G],W\bigr),
 \qquad (gf)(a)=f(ag).
\]
这里 \(g\in G\)、\(f\in\operatorname{Hom}_{k[H]}(k[G],W)\)、\(a\in k[G]\)。所得左 \(k[G]\)-模称为 \(W\) 的余诱导模。若 \(W\) 有限维，则此模也是有限维的，称为 \(W\) 的\term[yuyoudaobiaoshi]{余诱导表示}[coinduced representation]。
\end{definition}
\symidx{\(\operatorname{Coind}_H^G\)：从 \(H\) 到 \(G\) 的余诱导}

该作用保持左 \(H\)-线性，因为 \(f(hag)=hf(ag)\)；它满足
\(\bigl(g_1(g_2f)\bigr)(a)=f(ag_1g_2)=\bigl((g_1g_2)f\bigr)(a)\)。等价地，余诱导空间由函数 \(f:G\rightarrow W\) 组成，满足
\begin{equation}\label{b5:eq:coinduced-functions}
 f(hx)=hf(x),\qquad (gf)(x)=f(xg).
\end{equation}
函数在群元素上的取值唯一线性延拓到群代数。这时使用的是右陪集 \(H\backslash G\)；它与前面的左陪集模型不可混用。

\begin{proposition}\label{b5:prop:coinduction-adjunction}
设 \(k\) 为域，\(G\) 为有限群，\(H\le G\)，\(V\) 为左 \(k[G]\)-模，\(W\) 为左 \(k[H]\)-模。存在自然同构
\[
 \operatorname{Hom}_{k[G]}(V,\operatorname{Coind}_H^G W)
 \cong\operatorname{Hom}_{k[H]}(\operatorname{Res}_H^G V,W).
\]
其正向映射为 \(F\mapsto\bigl(v\mapsto F(v)(1)\bigr)\)，逆映射把 \(u\) 送到 \(F_u(v)(x)=u(xv)\)。
\end{proposition}
\begin{proof}
若 \(F\) 为 \(G\)-同态，则
\(F(hv)(1)=\bigl(hF(v)\bigr)(1)=F(v)(h)=hF(v)(1)\)，故取值于 \(1\) 得到 \(H\)-同态。反之，\(F_u(v)(hx)=u(hxv)=hu(xv)\)，所以 \(F_u(v)\) 满足\eqref{b5:eq:coinduced-functions}；并且
\(F_u(gv)(x)=u(xgv)=\bigl(gF_u(v)\bigr)(x)\)，故 \(F_u\) 为 \(G\)-同态。

在 \(1\) 处求值恢复 \(u(v)\)；反向恢复 \(F\) 则使用
\(F(xv)(1)=\bigl(xF(v)\bigr)(1)=F(v)(x)\)。前复合 \(G\)-同态与后复合 \(H\)-同态均可直接移入取值公式，因此该对应对两个变量自然。
\end{proof}

\begin{theorem}\label{b5:thm:induction-coinduction}
设 \(k\) 为任意域，\(G\) 为有限群，\(H\le G\)。对每个左 \(k[H]\)-模 \(W\)，存在自然的 \(k[G]\)-模同构
\[
 \operatorname{Ind}_H^G W\cong\operatorname{Coind}_H^G W.
\]
这一同构不要求 \(|G|\) 或 \(|H|\) 在 \(k\) 中可逆。
\end{theorem}
\begin{proof}
定义 \(p_H:k[G]\rightarrow k[H]\)，保留群代数元素在 \(H\) 上的系数，删去其余项。对 \(h_1,h_2\in H\)，有
\(p_H(h_1ah_2)=h_1p_H(a)h_2\)，因为左右乘 \(H\) 中元素均保持 \(H\) 及其补集。一个诱导分量 \(t\otimes W\) 应对应一个右陪集上的函数：将左陪集 \(tH\) 取逆得到 \(Ht^{-1}\)，并要求 \(t\otimes w\) 对应的函数在 \(ht^{-1}\) 处取 \(hw\)。条件 \(x\in Ht^{-1}\) 等价于 \(xt\in H\)，而 \(x=ht^{-1}\) 时恰有 \(xt=h\)。因此可以用 \(p_H(xt)\) 同时判断支集并给出作用于 \(w\) 的元素。对一般 \(a\in k[G]\) 线性延伸，令
\[
 \Theta(a\otimes w)(x)=p_H(xa)w\qquad(x\in G).
\]
在 \(a=t\in G\) 时，这个公式具体就是
\[
\Theta(t\otimes w)(x)=
\begin{cases}hw,&x=ht^{-1},\ h\in H,\\0,&x\notin Ht^{-1}.\end{cases}
\]
\(\Theta(a\otimes w)(hx)=hp_H(xa)w\)，故其值确在余诱导空间中。又
\(p_H(xah)w=p_H(xa)hw\)，所以它保持张量关系。\(G\) 等变性则来自
\[
 \Theta(ga\otimes w)(x)=p_H(xga)w
 =\bigl(g\Theta(a\otimes w)\bigr)(x).
\]

选择左陪集代表 \(T\)。其逆元组成右陪集代表集 \(T^{-1}\)。对 \(s,t\in T\)，\(s^{-1}t\in H\) 当且仅当 \(s=t\)，所以
\[
 \Theta(t\otimes w)(s^{-1})=
 \begin{cases}w,&s=t,\\0,&s\ne t.\end{cases}
\]
一个满足\eqref{b5:eq:coinduced-functions}的函数由它在 \(T^{-1}\) 上的值唯一确定。因此 \(\Theta\) 的逆为
\[
 f\longmapsto\sum_{t\in T}t\otimes f(t^{-1}).
\]
这条和式有限；它在 \(T^{-1}\) 上逐个恢复函数值，因而恢复整个函数。\(\Theta\) 的定义本身不依赖 \(T\)，故逆映射也不依赖所选代表。对 \(H\)-同态 \(W\rightarrow W'\)，保留系数再作用与先送到 \(W'\) 的结果相同，证明自然性。整个构造没有任何除法。
\end{proof}

由\cref{b5:prop:coinduction-adjunction}及\cref{b5:thm:induction-coinduction}，有限群情形还得到
\begin{equation}\label{b5:eq:frobenius-right-hom}
 \operatorname{Hom}_{k[G]}(V,\operatorname{Ind}_H^G W)
 \cong\operatorname{Hom}_{k[H]}(\operatorname{Res}_H^G V,W).
\end{equation}
因此诱导既是限制的左伴随，也是其右伴随；两个方向的 Hom 公式各有自己的具体映射，并非把箭头任意倒转所得。

\ProseParagraphBreak
Etingof 等人的 Definition 5.8.1 直接使用\eqref{b5:eq:coinduced-functions}作为诱导模型，其 Theorem 5.10.1 因而首先给出右侧伴随公式\cite[Sections 5.8, 5.10, pp. 105--109]{EtingofEtAl2011RepresentationTheory}。\cref{b5:thm:induction-coinduction}说明了两种叙述的对应。对一般群及无限指数子群，余诱导仍允许全部满足 \(f(hx)=hf(x)\) 的函数，而张量诱导对应的函数只在有限多个右 \(H\)-陪集上非零。这里的有限支集是相对于陪集空间 \(H\backslash G\) 而言；若 \(H\) 无限，函数在一个非零陪集内也可有无限多个非零值。逆映射的和式因而只对具有这种陪集有限支集的函数有定义。


\begin{chaptersummary}
诱导空间 \(k[G]\otimes_{k[H]}W\) 把子群作用写进平衡关系，左陪集代表则把这个商空间分成若干 \(W\) 的副本。由此得到维数、显式矩阵、正合性与诱导特征标公式；这些结论的条件并不完全相同。含有 \(|H|^{-1}\) 的诱导特征标第二公式要求 \(|H|\) 在底域 \(k\) 中可逆，在 \(\operatorname{char}k\bigm\vert |H|\) 时不能直接使用。

\ProseParagraphBreak
Frobenius 互反律首先是 Hom 空间的自然同构。复表示半单性把它转为重数计算，但同构本身不依赖半单性。置换特征标把迹变成固定点数，把内积变成轨道数；二重传递作用和循环子群上的一维表示因而能产生具体的不可约表示。

\ProseParagraphBreak
传递性与张量恒等式说明诱导与其他构造的相容性。由等变函数构成的余诱导是限制的右伴随；有限指数时，函数模型与张量诱导模型自然同构，因此诱导既是限制的左伴随，也是右伴随。下一章进一步计算诱导后限制的结构，其分块由双陪集决定。
\end{chaptersummary}

\BookExercises

\begin{exercise}\label{b5:ex:03-cyclic-matrix}
设 \(G=\langle r\rangle\) 为六阶循环群，\(H=\langle r^2\rangle\)，\(\lambda^3=1\) 于 \(\mathbb C\)。令一维 \(H\)-表示满足 \(r^2w=\lambda w\)。求其诱导表示中 \(r\) 的矩阵，并列出不可约分解。
\end{exercise}
\begin{solution}
取左陪集代表 \(1,r\)，基为 \(1\otimes w,r\otimes w\)，则
\[
 r\longmapsto\begin{pmatrix}0&\lambda\\1&0\end{pmatrix}.
\]
取 \(\mu\in\mathbb C\) 满足 \(\mu^2=\lambda\)。两特征值为 \(\mu,-\mu\)，互不相同，且均满足六次方为 \(1\)。相应特征向量可取 \(\mu(1\otimes w)+r\otimes w\) 与 \(-\mu(1\otimes w)+r\otimes w\)。诱导表示因此是两种一维 \(G\)-表示 \(r\mapsto\mu\)、\(r\mapsto-\mu\) 的直和；它们恰好都是给定 \(H\)-表示的一维延拓。
\end{solution}

\begin{exercise}\label{b5:ex:03-modular-cyclic}
设 \(k=\mathbb F_2\)、\(G=C_4=\langle r\rangle\)、\(H=\langle r^2\rangle\)。证明 \(\operatorname{Ind}_H^G\mathbf1_H\) 的维数为 \(2\)，有一个平凡子表示和一个平凡商表示，却不是二者的直和。
\end{exercise}
\begin{solution}
在陪集基 \(e_1,e_2\) 上，\(r\) 交换两基向量。\(r-I\ne0\)，但 \((r-I)^2=0\)。其核为直线 \(k(e_1+e_2)\)，且该直线也等于其像，因此这条子表示及其商上的作用都平凡。若表示是两个平凡表示之和，\(r\) 应为恒等，与实际矩阵矛盾。任意可能的平凡补空间都应包含于 \(\ker(r-I)\)，也直接排除了分裂。
\end{solution}

\begin{exercise}\label{b5:ex:03-representative-change}
设 \(k\) 为域，\(G\) 为有限群，\(H\le G\)，\(\rho:H\to\operatorname{GL}_k(W)\) 为有限维表示。设 \(T=\{t_i\}\) 是 \(G/H\) 的左陪集代表集，\(T'=\{t_ih_i\}\)，其中 \(h_i\in H\)。给 \(W\) 固定一组基，证明由这两组代表所得诱导矩阵满足 \(M'(g)=D^{-1}M(g)D\) 对 \(g\in G\) 成立，其中 \(D\) 的第 \(i\) 个对角块为 \(\rho(h_i)\)。
\end{exercise}
\begin{solution}
新基向量 \(t_ih_i\otimes w_j\) 在旧基中的坐标为 \(t_i\otimes\rho(h_i)w_j\)。因此若用列向量记录坐标，则旧坐标等于 \(D\) 乘新坐标。先施加作用再改回新坐标，得到 \(M'(g)=D^{-1}M(g)D\)。特别地，陪集代表改变造成的是同时相似变换，而不是改变表示的同构类。
\end{solution}

\begin{exercise}\label{b5:ex:03-regular-restriction}
设 \(k\) 为域，\(H\le G\) 为有限群。证明正则表示满足
\[
 \operatorname{Res}_H^G k[G]\cong k[H]^{\oplus[G:H]},
 \qquad \operatorname{Ind}_H^G k[H]\cong k[G].
\]
说明第一式应使用哪一侧的陪集代表。
\end{exercise}
\begin{solution}
第一式是左 \(k[H]\)-模的分解，故选择右陪集 \(H\backslash G\) 的代表集 \(S\)，得到 \(k[G]=\bigoplus_{s\in S}k[H]s\)。每项映射 \(b\mapsto bs\) 为左 \(k[H]\)-模同构。第二式由 \(a\otimes b\mapsto ab\) 给出，平衡关系保证良定义；其逆为 \(a\mapsto a\otimes1\)，因为 \(ab\otimes1=a\otimes b\)。两映射都与 \(G\) 的左乘相容。
\end{solution}

\begin{exercise}\label{b5:ex:03-induced-invariants}
设 \(k\) 为域，\(H\le G\) 为有限群，\(W\) 为 \(k[H]\)-模。给出 \((\operatorname{Ind}_H^G W)^G\cong W^H\) 的显式同构，不假定群阶可逆。
\end{exercise}
\begin{solution}
取包含 \(1\) 的左陪集代表集 \(T\)。对 \(w\in W^H\)，向量 \(\sum_{t\in T}t\otimes w\) 在 \(G\) 作用下只发生分量置换，因为每个出现的 \(H\) 元素都固定 \(w\)，故该向量 \(G\)-不变。

反之，设 \(v=\sum_t t\otimes w_t\) 不变。\(h\in H\) 固定陪集 \(H\)，所以 \(1\) 分量表明 \(hw_1=w_1\)。对 \(t\in T\)，在 \(tv=v\) 中比较 \(t\) 分量：只有原来的 \(1\) 分量会被 \(t\) 送入该分量，故 \(w_t=w_1\)。因此逆映射是取 \(1\) 分量。所有公式均不含平均除法。
\end{solution}

\begin{exercise}\label{b5:ex:03-induced-coinvariants}
设 \(k\) 为域。对群 \(L\) 的左 \(k[L]\)-模 \(M\)，定义余不变商空间
\(M_L=M/\operatorname{span}_k\{\ell m-m:\ell\in L,m\in M\}\)。若 \(H\le G\) 为有限群，\(W\) 为左 \(k[H]\)-模，证明 \((\operatorname{Ind}_H^G W)_G\cong W_H\)，并写出同构。
\end{exercise}
\symidx{\(M_L\)：群作用的余不变商空间}
\begin{solution}
令 \(\epsilon:k[G]\rightarrow k\) 为系数和映射。公式 \(a\otimes w\mapsto\epsilon(a)[w]\) 对 \(H\) 平衡，因为 \([hw]=[w]\)；对 \(G\) 作用不变，因而降到余不变商。逆映射为 \([w]\mapsto[1\otimes w]\)，其良定义来自
\([1\otimes hw]=[h\otimes w]=[1\otimes w]\)。在商中 \([g\otimes w]=[1\otimes w]\)，故两个方向的复合均为恒等。
\end{solution}

\begin{exercise}\label{b5:ex:03-induction-projective-injective}
设 \(k\) 为域，\(H\le G\) 为有限群。证明诱导把投射 \(k[H]\)-模送到投射 \(k[G]\)-模，也把内射 \(k[H]\)-模送到内射 \(k[G]\)-模。可分别用满射的提升性质和沿单射的延伸性质定义投射性、内射性。
\end{exercise}
\begin{solution}
若 \(P\) 投射，给定 \(G\)-满射 \(X\rightarrow Y\) 及 \(\operatorname{Ind}P\rightarrow Y\)，由\cref{b5:thm:frobenius-reciprocity}，后一个映射对应于 \(P\rightarrow\operatorname{Res}Y\)。限制保持满射，故投射性给出到 \(\operatorname{Res}X\) 的提升。互反同构的自然性把它送回到 \(X\) 的提升，所以 \(\operatorname{Ind}P\) 投射。

若 \(I\) 内射，给定 \(G\)-单射 \(X\rightarrow Y\) 及 \(X\rightarrow\operatorname{Ind}I\)，由\eqref{b5:eq:frobenius-right-hom}，后一个映射对应于 \(\operatorname{Res}X\rightarrow I\)。限制保持单射，故可延伸到 \(\operatorname{Res}Y\rightarrow I\)，再由自然性转回 \(Y\rightarrow\operatorname{Ind}I\)。这里使用有限群的诱导—余诱导同构，不能只靠左伴随公式证明第二项。
\end{solution}

\begin{exercise}\label{b5:ex:03-a4-induction}
设 \(G=A_4\)，\(H=\bigl\langle(123)\bigr\rangle\)，\(N=\bigl\{1,(12)(34),(13)(24),(14)(23)\bigr\}\)。令 \(\omega\) 为三次本原单位根，\(\lambda_j((123))=\omega^j\)，\(j=0,1,2\)。求三个 \(\operatorname{Ind}_H^G\lambda_j\) 的不可约复分解。
\end{exercise}
\begin{solution}
\(N\) 正规且 \(G/N\) 为三阶循环群；\(H\cap N=\{1\}\)，故 \(H\rightarrow G/N\) 为同构。于是 \(\lambda_j\) 各自延拓为 \(G\) 的一维表示 \(\theta_j\)，该延拓在 \(N\) 上平凡。

取四个字母的置换表示中坐标和为零的三维子表示 \(U\)。其特征标在单位元、三个非平凡 \(N\) 元素和八个三轮换上分别为 \(3,-1,0\)。因此特征标范数平方为 \((9+3)/12=1\)，\(U\) 不可约。三种一维表示与 \(U\) 的维数平方和为 \(3+9=12\)，故穷尽全部不可约表示。限制到 \(H\) 时，\(\chi_U\) 的值为 \(3,0,0\)，从而 \(\operatorname{Res}_H^G U\cong\lambda_0\oplus\lambda_1\oplus\lambda_2\)。Frobenius 互反律给出
\[
 \operatorname{Ind}_H^G\lambda_j\cong\theta_j\oplus U
 \qquad(j=0,1,2).
\]
每个右边维数为 \(4=[G:H]\)。
\end{solution}

\begin{exercise}\label{b5:ex:03-induced-determinant}
设 \(k\) 为域，\(G\) 为有限群，\(H\le G\)，\(W\) 为 \(d\) 维左 \(k[H]\)-模，\(\rho\) 为它在所选基下的矩阵表示。取左陪集代表 \(t_1,\ldots,t_m\)，对 \(g\in G\) 唯一写成 \(gt_i=t_{\sigma_g(i)}h_i(g)\)，其中 \(\sigma_g\in S_m\)、\(h_i(g)\in H\)。证明
\[
 \det\left(g\mid\operatorname{Ind}_H^G W\right)
 =\operatorname{sgn}(\sigma_g)^d\prod_i\det\rho\bigl(h_i(g)\bigr).
\]
符号 \(\pm1\) 在域 \(k\) 中理解。
\end{exercise}
\begin{solution}
诱导矩阵是先在第 \(i\) 个块上施加 \(\rho\bigl(h_i(g)\bigr)\)，再按 \(\sigma_g\) 置换维数为 \(d\) 的块。因此对角块部分的行列式是所列乘积。交换两个长为 \(d\) 的块产生 \(d^2\) 次基向量交换，符号为 \((-1)^{d^2}=(-1)^d\)；把 \(\sigma_g\) 写成换位之积，得到块置换符号 \(\operatorname{sgn}(\sigma_g)^d\)。在特征 \(2\) 时所有符号都等于 \(1\)，公式仍成立。
\end{solution}

\begin{exercise}\label{b5:ex:03-quaternion-induction}
设 \(G=Q_8=\{\pm1,\pm i,\pm j,\pm k\}\)，\(H=\langle i\rangle\)。令一维复表示 \(\lambda\) 满足 \(\lambda(i)=\sqrt{-1}\)。计算 \(\operatorname{Ind}_H^G\lambda\) 的特征标，证明它不可约。
\end{exercise}
\begin{solution}
\(H\) 为指数二正规子群，取代表 \(1,j\)。共轭 \(j^{-1}ij=i^{-1}\)，所以限制到 \(H\) 的特征标为 \(\lambda+\lambda^{-1}\)，在 \(1,-1,\pm i\) 上分别为 \(2,-2,0\)。\(H\) 外元素不共轭于 \(H\) 中元素，故诱导特征标在 \(\pm j,\pm k\) 上为零。范数平方为 \((4+4)/8=1\)，所以它不可约。特别地，中心元 \(-1\) 虽固定所有陪集，却作用为 \(-I\)，说明“固定陪集数”只适用于诱导平凡表示的特殊情形。
\end{solution}

\begin{exercise}\label{b5:ex:03-two-subsets-s4}
令 \(S_4\) 作用在四个字母的六个二元素子集上。求置换特征标，并证明该置换表示分解为平凡表示、三维标准表示和一个二维不可约表示的直和。
\end{exercise}
\begin{solution}
按循环型 \(1^4,2\,1^2,2^2,3\,1,4\) 排列，固定二元素子集数依次为 \(6,2,2,0,0\)，共轭类大小分别为 \(1,6,3,8,6\)，所以置换特征标 \(\pi\) 的范数平方为 \((36+24+12)/24=3\)。作用传递，平凡表示重数为 \(1\)。

三维标准表示 \(U\) 不可约，由\cref{b5:cor:two-transitive-augmentation}可知。其特征标为 \(3,1,-1,0,-1\)，故 \(\langle\chi_U,\pi\rangle=(18+12-6)/24=1\)。Maschke 定理可取这两项的补表示 \(W\)，其特征标为 \(2,0,2,-1,0\)，范数平方为 \((4+12+8)/24=1\)。所以 \(W\) 不可约，得到所求分解。
\end{solution}

\begin{exercise}\label{b5:ex:03-regular-augmentation}
令 \(C_n\) 在自身上正则作用，\(n\ge2\)。计算它在有序对上的轨道数，并判断复置换表示的坐标和为零子空间何时不可约。
\end{exercise}
\begin{solution}
把群写成加法群 \(\mathbb Z/n\mathbb Z\)。同时平移不改变 \((a,b)\) 的差 \(b-a\)，而差相同的两个有序对能经平移互相得到，故有 \(n\) 个轨道。由\cref{b5:prop:permutation-character}，置换特征标范数平方为 \(n\)，去掉唯一平凡分量后范数平方为 \(n-1\)。因此该 \((n-1)\)-维子空间不可约当且仅当 \(n=2\)。这说明传递作用本身不足以保证增广子空间不可约。
\end{solution}

\begin{exercise}\label{b5:ex:03-dihedral-tensor}
设 \(n\ge3\) 为整数，\(D_{2n}=\langle r,s\mid r^n=s^2=1,\ srs=r^{-1}\rangle\)，\(H=\langle r\rangle\)，\(\zeta\in\mathbb C\) 为本原 \(n\) 次单位根。对 \(j\in\mathbb Z/n\mathbb Z\)，令 \(\lambda_j(r)=\zeta^j\)，\(V_j=\operatorname{Ind}_H^{D_{2n}}\lambda_j\)；这里允许 \(2j=0\)，因而 \(V_j\) 未必不可约。对 \(i,j\in\mathbb Z/n\mathbb Z\)，证明
\[
 V_i\otimes_{\mathbb C}V_j\cong V_{i+j}\oplus V_{i-j}.
\]
进一步在 \(D_{10}\) 中分解 \(V_1\otimes V_1\) 与 \(V_1\otimes V_2\)。
\end{exercise}
\begin{solution}
在旋转 \(r^a\) 上，\(\chi_{V_j}(r^a)=\zeta^{ja}+\zeta^{-ja}\)。两这样的值相乘恰好为 \(\chi_{V_{i+j}}(r^a)+\chi_{V_{i-j}}(r^a)\)。在反射上四个 \(V\) 的特征标均为零，等式仍成立。复表示由特征标决定，故得同构。对 \(n=5\)，有 \(V_{-j}\cong V_j\)，\(V_0=\mathbf1\oplus\varepsilon\)，其中 \(\varepsilon\) 在旋转上为 \(1\)、在反射上为 \(-1\)。于是
\[
 V_1\otimes V_1\cong V_2\oplus\mathbf1\oplus\varepsilon,
 \qquad
 V_1\otimes V_2\cong V_2\oplus V_1.
\]
\end{solution}

\begin{exercise}\label{b5:ex:03-tower-d8}
在 \(D_8=\langle r,s\mid r^4=s^2=1,srs=r^{-1}\rangle\) 中，设 \(K=\langle r^2\rangle\le H=\langle r\rangle\)。令 \(K\) 的复一维表示 \(\eta\) 满足 \(\eta(r^2)=-1\)。利用诱导传递性求 \(\operatorname{Ind}_K^{D_8}\eta\) 的分解。
\end{exercise}
\begin{solution}
\(\operatorname{Ind}_K^H\eta\) 中 \(r\) 的矩阵为 \(\left(\begin{smallmatrix}0&-1\\1&0\end{smallmatrix}\right)\)，其两个特征值为 \(\sqrt{-1}\) 与 \(-\sqrt{-1}\)，故分解为 \(H\) 的一维表示 \(\lambda_1\oplus\lambda_3\)。由\cref{b5:thm:induction-transitivity}及诱导保持直和，所求表示为 \(V_1\oplus V_3\)。\cref{b5:exm:dihedral-induced-family}说明 \(V_1\cong V_3\)，且它们为二维不可约表示，故答案为 \(V_1^{\oplus2}\)。总维数为 \(4=[D_8:K]\)。
\end{solution}

\begin{exercise}\label{b5:ex:03-counit-splitting}
设 \(k\) 为域，\(H\le G\) 为有限群，\(V\) 为 \(k[G]\)-模。证明若 \([G:H]\) 在 \(k\) 中可逆，则映射 \(\epsilon_V:\operatorname{Ind}_H^G\operatorname{Res}_H^G V\rightarrow V\)、\(g\otimes v\mapsto gv\) 有 \(G\)-线性右逆。再给出指数不可逆时它不分裂的例子。
\end{exercise}
\begin{solution}
取左陪集代表集 \(T\)，定义 \(j(v)=\sum_{t\in T}t\otimes t^{-1}v\)。把 \(t\) 换成 \(th\) 时相应项不变；对 \(g\in G\)，用 \(gt=t'h\) 重编号可得
\(g(t\otimes t^{-1}v)=t'\otimes t'^{-1}gv\)，故 \(j\) 为 \(G\)-同态。又 \(\epsilon_Vj(v)=[G:H]v\)，所以 \([G:H]^{-1}j\) 为所求右逆。

取 \(G=C_p\)、\(H=\{1\}\)、\(k=\mathbb F_p\)、\(V=k\) 平凡。此时 \(\epsilon_V\) 是正则模的系数和映射。正则模中的不变向量都是 \(\sum_{g\in G}g\) 的标量倍数，其系数和为 \(p=0\)。若存在右逆，\(1\in k\) 必须映到系数和为 \(1\) 的不变向量，矛盾。
\end{solution}

\begin{exercise}\label{b5:ex:03-induced-dual}
设 \(k\) 为域，\(H\le G\) 为有限群，\(W\) 为有限维 \(k[H]\)-模；对偶作用为 \((h\varphi)(w)=\varphi(h^{-1}w)\)。构造自然同构
\(\operatorname{Ind}_H^G(W^*)\cong(\operatorname{Ind}_H^G W)^*\)。
\end{exercise}
\begin{solution}
在两个诱导空间之间定义配对，对 \(g,t\in G\) 令
\[
 B(g\otimes\varphi,t\otimes w)=
 \begin{cases}\varphi(g^{-1}tw),&g^{-1}t\in H,\\0,&g^{-1}t\notin H.\end{cases}
\]
第一变量的关系 \(gh\otimes\varphi=g\otimes h\varphi\) 对应于两边共同的值 \(\varphi(h^{-1}g^{-1}tw)\)，第二变量的平衡关系由 \(t h\otimes w=t\otimes hw\) 直接保持。左乘同一个 \(x\in G\) 时，\((xg)^{-1}(xt)=g^{-1}t\)，所以 \(B(xu,xv)=B(u,v)\)。

选择同一组左陪集代表 \(T\)，由\cref{b5:prop:induced-coset-model}把两个诱导空间分别写成 \(\bigoplus_{t\in T}W^*\) 与 \(\bigoplus_{t\in T}W\)。对 \(s,t\in T\)，\(s\ne t\) 时配对为零，而
\[
 B(t\otimes\varphi,t\otimes w)=\varphi(w).
\]
因此每一对同代表分量之间都是非退化的求值配对。\(T\) 与 \(\dim_k W\) 均有限，故 \(B\) 给出线性同构
\[
 \Phi_W:\operatorname{Ind}_H^G(W^*)\longrightarrow
 (\operatorname{Ind}_H^G W)^*,\qquad
 \Phi_W(u)(v)=B(u,v).
\]
对偶作用满足 \((x\lambda)(v)=\lambda(x^{-1}v)\)，于是
\[
 \Phi_W(xu)(v)=B(xu,v)=B(u,x^{-1}v)
 =(x\Phi_W(u))(v).
\]
这说明 \(\Phi_W\) 为 \(G\)-同态。

若 \(W'\) 也是有限维 \(k[H]\)-模，\(f:W\rightarrow W'\) 为 \(H\)-同态，则对 \(u'\in\operatorname{Ind}_H^G(W'^*)\)、\(v\in\operatorname{Ind}_H^G W\)，配对公式与 \(f\) 的 \(H\)-等变性给出
\[
 B_W\bigl(\operatorname{Ind}(f^*)(u'),v\bigr)
 =B_{W'}\bigl(u',\operatorname{Ind}(f)(v)\bigr).
\]
因而
\[
 \Phi_W\circ\operatorname{Ind}(f^*)
 =(\operatorname{Ind}(f))^*\circ\Phi_{W'}.
\]
这正是两种反变构造之间的自然性。
\end{solution}

\begin{exercise}\label{b5:ex:03-function-basis}
对有限群 \(G\) 及平凡子群 \(H=\{1\}\)，把 \(\operatorname{Coind}_1^G k\) 看成所有函数 \(G\rightarrow k\)。记 \(\delta_s\) 为在 \(s\) 处取 \(1\)、其余处为零的函数。求 \(g\delta_s\)，并说明正则模的群元素基 \(t\) 在\cref{b5:thm:induction-coinduction}下对应哪个函数。
\end{exercise}
\begin{solution}
\((g\delta_s)(x)=\delta_s(xg)\) 非零当且仅当 \(x=sg^{-1}\)，故 \(g\delta_s=\delta_{sg^{-1}}\)。自然同构把 \(t\) 送到 \(x\mapsto p_1(xt)\)，这正是 \(\delta_{t^{-1}}\)。于是 \(g\delta_{t^{-1}}=\delta_{t^{-1}g^{-1}}=\delta_{(gt)^{-1}}\)，与左正则作用吻合。若误把 \(t\) 对应到 \(\delta_t\)，一般不能得到这种等变性。
\end{solution}

\begin{exercise}\label{b5:ex:03-right-frobenius-explicit}
设 \(k\) 为域，\(G\) 为有限群，\(H\le G\)，\(V\) 为左 \(k[G]\)-模，\(W\) 为左 \(k[H]\)-模，\(u:\operatorname{Res}_H^G V\rightarrow W\) 为 \(H\)-同态，\(T\) 为 \(G/H\) 的左陪集代表集。证明右侧互反同构\eqref{b5:eq:frobenius-right-hom}的逆把 \(u\) 送到
\[
 v\longmapsto\sum_{t\in T}t\otimes u(t^{-1}v),
\]
并写出从 \(G\)-同态到 \(u\) 的反向公式。
\end{exercise}
\begin{solution}
由\cref{b5:prop:coinduction-adjunction}，\(u\) 首先给出函数 \(x\mapsto u(xv)\)。\cref{b5:thm:induction-coinduction}的逆把该函数送到 \(\sum_t t\otimes u(t^{-1}v)\)，正是所求公式。若替换 \(t\) 为 \(th\)，其项变成 \(th\otimes u(h^{-1}t^{-1}v)=t\otimes u(t^{-1}v)\)，所以不依赖代表。

反向对 \(G\)-同态 \(F:V\rightarrow\operatorname{Ind}_H^G W\) 先施加\cref{b5:thm:induction-coinduction}中的 \(\Theta\)，再在 \(1\) 处取值。设 \(t_0\in T\) 是陪集 \(H\) 的代表，故 \(t_0\in H\)。把 \(F(v)\) 唯一写成 \(\sum_{t\in T}t\otimes w_t(v)\)，则
\[
 u(v)=\Theta(F(v))(1)=t_0w_{t_0}(v).
\]
这里除 \(t_0\) 外的代表均不在 \(H\) 中，因而它们的项在 \(1\) 处取值为零。若 \(1\in T\)，公式才简化为 \(u(v)=w_1(v)\)。
\end{solution}

\begin{exercise}\label{b5:ex:03-relative-trace}
设 \(k\) 为域，\(G\) 为有限群，\(H\le G\)，\(V\) 为左 \(k[G]\)-模，\(a\in\operatorname{End}_{k[H]}(V)\)，\(T\) 为 \(G/H\) 的左陪集代表集。证明
\[
 \operatorname{Tr}_H^G(a)\coloneqq\sum_{t\in T}t a t^{-1}
\]
不依赖左陪集代表集 \(T\)，并属于 \(\operatorname{End}_{k[G]}(V)\)。此处群元素表示其在 \(V\) 上的算子。若 \(a\) 已为 \(G\)-同态，求此和式。
\end{exercise}
\begin{solution}
因为 \(a\) 与每个 \(h\in H\) 交换，\((th)a(th)^{-1}=tat^{-1}\)，故代表元选择无关。对 \(g\in G\)，\(g\operatorname{Tr}_H^G(a)g^{-1}\) 是用左陪集代表集 \(gT\) 计算的同一和式，所以等于原和式。这恰好表示它与全部 \(G\) 作用交换。若 \(a\) 本来就与 \(G\) 交换，则每一项等于 \(a\)，和式为 \([G:H]a\)。这个相对迹是算子的和；不要把它误认为标量值的矩阵迹。
\end{solution}

\begin{exercise}\label{b5:ex:03-induction-reflects-splitting}
设 \(0\rightarrow W'\rightarrow W\rightarrow W''\rightarrow0\) 为有限群 \(H\) 在域 \(k\) 上的表示短正合列，\(H\le G\) 且 \(G\) 有限。证明该列分裂，当且仅当诱导到 \(G\) 后的短正合列分裂。
\end{exercise}
\begin{solution}
若原列有 \(H\)-线性截面，诱导它便得到 \(G\)-线性截面。反之，对任何 \(H\)-模 \(M\)，在\cref{b5:prop:induced-coset-model}中选择含 \(1\) 的左陪集代表集，则 \(\operatorname{Ind}_H^G M\) 是各陪集对应的一份 \(M\) 的直和。映射 \(i_M:m\mapsto1\otimes m\) 正是陪集 \(H\) 对应分量的包含，因而单射，且有 \(H\)-线性投影 \(p_M\) 返回此分量：在陪集 \(H\) 上取原向量，在其他陪集上取零。\(H\) 的左乘同时保持 \(H\) 和 \(G\setminus H\)，因此投影确为 \(H\)-线性，并对 \(M\) 中的同态自然。

若诱导列有截面 \(s:\operatorname{Ind}W''\rightarrow\operatorname{Ind}W\)，则 \(p_W s i_{W''}:W''\rightarrow W\) 为 \(H\)-同态。记原满射为 \(q\)，利用自然性及诱导截面条件可得
\(q p_W s i_{W''}=p_{W''}(\operatorname{Ind}q)s i_{W''}=p_{W''}i_{W''}=1\)。故原列分裂。
\end{solution}

\begin{exercise}\label{b5:ex:03-real-multiplicity}
设 \(G=C_3\)、\(H=\{1\}\)，在实数域上取二维旋转表示 \(U\)，生成元作用为平面中转角 \(2\pi/3\) 的旋转。证明 \(U\) 不可约，\(\mathbb R[G]\cong\mathbb R\oplus U\)，但
\(\dim_{\mathbb R}\operatorname{Hom}_{\mathbb R[G]}\bigl(U,\mathbb R[G]\bigr)=2\)。解释为什么这不违背 Frobenius 互反律。
\end{exercise}
\begin{solution}
旋转矩阵没有实特征值，因而没有一维实不变子空间，\(U\) 不可约。又
\[
 \mathbb R[G]\cong\mathbb R[t]/(t^3-1)
 \cong\mathbb R\times\mathbb R[t]/(t^2+t+1)
 \cong\mathbb R\times\mathbb C.
\]
第二因子作为实表示同构于 \(U\)，故重数为 \(1\)。与该旋转交换的实线性算子，识别平面为 \(\mathbb C\) 后恰为复数乘法，所以 \(\operatorname{End}_{\mathbb R[G]}(U)\cong\mathbb C\)，实维数为 \(2\)。从 \(U\) 到平凡表示的同态为零，于是所述 Hom 空间维数为 \(2\)。

右侧互反同构给出
\(\operatorname{Hom}_{\mathbb R[G]}(U,\operatorname{Ind}_1^G\mathbb R)\cong\operatorname{Hom}_{\mathbb R}(U,\mathbb R)\)，右边也有实维数 \(2\)。互反同构完全成立；失效的是把实 Hom 维数不加修正地当作不可约重数，而复数域中 Schur 引理给出的标量自同态条件在这里不成立。
\end{solution}

\begin{exercise}\label{b5:ex:03-infinite-index}
把群代数张量诱导与 Hom 余诱导的定义用于 \(G=\mathbb Z\)、\(H=\{0\}\)、\(k=\mathbb F_2\)。证明两空间不可能同构为 \(k\)-向量空间，并指出\cref{b5:thm:induction-coinduction}证明中哪一步不再适用。
\end{exercise}
\begin{solution}
诱导空间是 \(\mathbb F_2[\mathbb Z]\)，即整数集上的有限支集函数。对每个有限支集只有有限种系数选择，而整数的有限子集全体可数，所以它的底集合可数。余诱导是全部函数 \(\mathbb Z\rightarrow\mathbb F_2\)，其底集合不可数：枚举整数后用对角法，就能构造一个不属于任何给定可数函数表的函数。故两者连底集合都不能双射，更不可能有线性同构。

原来的 \(\Theta\) 仍把诱导元素送到相应的有限支集函数，但不能覆盖任意函数。逆公式中的 \(\sum_{t\in T}t\otimes f(t^{-1})\) 对无限支集函数成为无限和，而代数张量积没有定义这种无限求和。
\end{solution}

% !TeX root = ../../Book5.tex
% 纲要标题：## 第4章　双陪集、Mackey 理论与 Clifford 理论
% 纲要位置：工作记录/计划纲要.md，第 3817 行。


\chapter{双陪集、Mackey 理论与 Clifford 理论}
\label{b5:ch:04}
\BookChapterNumber{b5:ch:04}

\begin{chapterabstract}
本章研究诱导与限制怎样组合，以及不可约表示在正规子群上如何分解。双陪集把陪集作用分成轨道，Mackey 公式把每个轨道写成一个新的诱导表示，由此计算交织空间并判定诱导表示的不可约性。Clifford 理论进一步说明，正规子群上的不可约成分组成一个共轭轨道，整个表示可从惯性群上的等型分量恢复。最后用因子集描述表示延拓的障碍，并在重数空间上建立射影表示。
\end{chapterabstract}

\begin{learninggoals}
\begin{itemize}
\item 能由陪集作用求双陪集、交集子群及每个分块的维数。
\item 写出 Mackey 分解的显式映射，说明两种伴随在交织空间公式中的作用。
\item 能计算正规子群上不可约成分的轨道、惯性群与重数，并应用 Clifford 对应。
\item 能构造标量因子集，检验余圈恒等式，判断延拓障碍及其对重数空间的影响。
\end{itemize}
预备知识为有限群的陪集与正规子群、复表示的完全可约性和 Schur 引理，以及\chapref{b5:ch:03}中的诱导、限制、余诱导和 Frobenius 互反性。前两节的模分解允许任意基域；不可约性判据和后两节的结构定理在复数域上讨论。
\end{learninggoals}

诱导由子群的表示构造整个群的表示，限制则保留子群元素的作用。计算诱导后限制的表示时，先求陪集空间中的轨道，再分析每个轨道对应的模。对正规子群，共轭还作用于不可约表示的同构类。陪集轨道与不可约类型的共轭轨道分别给出 Mackey 分解和 Clifford 理论。

% !TeX root = ../../Book5.tex
% 纲要标题：### 4.1 双陪集
% 纲要位置：工作记录/计划纲要.md，第 3819 行。
% 纲要细目（写作范围）：
% * 双陪集分解
% * 诱导后限制的几何分块
% * 小群的实际计算


\section{双陪集}
\label{b5:sec:0401}

把一个 \(H\)-表示诱导到 \(G\)，再只观察子群 \(K\) 的作用时，原来由 \(G/H\) 标记的各个向量空间并不会任意混合。哪些空间能够互相到达，取决于 \(K\) 在陪集集上的轨道。双陪集正是这些轨道的另一种写法。

\subsection{双陪集分解}
\label{b5:subsec:0401:01}

\begin{definition}\label{b5:def:double-coset}
设 \(G\) 为有限群，\(H,K\leq G\)，\(x\in G\)。集合
\[
KxH=\{kxh:k\in K,\ h\in H\}
\]
称为由 \(x\) 确定的 \term[shuangpeiji]{双陪集}[double coset]。所有这些集合组成的集合记为 \(K\backslash G/H\)。这个记号表示一个商集，一般没有商群结构。
\end{definition}
\symidx{\(K\backslash G/H\)：双陪集集}

关系 \(y=kxh\) 是等价关系：单位元给出自反性，等式 \(x=k^{-1}yh^{-1}\) 给出对称性，两次这样的表达式相接给出传递性。因此，不同双陪集互不相交，它们的并为 \(G\)。

\begin{proposition}\label{b5:prop:double-coset-orbits}
设 \(G\) 为有限群，\(H,K\leq G\)，\(x\in G\)。在左陪集集 \(G/H\) 上令 \(K\) 左乘作用。包含 \(xH\) 的轨道为 \(KxH/H\)，其稳定子群为
\[
L_x=K\cap xHx^{-1}.
\]
映射 \(kL_x\mapsto kxH\) 给出 \(K\)-集合的同构
\[
K/L_x\longrightarrow KxH/H.
\]
特别地，
\begin{equation}\label{b5:eq:double-coset-size}
|KxH|=\frac{|K|\cdot|H|}{|K\cap xHx^{-1}|}.
\end{equation}
\end{proposition}
\begin{proof}
元素 \(k\in K\) 固定 \(xH\)，当且仅当 \(x^{-1}kx\in H\)，即 \(k\in L_x\)。而 \(k_1xH=k_2xH\) 当且仅当 \(k_2^{-1}k_1\in L_x\)，所以所给映射良定义、单射，并由轨道的定义而满射。它与左乘作用相容。轨道含有 \([K:L_x]\) 个左 \(H\)-陪集，每个陪集含有 \(|H|\) 个元素，得到公式。
\end{proof}

于是，选取双陪集代表元，就是在 \(K\) 作用于 \(G/H\) 的每个轨道中选一个点。即使 \(H=K\)，交集也仍须写成 \(H\cap xHx^{-1}\)，不能把共轭后的子群忽略。

\subsection{诱导后限制的几何分块}
\label{b5:subsec:0401:02}

设 \(k\) 为任意域，\(W\) 为有限维左 \(k[H]\)-模。诱导表示采用
\[
\operatorname{Ind}_H^G W=k[G]\otimes_{k[H]}W.
\]
选取左陪集代表元 \(t\) 后，它作为向量空间等于 \(\bigoplus_{tH\in G/H}t\otimes W\)。若 \(a\in K\)，则 \(a\) 把 \(t\otimes W\) 送到 \(atH\) 所对应的空间。因此，每一个 \(K\)-轨道内的全部空间之和都是 \(K\)-不变子空间。

\begin{proposition}\label{b5:prop:double-coset-blocks}
设 \(G\) 为有限群，\(H,K\leq G\)，\(k\) 为域，\(W\) 为有限维左 \(k[H]\)-模。对每个双陪集 \(D=KxH\)，令 \(k[D]\) 为以 \(D\) 为基的子空间。它具有左 \(k[K]\)、右 \(k[H]\) 作用，且有左 \(k[K]\)-模直和分解
\[
\operatorname{Res}_K^G\operatorname{Ind}_H^G W
=\bigoplus_{D\in K\backslash G/H}k[D]\otimes_{k[H]}W.
\]
其中 \(D=KxH\) 所对应的分量的维数为
\[
[K:K\cap xHx^{-1}]\dim_k W.
\]
\end{proposition}
\begin{proof}
群元素基按互不相交的双陪集分组，得到
\(k[G]=\bigoplus_D k[D]\)。左右作用都保持各个 \(D\)，故这是双模的直和。对右 \(k[H]\)-模作张量积，直和的包含与投影仍互为分裂，得到所述直和。

每个 \(D\) 是 \([K:K\cap xHx^{-1}]\) 个左 \(H\)-陪集之并。选取这些陪集的代表元后，\(k[D]\) 是同样秩的自由右 \(k[H]\)-模；张量积因而是同样多个 \(W\) 的直和。
\end{proof}

这个分块不要求特征不整除群阶。它来自群元素基的划分，与表示是否完全可约无关；下节将在每一块内重新识别一个诱导表示。

\subsection{小群的双陪集计算}
\label{b5:subsec:0401:03}

先取 \(G=S_3\)、\(H=K=\bigl\langle(12)\bigr\rangle\)。左陪集 \(gH\) 可由点 \(g(3)\) 标记，\(H\) 在三点集上的轨道为 \(\{3\}\) 与 \(\{1,2\}\)。因此双陪集代表元可取 \(1\) 和 \(x=(13)\)。对应的交集分别为 \(H\) 和
\[
H\cap xHx^{-1}=\bigl\langle(12)\bigr\rangle\cap\bigl\langle(23)\bigr\rangle=\{1\}.
\]
两个双陪集的大小为 \(2\) 与 \(4\)，恰好覆盖 \(S_3\)。若 \(W=k\) 是平凡表示，则三维置换表示限制到 \(H\) 后分成一个固定点的平凡表示，以及两点自由轨道的正则表示。

\ProseParagraphBreak

再取 \(G=S_4\)，让 \(H=K=S_{\{1,2,3\}}\) 固定 \(4\)。\(G/H\) 同样可识别为四点集，轨道为 \(\{4\}\) 与 \(\{1,2,3\}\)。代表元取 \(1\) 和 \(x=(34)\)，后者的交集为固定 \(3,4\) 的二阶子群 \(\bigl\langle(12)\bigr\rangle\)。所以双陪集大小为 \(6\) 与 \(18\)，诱导平凡表示的两块维数为 \(1\) 与 \(3\)。

\ProseParagraphBreak

这些计算提供了一种实际方法：先寻找容易理解的陪集作用，再算轨道和稳定子群。直接枚举 \(kxh\) 往往会重复产生同一个元素，而轨道方法同时给出分块数、代表元和每块大小。

% !TeX root = ../../Book5.tex
% 纲要标题：### 4.2 Mackey 分解
% 纲要位置：工作记录/计划纲要.md，第 3825 行。
% 纲要细目（写作范围）：
% * 双陪集公式
% * 交织算子空间的计算
% * 不可约性判据


\section{Mackey 分解}
\label{b5:sec:0402}

双陪集分块中的 \(x\otimes W\) 受交集 \(K\cap xHx^{-1}\) 作用。这个作用不是原封不动地把 \(W\) 限制到交集：必须先通过 \(x\) 把交集搬回 \(H\)。将这一点写清以后，Mackey 公式便可由一个显式张量映射得到。

\subsection{双陪集公式}
\label{b5:subsec:0402:01}

\begin{definition}\label{b5:def:conjugate-representation}
设 \(G\) 为有限群，\(H\leq G\)，\(k\) 为域，\(\rho:H\to\operatorname{GL}_k(W)\) 为有限维表示，\(x\in G\)。在同一向量空间 \(W\) 上定义 \(xHx^{-1}\) 的作用
\[
\rho^x(a)=\rho(x^{-1}ax)\qquad(a\in xHx^{-1}).
\]
所得表示记为 \({}^xW\)，称为 \(W\) 的\term[gonge-biaoshi]{共轭表示}[conjugate representation]。
\end{definition}
\symidx{\({}^xW\)：共轭表示}

这个方向约定与左陪集相配合：在诱导表示中，
\(a(x\otimes w)=x\otimes\rho(x^{-1}ax)w\)。此外，先取 \({}^yW\) 再取其 \(x\)-共轭，所得表示为 \({}^{xy}W\)。

\begin{theorem}[Mackey 分解]\label{b5:thm:mackey-decomposition}
设 \(G\) 为有限群，\(H,K\leq G\)，\(k\) 为任意域，\(W\) 为有限维左 \(k[H]\)-模。取 \(K\backslash G/H\) 的一组代表元 \(X\)，并令 \(L_x=K\cap xHx^{-1}\)。则有左 \(k[K]\)-模同构
\begin{equation}\label{b5:eq:mackey-decomposition}
\operatorname{Res}_K^G\operatorname{Ind}_H^G W
\cong
\bigoplus_{x\in X}
\operatorname{Ind}_{L_x}^K
\operatorname{Res}_{L_x}^{xHx^{-1}}({}^xW).
\end{equation}
\end{theorem}
这一分解称为 \term[mackey-fenjie]{Mackey 分解}[Mackey decomposition]，公式也称为 Mackey 双陪集公式。
\begin{proof}
根据\cref{b5:prop:double-coset-blocks}，只须识别每一块。固定 \(x\)，把 \({}^xW\) 限制到 \(L_x\)，定义
\[
\begin{aligned}
\Phi_x:k[K]\otimes_{k[L_x]}{}^xW
&\longrightarrow k[KxH]\otimes_{k[H]}W,\\
a\otimes w&\longmapsto ax\otimes w.
\end{aligned}
\]
对 \(\ell\in L_x\)，有
\[
a\ell x\otimes w
=ax(x^{-1}\ell x)\otimes w
=ax\otimes\rho(x^{-1}\ell x)w.
\]
这正是源端平衡关系的像，故映射良定义；左 \(K\)-线性直接来自左乘。目标中的生成元 \(axh\otimes w\) 等于 \(\Phi_x(a\otimes hw)\)，所以映射满射。两端维数都为 \([K:L_x]\dim_k W\)，因而它是同构。对所有 \(x\) 取直和即得结论。
\end{proof}

选取别的代表元会改变上式各项的具体写法，却不改变\cref{b5:prop:double-coset-blocks}中由双陪集本身确定的子空间。公式的作用是把这些子空间写成熟悉的诱导模。

\subsection{交织算子空间的计算}
\label{b5:subsec:0402:02}

\begin{theorem}[交织空间公式]\label{b5:thm:mackey-intertwining}
设 \(G\) 为有限群，\(H,K\leq G\)，\(k\) 为域，\(U,W\) 分别为有限维左 \(k[K]\)-模、左 \(k[H]\)-模。取双陪集集 \(K\backslash G/H\) 的代表元集 \(X\)，对 \(x\in X\) 令 \(L_x=K\cap xHx^{-1}\)。有 \(k\) 向量空间同构
\begin{equation}\label{b5:eq:mackey-intertwining}
\begin{aligned}
&\operatorname{Hom}_{k[G]}
\left(\operatorname{Ind}_K^G U,\operatorname{Ind}_H^G W\right)\\
&\qquad\cong
\bigoplus_{x\in X}\operatorname{Hom}_{k[L_x]}
\bigl(\operatorname{Res}_{L_x}^K U,
\operatorname{Res}_{L_x}^{xHx^{-1}}({}^xW)\bigr).
\end{aligned}
\end{equation}
\end{theorem}
\begin{proof}
\cref{b5:thm:frobenius-reciprocity}把左侧识别为
\(\operatorname{Hom}_{k[K]}(U,\operatorname{Res}_K^G\operatorname{Ind}_H^G W)\)。
将\cref{b5:thm:mackey-decomposition}代入第二个变量，便得到各个
\(\operatorname{Hom}_{k[K]}(U,\operatorname{Ind}_{L_x}^K{}^xW)\) 的直和。

此时诱导出现在 Hom 的第二个变量，不能再次直接使用同一方向的左伴随。\cref{b5:thm:induction-coinduction}说明有限群的诱导与余诱导自然同构；\cref{b5:prop:coinduction-adjunction}给出的右伴随性质把这一项识别为
\(\operatorname{Hom}_{k[L_x]}(\operatorname{Res}_{L_x}^K U,{}^xW)\)。
这里每一步只用有限群的模构造，未用完全可约性。
\end{proof}

在复数域上，取特征标后，公式也计算两个诱导特征标的内积。取 \(K=H\)、\(U=W=\mathbf1_H\)，每一项都是一维空间，故
\[
\dim_{\mathbb C}\operatorname{End}_{\mathbb C[G]}
\bigl(\mathbb C[G/H]\bigr)=|H\backslash G/H|.
\]
例如，\(S_n\) 在 \(n\) 个点上的置换表示对应 \(H=S_{n-1}\)；当 \(n\geq2\) 时只有两个双陪集，所以它的自同态空间维数为 \(2\)。平凡表示出现一次，其余直和项的重数平方和因而为 \(1\)，得到一个不可约的 \((n-1)\) 维分量。

\subsection{不可约性判据}
\label{b5:subsec:0402:03}

\begin{theorem}[Mackey 不可约性判据]\label{b5:thm:mackey-irreducibility}
设 \(G\) 为有限群，\(H\leq G\)，\(W\) 为非零有限维复 \(H\)-表示。则 \(\operatorname{Ind}_H^G W\) 不可约，当且仅当 \(W\) 不可约，且对每个不等于 \(H\) 的双陪集 \(HxH\)，令 \(L_x=H\cap xHx^{-1}\) 后有
\[
\operatorname{Hom}_{\mathbb C[L_x]}
\bigl(\operatorname{Res}_{L_x}^H W,
\operatorname{Res}_{L_x}^{xHx^{-1}}({}^xW)\bigr)=0.
\]
\end{theorem}
\begin{proof}
\cref{b5:thm:maschke}说明有限群的复表示完全可约。若一个非零表示分解为互不同构不可约表示 \(V_i\) 的 \(\bigoplus_iV_i^{\oplus m_i}\)，则\cref{b5:prop:isotypic-decomposition}给出其自同态空间维数为 \(\sum_i m_i^2\)；这个数等于 \(1\)，恰好意味着表示不可约。

对\cref{b5:thm:mackey-intertwining}取 \(K=H\)、\(U=W\)，代表元 \(x=1\) 的项就是 \(\operatorname{End}_{\mathbb C[H]}W\)。所有直和项的维数非负，且这一项至少为 \(1\)。因此总维数为 \(1\)，当且仅当这一项维数为 \(1\) 而其余项都为零。再次使用完全可约性与上述维数公式，前一条件等价于 \(W\) 不可约。
\end{proof}

不可约性判据中的基域条件不能从前面的模同构公式中自动删去。在模特征下，表示可能不是完全可约的；此时自同态空间一维并不足以推出不可约。

\begin{example}\label{b5:exm:dihedral-induction}
设 \(G=D_{2n}=\langle r,s:r^n=s^2=1,\ srs=r^{-1}\rangle\)，其中 \(n\geq3\)，且 \(H=\langle r\rangle\)。令 \(\zeta\) 为本原 \(n\) 次单位根，取 \(\lambda_j(r)=\zeta^j\)。因为 \(H\) 正规，双陪集代表元为 \(1,s\)，交集均为 \(H\)，并且 \({}^s\lambda_j=\lambda_{-j}\)。所以
\(\operatorname{Ind}_H^G\lambda_j\) 不可约，当且仅当 \(2j\not\equiv0\pmod n\)。

\ProseParagraphBreak
在基 \(1\otimes1,s\otimes1\) 下，两个生成元的矩阵为
\[
r\longmapsto
\begin{pmatrix}\zeta^j&0\\0&\zeta^{-j}\end{pmatrix},
\qquad
s\longmapsto
\begin{pmatrix}0&1\\1&0\end{pmatrix}.
\]
当 \(2j\equiv0\) 时，向量 \((1,1)\) 和 \((1,-1)\) 各自生成不变直线，诱导表示分成两个一维表示。这也从矩阵上解释了判据中的例外。
\end{example}

% !TeX root = ../../Book5.tex
% 纲要标题：### 4.3 正规子群与 Clifford 理论
% 纲要位置：工作记录/计划纲要.md，第 3831 行。
% 纲要细目（写作范围）：
% * 不可约表示限制到正规子群
% * 共轭轨道与惯性群
% * Clifford 对应


\section{正规子群与 Clifford 理论}
\label{b5:sec:0403}

一个不可约表示限制到子群后通常会分解。子群正规时，共轭作用把这些分量联系起来：它们不能来自彼此无关的若干共轭轨道。本节在复数域上证明这一结构，并由其中一个等型分量恢复整个表示。

\subsection{不可约表示限制到正规子群}
\label{b5:subsec:0403:01}

先回顾\cref{b5:prop:isotypic-decomposition}中的等型分量。对有限群 \(N\) 的有限维复表示 \(V\)，以及不可约复 \(N\)-表示 \(W\)，令 \(V_W\) 为 \(V\) 中所有同构于 \(W\) 的 \(N\)-子模之和；若没有这种子模，就令 \(V_W=0\)。

\ProseParagraphBreak

由\cref{b5:thm:maschke}，\(V\) 是各个非零 \(V_W\) 的直和，而且 \(V_W\cong W^{\oplus e_W}\)。这个子空间不依赖所选的不可约直和分解：任意简单子模到不同类型分量的投影均为零，因为\cref{b5:lem:schur}说明不同不可约表示之间没有非零同态。同样的理由说明，任意 \(N\)-子模 \(U\subseteq V\) 都有
\begin{equation}\label{b5:eq:isotypic-submodule}
U=\bigoplus_{[W]}(U\cap V_W).
\end{equation}
事实上，先把 \(U\) 分解为简单子模，每个简单子模都只落在一种等型分量中。

\begin{theorem}[Clifford 限制定理]\label{b5:thm:clifford-restriction}
设 \(G\) 为有限群，\(N\trianglelefteq G\)，\(V\) 为不可约有限维复 \(G\)-表示，\(W\) 为 \(\operatorname{Res}_N^G V\) 的一个不可约成分。则限制表示中的不可约类型恰好组成 \(W\) 在 \(G\) 共轭作用下的一个轨道，而且所有类型的重数相等。
\end{theorem}
这一结论及其后续对应构成 \term[clifford-lilun]{Clifford 理论}[Clifford theory]的基本部分。
\begin{proof}
由于 \(N\) 正规，\({}^gW\) 仍是 \(N\)-表示。若 \(A\subseteq V\) 是同构于 \(W\) 的 \(N\)-子模，则
\[
n(ga)=g(g^{-1}ng)a\qquad(n\in N,\ a\in A)
\]
说明 \(gA\cong{}^gW\)。对所有这样的 \(A\) 求和，得到
\begin{equation}\label{b5:eq:clifford-component-transport}
gV_W=V_{{}^gW}.
\end{equation}
因此，同一共轭轨道对应的等型分量之和是非零 \(G\)-子模。\(V\) 不可约，迫使这个和等于 \(V\)，故不可能出现第二个轨道。

若 \(V_W\cong W^{\oplus e}\)，则等式右边由左边的线性同构得到
\(V_{{}^gW}\cong({}^gW)^{\oplus e}\)。所以各个类型出现相同的重数 \(e\)。
\end{proof}

证明中的完全可约性用于限制到 \(N\) 的表示和它的子模，由复数域上的 Maschke 定理保证。若要推广到其他特征，需要另行检验这一步。

\subsection{共轭轨道与惯性群}
\label{b5:subsec:0403:02}

\begin{definition}\label{b5:def:inertia-group}
设 \(G\) 为有限群，\(N\trianglelefteq G\)，\(W\) 为不可约有限维复 \(N\)-表示。子集
\[
I_G(W)=\{g\in G:{}^gW\cong W\}
\]
称为 \(W\) 在 \(G\) 中的\term[guanxing-qun]{惯性群}[inertia group]。
\end{definition}
\symidx{\(I_G(W)\)：不可约表示的惯性群}

共轭构成 \(G\) 在不可约同构类上的作用，所以 \(I_G(W)\) 是子群。每个 \(n\in N\) 都属于它：若 \(W\) 的作用为 \(\rho\)，则 \(\rho(n)\) 把 \({}^nW\) 交织到 \(W\)。因此 \(N\subseteq I_G(W)\)，轨道由左陪集 \(G/I_G(W)\) 标记。

\begin{corollary}\label{b5:cor:clifford-character}
设 \(G\) 为有限群，\(N\trianglelefteq G\)，\(V\) 为不可约有限维复 \(G\)-表示，\(W\) 为 \(\operatorname{Res}_N^G V\) 的一个不可约成分。令 \(I=I_G(W)\)，令 \(e\) 为 \(W\) 在该限制表示中的重数。若 \(\chi,\theta\) 分别是 \(V,W\) 的特征标，则
\[
\operatorname{Res}_N^G V\cong
\bigoplus_{gI\in G/I}({}^gW)^{\oplus e},
\qquad
\chi|_N=e\sum_{gI\in G/I}\theta^g,
\]
其中 \(\theta^g(n)=\theta(g^{-1}ng)\)、\(n\in N\)，各和取 \(G/I\) 的一组左陪集代表元。特别地，
\[
\dim_{\mathbb C}V=e[G:I]\dim_{\mathbb C}W.
\]
\end{corollary}
\begin{proof}
\cref{b5:thm:clifford-restriction}给出一个共轭轨道及相同重数。轨道中的类型恰由 \(G/I\) 标记，因而得到模分解。对作用矩阵取迹得到特征标等式，在单位元处取值得到维数等式。
\end{proof}

例如，\(S_3\) 的二维不可约表示限制到 \(N=A_3\) 后，是两个非平凡一维表示之和。若 \(\theta((123))=\exp(2\pi i/3)\)，则换位把 \(\theta\) 变成 \(\theta^{-1}\)，二者不同，因而 \(I_{S_3}(\theta)=A_3\)。这正是\cref{b5:exm:dihedral-induction}中 \(n=3\) 的情形。

\subsection{Clifford 对应}
\label{b5:subsec:0403:03}

\begin{theorem}[Clifford 对应]\label{b5:thm:clifford-correspondence}
设 \(G\) 为有限群，\(N\trianglelefteq G\)，\(W\) 为不可约有限维复 \(N\)-表示，\(I=I_G(W)\)。诱导给出以下两类有限维表示的同构类之间的双射：
\begin{enumerate}
\item 限制到 \(N\) 时含有 \(W\) 的不可约复 \(I\)-表示 \(U\)；
\item 限制到 \(N\) 时含有 \(W\) 的不可约复 \(G\)-表示 \(V\)。
\end{enumerate}
正向映射为 \(U\mapsto\operatorname{Ind}_I^G U\)，逆映射为 \(V\mapsto V_W\)，其中 \(V_W\) 为 \(\operatorname{Res}_N^G V\) 中的 \(W\)-等型分量。特别地，\(V_W\) 是不可约 \(I\)-表示，而自然映射
\begin{equation}\label{b5:eq:clifford-reconstruction}
\mathbb C[G]\otimes_{\mathbb C[I]}V_W\longrightarrow V,
\qquad g\otimes v\longmapsto gv
\end{equation}
是 \(G\)-同构。
\end{theorem}
\begin{proof}
先给定第二类中的 \(V\)。等式\eqref{b5:eq:clifford-component-transport}说明 \(V_W\) 是 \(I\)-不变的。若 \(0\neq U\subseteq V_W\) 是 \(I\)-子模，则选取 \(G/I\) 的代表元后，\(\bigoplus_{gI}gU\) 是 \(G\)-子模：任何群元素左乘代表元，只会改变代表元并带出一个 \(I\) 中的因子。这个和非零，所以等于 \(V\)。不同 \(gU\) 落在不同的 \(N\)-等型分量，取其中的 \(W\)-等型部分便有 \(U=V_W\)。故 \(V_W\) 不可约。

映射\eqref{b5:eq:clifford-reconstruction}的平衡关系由 \(I\)-作用保证。它的像包含所有 \(gV_W\)，故满射；源端的各个陪集分量映到目标中不同的等型分量，而且每一分量上的映射均为线性同构。因此它也单射。

反过来，设 \(U\) 属于第一类。对有限群 \(I\) 及其正规子群 \(N\) 应用\cref{b5:thm:clifford-restriction}，由于 \(I\) 固定 \(W\) 的同构类，得到 \(U|_N\cong W^{\oplus e}\)。令 \(Y=\operatorname{Ind}_I^G U\)。作为 \(N\)-模，它的陪集分解为
\[
Y=\bigoplus_{gI\in G/I}g\otimes U,
\]
每一项的类型是 \({}^gW\)，不同陪集的类型互不相同。

若 \(0\neq Z\subseteq Y\) 是 \(G\)-子模，等式\eqref{b5:eq:isotypic-submodule}说明它与某个 \(g\otimes U\) 有非零交。用 \(g^{-1}\) 作用，把非零交移到 \(1\otimes U\)。于是 \(Z\cap(1\otimes U)\) 是非零 \(I\)-子模，必等于 \(1\otimes U\)。再用 \(G\) 作用，得到 \(Z=Y\)。所以 \(Y\) 不可约。它的 \(W\)-等型分量恰为 \(1\otimes U\)，故两种构造互为逆；同构也在取等型分量时被保留。
\end{proof}

Clifford 对应把问题从 \(G\) 缩到惯性群 \(I\)。当 \(I=N\) 时，\(W\) 本身已经诱导出不可约 \(G\)-表示；当 \(I=G\) 时，限制只含一种类型，但重数 \(e\) 仍可能大于 \(1\)。后一种情形需要研究 \(W\) 能否延拓，以及重数空间上究竟有怎样的作用。

% !TeX root = ../../Book5.tex
% 纲要标题：### 4.4 射影表示初步
% 纲要位置：工作记录/计划纲要.md，第 3837 行。
% 纲要细目（写作范围）：
% * 延拓表示的障碍
% * 因子集与2-余圈
% * 与第四卷群上同调和第二卷交叉积的联系
% ---


\section{射影表示初步}
\label{b5:sec:0404}

惯性群中的元素都把 \(W\) 变成同构的表示，但为每个元素分别选一个同构，并不保证这些同构按照群乘法相乘。差出的标量恰好记下延拓的障碍。本节直接由这些标量建立所需理论，不预先要求群上同调。

\subsection{表示延拓的障碍}
\label{b5:subsec:0404:01}

\begin{definition}\label{b5:def:representation-extension}
设 \(I\) 为有限群，\(N\trianglelefteq I\)，\(W\) 为有限维复向量空间，\(\rho:N\to\operatorname{GL}_{\mathbb C}(W)\) 为表示。若表示 \(\widetilde\rho:I\to\operatorname{GL}_{\mathbb C}(W)\) 在 \(N\) 上等于 \(\rho\)，则称它为 \(\rho\) 到 \(I\) 的\term[biaoshi-yantuo]{延拓}[extension of a representation]。
\end{definition}

延拓保持原来的向量空间，维数不变；它与一般会增大维数的诱导是不同的构造。如果 \(\rho\) 能延拓，则对 \(g\in I\)，矩阵 \(\widetilde\rho(g)\) 实现 \(W\cong{}^gW\)，故 \(W\) 的同构类必须被 \(I\) 固定。下面精确考察这个必要条件还差什么。

\begin{theorem}[延拓的标量障碍]\label{b5:thm:extension-obstruction}
设 \(I\) 为有限群，\(N\trianglelefteq I\)，\(\rho:N\to\operatorname{GL}_{\mathbb C}(W)\) 为不可约有限维复表示，并假定 \({}^gW\cong W\) 对所有 \(g\in I\) 成立。令 \(Q=I/N\)，选取商映射的集合截面 \(s:Q\to I\)，满足 \(s(1)=1\)，并对 \(p,q\in Q\) 记
\[
c(p,q)=s(p)s(q)s(pq)^{-1}\in N.
\]
可以选取 \(T_q\in\operatorname{GL}_{\mathbb C}(W)\)，使 \(T_1=1_W\)，且
\begin{equation}\label{b5:eq:extension-intertwiner}
T_q\rho(n)T_q^{-1}
=\rho\bigl(s(q)ns(q)^{-1}\bigr)\qquad(n\in N).
\end{equation}
存在唯一标量 \(\alpha(p,q)\in\mathbb C^\times\)，使
\begin{equation}\label{b5:eq:extension-factor}
T_pT_q=\alpha(p,q)\rho\bigl(c(p,q)\bigr)T_{pq}.
\end{equation}
这些标量满足
\begin{equation}\label{b5:eq:factor-cocycle}
\alpha(1,q)=\alpha(q,1)=1,\qquad
\alpha(p,q)\alpha(pq,r)=\alpha(q,r)\alpha(p,qr)\quad(p,q,r\in Q).
\end{equation}
表示 \(W\) 能延拓到 \(I\)，当且仅当存在 \(b:Q\to\mathbb C^\times\)、\(b(1)=1\)，使
\begin{equation}\label{b5:eq:extension-rescaling}
b(p)b(q)\alpha(p,q)=b(pq)\qquad(p,q\in Q).
\end{equation}
\end{theorem}
\begin{proof}
共轭不变性给出满足\eqref{b5:eq:extension-intertwiner}的可逆交织映射。\(T_pT_q\) 和 \(\rho\bigl(c(p,q)\bigr)T_{pq}\) 都实现由 \(s(p)s(q)\) 所给的共轭，所以二者之比与 \(\rho(N)\) 对易。\cref{b5:cor:schur-algebraically-closed}说明该比值为唯一非零标量，得到\eqref{b5:eq:extension-factor}。

每个 \(g\in I\) 唯一写成 \(ns(q)\)。定义可逆算子
\[
P\bigl(ns(q)\bigr)=\rho(n)T_q.
\]
若 \(g=ns(p)\)、\(h=ms(q)\)，则
\[
\begin{aligned}
P(g)P(h)
&=\rho\bigl(n s(p)m s(p)^{-1}\bigr)T_pT_q\\
&=\alpha(p,q)
\rho\bigl(n s(p)m s(p)^{-1}c(p,q)\bigr)T_{pq}\\
&=\alpha(p,q)P(gh).
\end{aligned}
\]
对三个元素比较 \(\bigl(P(g)P(h)\bigr)P(\ell)\) 与 \(P(g)\bigl(P(h)P(\ell)\bigr)\)，并消去可逆的 \(P(gh\ell)\)，就得到\eqref{b5:eq:factor-cocycle}中的第二个等式；单位元处的等式来自 \(T_1=1_W\)。

如果存在 \(b\) 满足\eqref{b5:eq:extension-rescaling}，则
\(\widetilde\rho\bigl(ns(q)\bigr)=\rho(n)b(q)T_q\) 按上述乘法计算是群表示，而且在 \(N\) 上等于 \(\rho\)。反过来，若有延拓，则 \(\widetilde\rho\bigl(s(q)\bigr)\) 与 \(T_q\) 实现同一个共轭。Schur 引理给出 \(\widetilde\rho\bigl(s(q)\bigr)=b(q)T_q\)，再比较两者乘法便得到\eqref{b5:eq:extension-rescaling}。
\end{proof}

不能把“不变”直接换成“可延拓”。取
\[
I=Q_8=\{\pm1,\pm i,\pm j,\pm ij\},\qquad
N=\{\pm1\},
\]
并令 \(\rho(-1)=-1\)。\(N\) 在中心，故这个一维表示被全部 \(I\) 固定。但任何一维 \(Q_8\)-表示都把交换子送到 \(1\)，而 \(iji^{-1}j^{-1}=-1\)，所以它不能延拓。这个障碍在下一小节可以直接算成一个不对称的标量函数。

\subsection{因子集与 2-余圈}
\label{b5:subsec:0404:02}

\begin{definition}\label{b5:def:projective-representation}
设 \(Q\) 为有限群，\(E\neq0\) 为有限维复向量空间。若映射 \(R:Q\to\operatorname{GL}_{\mathbb C}(E)\) 满足 \(R(1)=1_E\)，并对每个 \(p,q\in Q\) 存在非零标量 \(\alpha(p,q)\)，使
\[
R(p)R(q)=\alpha(p,q)R(pq),
\]
则称 \(R\) 为 \(Q\) 的\term[sheying-biaoshi]{射影表示}[projective representation]，称 \(\alpha\) 为所选矩阵的\term[yinziji]{因子集}[factor set]。若 \(E\) 没有非零真子空间同时被所有 \(R(q)\) 保持，则称这个射影表示不可约。
\end{definition}

把矩阵模去非零标量后，射影表示给出群同态
\(Q\to\operatorname{PGL}(E)\)。反之，这样的群同态选取矩阵代表后就给出因子集；代表不同，因子集也可能不同。

\begin{definition}\label{b5:def:scalar-two-cocycle}
设 \(Q\) 为有限群，\(\alpha:Q\times Q\to\mathbb C^\times\) 为函数。若对所有 \(p,q,r\in Q\) 有
\[
\alpha(1,q)=\alpha(q,1)=1,\qquad
\alpha(p,q)\alpha(pq,r)=\alpha(q,r)\alpha(p,qr),
\]
则称 \(\alpha\) 为系数作用平凡的\term[guiyihua-eryuquan]{归一化 \(2\)-余圈}[normalized 2-cocycle]。若存在 \(a:Q\to\mathbb C^\times\)、\(a(1)=1\)，使
\[
\alpha(p,q)=\frac{a(p)a(q)}{a(pq)}\qquad(p,q\in Q),
\]
则称 \(\alpha\) 为一个\term[eryubian]{\(2\)-余边}[2-coboundary]。
\end{definition}

矩阵乘法的结合律说明每个因子集都是归一化 \(2\)-余圈。将 \(R(q)\) 改成 \(b(q)R(q)\)，因子集变成
\begin{equation}\label{b5:eq:factor-coboundary-change}
\alpha'(p,q)=\frac{b(p)b(q)}{b(pq)}\alpha(p,q).
\end{equation}
所以\cref{b5:thm:extension-obstruction}可表述为：延拓存在，当且仅当所得因子集是一个 \(2\)-余边。

\begin{proposition}\label{b5:prop:factor-choice-independence}
设 \(I\) 为有限群，\(N\trianglelefteq I\)，\(\rho:N\to\operatorname{GL}_{\mathbb C}(W)\) 为同构类被 \(I\) 固定的不可约有限维复表示。令 \(Q=I/N\)，取商映射的截面 \(s:Q\to I\)、\(s(1)=1\)，以及算子 \(T_q\in\operatorname{GL}_{\mathbb C}(W)\)，使 \(T_1=1_W\) 且 \(T_q\rho(n)T_q^{-1}=\rho(s(q)ns(q)^{-1})\) 对 \(q\in Q,n\in N\) 成立。对 \(p,q\in Q\)，记 \(c(p,q)=s(p)s(q)s(pq)^{-1}\)，由 \(T_pT_q=\alpha(p,q)\rho(c(p,q))T_{pq}\) 确定因子集 \(\alpha:Q\times Q\to\mathbb C^\times\)。更换交织算子或截面后，新因子集 \(\alpha'\) 满足 \(\alpha'/\alpha\) 是一个 \(2\)-余边。因此，“能否延拓”的判断不依赖这些选择。
\end{proposition}
\begin{proof}
固定截面时，任意另一组交织算子由 Schur 引理写成 \(T'_q=b(q)T_q\)，其因子集由\eqref{b5:eq:factor-coboundary-change}给出。

若改用 \(s'(q)=d(q)s(q)\)，其中 \(d(q)\in N\)、\(d(1)=1\)，先选 \(T'_q=\rho\bigl(d(q)\bigr)T_q\)。对同一个群元素，写成 \(ns(q)=n's'(q)\) 就有 \(n=n'd(q)\)，从而 \(\rho(n)T_q=\rho(n')T'_q\)。所以证明中定义的 \(P(g)\) 完全不变，因子集也不变。再作其他交织算子选择时，才按前一段改变一个余边。
\end{proof}

回到 \(Q_8\)。把 \(Q=Q_8/N\) 写成 \(\mathbb F_2^2\)，取 \(s(a,b)=i^aj^b\) 及 \(T_{(a,b)}=1\)。这里
\[
\alpha\bigl((a,b),(c,d)\bigr)=(-1)^{ac+bd+bc}.
\]
指数按模 \(2\) 计算。这是双线性式，所以直接展开也可验证余圈恒等式。取 \(p=(1,0)\)、\(q=(0,1)\)，得到 \(\alpha(p,q)=1\) 而 \(\alpha(q,p)=-1\)。交换群上的余边 \(a(p)a(q)/a(pq)\) 对称，故这个因子集不是余边。

\subsection{与群上同调、交叉积的联系}
\label{b5:subsec:0404:03}

把所有归一化 \(2\)-余圈按逐点相乘组成交换群，再模去 \(2\)-余边子群，就得到记为 \(H^2(Q,\mathbb C^\times)\) 的群。这里 \(\mathbb C^\times\) 上的 \(Q\)-作用是平凡的。它是群上同调的一个具体低阶形式；当前的延拓障碍就是 \([\alpha]\)，不需要借助消解或导出函子才能计算。

\ProseParagraphBreak

同一个因子集还可写成一个代数的乘法。给定有限群 \(Q\) 和归一化 \(2\)-余圈 \(\alpha\)，在以 \(u_q\) 为基的复向量空间上定义
\[
u_pu_q=\alpha(p,q)u_{pq}.
\]
归一化条件使 \(u_1\) 为单位元，余圈恒等式恰好使乘法结合；各个 \(u_q\) 都可逆。这个代数称为\term[niuqu-qundaisu]{扭曲群代数}[twisted group algebra]，记为 \(\mathbb C^\alpha[Q]\)。它是系数域上作用平凡的交叉积。给出一个非零有限维左 \(\mathbb C^\alpha[Q]\)-模，等价于给出因子集为 \(\alpha\) 的射影表示：令 \(u_q\) 按 \(R(q)\) 作用即可。重标定基 \(u_q\) 与\eqref{b5:eq:factor-coboundary-change}完全对应。
\symidx{\(\mathbb C^\alpha[Q]\)：扭曲群代数}

\begin{theorem}[重数空间上的射影表示]\label{b5:thm:clifford-projective-multiplicity}
设 \(I\) 为有限群，\(N\trianglelefteq I\)，\(\rho:N\to\operatorname{GL}_{\mathbb C}(W)\) 为同构类被 \(I\) 固定的不可约有限维复表示。令 \(Q=I/N\)，取商映射的截面 \(s:Q\to I\)、\(s(1)=1\)，以及算子 \(T_q\in\operatorname{GL}_{\mathbb C}(W)\)，使 \(T_1=1_W\) 且
\[
T_q\rho(n)T_q^{-1}=\rho(s(q)ns(q)^{-1})\qquad(n\in N,\ q\in Q).
\]
对 \(p,q\in Q\)，记 \(c(p,q)=s(p)s(q)s(pq)^{-1}\)，以
\[
T_pT_q=\alpha(p,q)\rho(c(p,q))T_{pq}
\]
确定因子集 \(\alpha:Q\times Q\to\mathbb C^\times\)。给定正整数 \(e\)，满足 \(\operatorname{Res}_N^I U\cong W^{\oplus e}\) 的非零有限维复 \(I\)-表示 \(U\)，等价地由 \(e\) 维复向量空间 \(E\) 上因子集为 \(\alpha^{-1}\) 的射影 \(Q\)-表示 \(S:Q\to\operatorname{GL}_{\mathbb C}(E)\) 给出，其作用为
\begin{equation}\label{b5:eq:clifford-projective-tensor}
ns(q)\longmapsto \rho(n)T_q\otimes S(q)
\quad\hbox{作用在 }W\otimes_{\mathbb C}E\hbox{ 上}.
\end{equation}
而且 \(U\) 不可约，当且仅当 \(S\) 不可约。
\end{theorem}
\begin{proof}
先给定 \(U\)，记其表示为 \(\pi\)，并令
\(E=\operatorname{Hom}_{\mathbb C[N]}(W,U)\)。由于 \(U|_N\cong W^{\oplus e}\) 且 \(\operatorname{End}_N(W)=\mathbb C\)，求值映射
\[
W\otimes E\longrightarrow U,\qquad w\otimes f\longmapsto f(w)
\]
由\cref{b5:prop:isotypic-decomposition}是 \(N\)-同构，且 \(\dim E=e\)。定义
\[
S(q)f=\pi\bigl(s(q)\bigr)\,f\,T_q^{-1}.
\]
由\eqref{b5:eq:extension-intertwiner}，对 \(n\in N\) 有
\[
\begin{aligned}
\bigl(S(q)f\bigr)\rho(n)
&=\pi\bigl(s(q)\bigr)f
\rho\bigl(s(q)^{-1}ns(q)\bigr)T_q^{-1}\\
&=\pi(n)\pi\bigl(s(q)\bigr)fT_q^{-1}
=\pi(n)\bigl(S(q)f\bigr).
\end{aligned}
\]
所以此复合仍是 \(N\)-线性映射，确实属于 \(E\)。求值映射把 \(T_q\otimes S(q)\) 的作用变成 \(\pi\bigl(s(q)\bigr)\)，从而得到\eqref{b5:eq:clifford-projective-tensor}。

将这个等式对 \(p,q\) 相乘。左侧的群作用相乘为 \(\pi\bigl(c(p,q)\bigr)\pi\bigl(s(pq)\bigr)\)，右侧的第一个张量因子满足\eqref{b5:eq:extension-factor}。消去可逆算子 \(\rho\bigl(c(p,q)\bigr)T_{pq}\)，得到
\[
S(p)S(q)=\alpha(p,q)^{-1}S(pq).
\]
反过来，若 \(S\) 满足此等式，则\eqref{b5:eq:clifford-projective-tensor}中两个因子产生的标量互相抵消，因而给出普通的 \(I\)-表示；在 \(N\) 上它是 \(\rho(n)\otimes1_E\)。

最后，\(W\otimes E\) 的每个 \(N\)-子模都是 \(W\otimes F\) 的形式，其中 \(F\subseteq E\)：将该子模分解为 \(W\) 的直和，并对它应用同一个求值同构即可。由于每个 \(T_q\) 可逆，子空间 \(W\otimes F\) 被 \(I\) 保持，当且仅当 \(F\) 被所有 \(S(q)\) 保持。这就证明了不可约性的等价。
\end{proof}

结合\cref{b5:thm:clifford-correspondence}，先在惯性群中求这些射影表示，再诱导到 \(G\)，便得到限制中含有 \(W\) 的全部不可约表示。若障碍为零，可以重标定 \(T_q\)，使它们给出 \(W\) 的延拓；此时 \(\alpha=1\)，重数空间上的表示也成为普通的商群表示。


\begin{chaptersummary}
双陪集 \(K\backslash G/H\) 描述 \(K\) 在 \(G/H\) 上的轨道，交集 \(K\cap xHx^{-1}\) 则是相应点的稳定子群。Mackey 分解把这一集合分解提升为任意域上的模同构；配合诱导与余诱导的伴随性质，便可把两个诱导表示之间的交织空间化为交集子群上的计算。在复数域上，自同态空间的维数又给出不可约性判据。

\ProseParagraphBreak

当 \(N\) 正规时，一个不可约 \(G\)-表示限制到 \(N\) 后，只出现一个共轭轨道，轨道中各类型的重数相等。Clifford 对应说明，其中一个等型分量是惯性群的不可约表示，诱导它就能恢复原表示。惯性群固定不可约类型，却未必允许在原空间上延拓作用；这种差别由一个标量 \(2\)-余圈记录。重数空间上的射影表示以该余圈的倒数为因子集，使两个射影作用的标量相互抵消，最终形成普通表示。
\end{chaptersummary}

\BookExercises

\begin{exercise}\label{b5:ex:04-normal-double-cosets}
设 \(G\) 为有限群，\(H\trianglelefteq G\)，\(K\leq G\)。证明 \(KH\) 是子群，且每个双陪集 \(KxH\) 等于右陪集 \(KHx\)。求双陪集个数以及每块的大小，并解释为什么所有交集子群相同。
\end{exercise}
\begin{solution}
正规性给出
\((k_1h_1)(k_2h_2)=k_1k_2(k_2^{-1}h_1k_2)h_2\in KH\)，逆元也可写成 \(K\) 中元素乘 \(H\) 中元素，所以 \(KH\leq G\)。由于 \(xHx^{-1}=H\)，有 \(KxH=KHx\)。这实际上是 \(KH\) 的右陪集，因而双陪集个数为 \([G:KH]\)，每块大小为 \(|KH|=|K|\cdot|H|/|K\cap H|\)。交集 \(K\cap xHx^{-1}=K\cap H\) 不依赖 \(x\)。
\end{solution}

\begin{exercise}\label{b5:ex:04-dihedral-double-cosets}
在 \(D_8=\langle r,s:r^4=s^2=1,\ srs=r^{-1}\rangle\) 中取 \(H=\langle s\rangle\)。求 \(H\backslash D_8/H\)。若 \(\varepsilon(s)=-1\)，把 \(\operatorname{Res}_H^{D_8}\operatorname{Ind}_H^{D_8}\varepsilon\) 分解为不可约复 \(H\)-表示。
\end{exercise}
\begin{solution}
三个双陪集为 \(H\)、\(HrH=\{r,r^3,rs,r^3s\}\)、\(Hr^2H=r^2H\)，大小为 \(2,4,2\)。代表元 \(1,r,r^2\) 对应的交集分别为 \(H,\{1\},H\)，因为 \(rsr^{-1}=r^2s\)，而 \(r^2\) 在中心。Mackey 分解给出
\[
\varepsilon\oplus\operatorname{Ind}_{\{1\}}^H\mathbb C\oplus\varepsilon
\cong\mathbf1_H\oplus\varepsilon^{\oplus3}.
\]
中间一项是二阶群的正则表示，其两个特征向量分别为两个基向量之和与之差。
\end{solution}

\begin{exercise}\label{b5:ex:04-unequal-subgroups}
令 \(G=S_4\)，\(H\) 为固定 \(4\) 的子群，\(K=\bigl\langle(12)\bigr\rangle\)。求 \(K\backslash G/H\) 的代表元和交集子群，并计算
\(\dim_k\operatorname{Hom}_{k[G]}\bigl(k[G/K],k[G/H]\bigr)\)，其中 \(k\) 是任意域。
\end{exercise}
\begin{solution}
把 \(G/H\) 识别为四点集，\(K\) 的轨道为 \(\{1,2\}\)、\(\{3\}\)、\(\{4\}\)。可取代表元 \((14),(34),1\)。第一点的稳定子群为 \(\{1\}\)，后两点的稳定子群为 \(K\)，所以这些也是对应的交集子群。对\cref{b5:thm:mackey-intertwining}取两个平凡表示，每个直和项均为 \(\operatorname{Hom}_{k[L_x]}(k,k)=k\)，故维数为 \(3\)，不受 \(k\) 的特征影响。
\end{solution}

\begin{exercise}\label{b5:ex:04-modular-mackey}
设 \(k\) 的特征为 \(3\)，\(G=S_3\)，\(H=\bigl\langle(12)\bigr\rangle\)，\(N=A_3\)。用 Mackey 公式计算 \(\operatorname{Res}_N^G\operatorname{Ind}_H^G k\)，并证明所得 \(N\)-模不是完全可约的。
\end{exercise}
\begin{solution}
\(N\) 在三点集 \(G/H\) 上传递，且稳定子群为 \(N\cap H=\{1\}\)。因此只有一个双陪集，限制表示同构于 \(k[N]\)。若 \(a=(123)\)，则
\[
k[N]\cong k[t]/(t^3-1)=k[t]/(t-1)^3.
\]
乘 \(t-1\) 是非零幂零算子。在任何简单 \(k[N]\)-模上，它的核非零且为子模，故它只能为零；所以若正则模完全可约，\(t-1\) 必在整个正则模上为零，与上述计算矛盾。这说明 Mackey 同构本身并不以半单性为前提。
\end{solution}

\begin{exercise}\label{b5:ex:04-trivial-intersection-obstruction}
设 \(G\) 为有限群，\(H<G\)。若存在 \(x\notin H\) 使 \(H\cap xHx^{-1}=\{1\}\)，证明对每个非零有限维复 \(H\)-表示 \(W\)，\(\operatorname{Ind}_H^G W\) 都可约。
\end{exercise}
\begin{solution}
在交织空间公式中，\(x=1\) 的项至少一维，所给 \(x\) 的项是底层向量空间之间的 \(\operatorname{Hom}_{\mathbb C}(W,W)\)，维数为 \((\dim W)^2>0\)。所以诱导表示的自同态空间维数至少为 \(2\)。它是完全可约的，故由\cref{b5:prop:isotypic-decomposition}不能不可约。这里不必事先假设 \(W\) 不可约。
\end{solution}

\begin{exercise}\label{b5:ex:04-normal-induction-end}
设 \(G\) 为有限群，\(N\trianglelefteq G\)，\(W\) 为不可约复 \(N\)-表示，\(I=I_G(W)\)。证明
\[
\dim_{\mathbb C}\operatorname{End}_G
\left(\operatorname{Ind}_N^G W\right)=[I:N].
\]
若 \(W'\) 是另一个不可约复 \(N\)-表示，证明 \(\operatorname{Ind}_N^G W\) 与 \(\operatorname{Ind}_N^G W'\) 之间存在非零同态，当且仅当 \(W'\) 与 \(W\) 共轭。
\end{exercise}
\begin{solution}
双陪集 \(N\backslash G/N\) 就是 \(G/N\)，每个交集都是 \(N\)。交织公式中的项为 \(\operatorname{Hom}_N(W,{}^gW)\)，由 Schur 引理，它在 \(g\in I\) 时一维，在其他情形为零。故总维数为 \([I:N]\)。对两个不同的表示应用同一公式，直和项为 \(\operatorname{Hom}_N(W,{}^gW')\)，其非零条件恰为 \(W\cong{}^gW'\)，得到第二个断言。
\end{solution}

\begin{exercise}\label{b5:ex:04-nonnormal-restriction}
设 \(V=\left\{(z_1,z_2,z_3,z_4)\in\mathbb C^4:\sum_i z_i=0\right\}\)，赋予 \(S_4\) 的坐标置换作用。证明 \(V\) 不可约，而它限制到固定 \(4\) 的 \(S_3\) 后出现维数不同的不可约成分。说明这与 Clifford 限制定理没有矛盾。
\end{exercise}
\begin{solution}
\(S_4\) 在四点上的置换表示包含一次平凡表示，且由双陪集计算，其自同态空间维数为 \(2\)。故其补空间 \(V\) 不可约。限制到 \(S_3\) 时，有分解
\[
V=\mathbb C(1,1,1,-3)
\oplus\bigl\{(z_1,z_2,z_3,0):z_1+z_2+z_3=0\bigr\}.
\]
前一项平凡，后一项是三点置换表示去掉平凡项后的二维不可约表示，其不可约性由相同的两双陪集计算得到。这个 \(S_3\) 不是 \(S_4\) 的正规子群，所以不可约成分不必形成同维的共轭轨道。
\end{solution}

\begin{exercise}\label{b5:ex:04-s4-klein-four}
令 \(N=\{1,a,b,c\}\trianglelefteq S_4\)，其中
\(a=(12)(34)\)、\(b=(13)(24)\)、\(c=(14)(23)\)。取上题的三维表示 \(V\)，并令 \(\lambda(a)=1\)、\(\lambda(b)=\lambda(c)=-1\)。求 \(V_\lambda\) 和 \(I_{S_4}(\lambda)\)，并把 \(V\) 写成一个一维表示的诱导。
\end{exercise}
\begin{solution}
向量 \(v=(1,1,-1,-1)\) 满足 \(av=v\)、\(bv=cv=-v\)，故 \(\mathbb Cv\subseteq V_\lambda\)。对另两个分拆 \(13|24\)、\(14|23\) 取对应的正负向量，得到三条线性无关的特征直线，因此 \(V_\lambda=\mathbb Cv\)，其余两条是另外两个非平凡特征标的等型分量。

惯性群必须固定 \(\ker\lambda=\{1,a\}\)，等价于在共轭下固定 \(a\)。它就是保持无序分拆 \(\bigl\{\{1,2\},\{3,4\}\bigr\}\) 的子群 \(I\)，阶为 \(8\)。\(I\) 在 \(\mathbb Cv\) 上的特征标 \(\eta\) 对保持两块的元素取 \(1\)，对交换两块的元素取 \(-1\)。Clifford 对应给出 \(V\cong\operatorname{Ind}_I^{S_4}\eta\)，维数为 \([S_4:I]=3\)。
\end{solution}

\begin{exercise}\label{b5:ex:04-quaternion-restrictions}
在 \(\mathbb C^2\) 上令
\[
\rho(i)=\begin{pmatrix}\mathrm i&0\\0&-\mathrm i\end{pmatrix},
\qquad
\rho(j)=\begin{pmatrix}0&1\\-1&0\end{pmatrix},
\]
其中 \(\mathrm i^2=-1\)。证明它们给出 \(Q_8\) 的不可约表示，并分别对正规子群 \(\langle i\rangle\) 和 \(\{\pm1\}\) 计算 Clifford 轨道、惯性群和重数。
\end{exercise}
\begin{solution}
两个矩阵的平方都是 \(-I_2\)，且反对易，因此满足 \(Q_8\) 的关系。第一个矩阵的两个特征值不同，任何共同不变直线只能是其一条坐标特征线，但第二个矩阵交换这两条线，所以表示不可约。

对 \(\langle i\rangle\)，限制为特征标 \(i\mapsto\mathrm i\) 与 \(i\mapsto-\mathrm i\) 之和。\(j\) 交换二者，惯性群为 \(\langle i\rangle\)，重数 \(e=1\)。对中心 \(\{\pm1\}\)，只有特征标 \(-1\mapsto-1\)，出现重数 \(e=2\)，其惯性群为全部 \(Q_8\)。同一个表示的轨道长度和重数会随正规子群的选择而改变。
\end{solution}

\begin{exercise}\label{b5:ex:04-central-subgroup}
设 \(G\) 为有限群，\(N\leq Z(G)\)，\(V\) 为不可约复 \(G\)-表示。证明存在唯一一维特征标 \(\lambda:N\to\mathbb C^\times\)，使 \(n\) 在 \(V\) 上按 \(\lambda(n)I_V\) 作用。给出它与 Clifford 限制公式中 \(e\) 和 \(I\) 的关系。
\end{exercise}
\begin{solution}
每个 \(\rho(n)\) 与整个 \(\rho(G)\) 对易，故\cref{b5:cor:schur-algebraically-closed}使它为唯一标量 \(\lambda(n)\)。群乘法说明 \(\lambda(nm)=\lambda(n)\lambda(m)\)，所以它是一维特征标。限制为 \(\lambda^{\oplus\dim V}\)，从而 \(I=G\)、\(e=\dim V\)。中心上的同构类型虽然只有一个，重数仍可大于 \(1\)，上题的 \(Q_8\) 就给出例子。
\end{solution}

\begin{exercise}\label{b5:ex:04-extension-torsor}
设 \(N\trianglelefteq I\) 为有限群的正规子群，\(\rho\) 为不可约复 \(N\)-表示，并已有一个延拓 \(\widetilde\rho\)。证明所有延拓都唯一写成
\[
g\longmapsto\lambda(gN)\widetilde\rho(g),
\]
其中 \(\lambda:I/N\to\mathbb C^\times\) 为一维特征标。
\end{exercise}
\begin{solution}
给定另一延拓 \(\sigma\)，\(\sigma(g)\) 与 \(\widetilde\rho(g)\) 实现同一个 \(N\) 的共轭作用，因此它们的比值与 \(\rho(N)\) 对易。Schur 引理给出唯一标量 \(\lambda(g)\)，使 \(\sigma(g)=\lambda(g)\widetilde\rho(g)\)。比较 \(gh\) 上的乘法得到 \(\lambda(gh)=\lambda(g)\lambda(h)\)，在 \(N\) 上两延拓相同使 \(\lambda|_N=1\)，故它下降为 \(I/N\) 的特征标。反过来，用任意这样的特征标相乘仍是延拓，且表达式唯一。
\end{solution}

\begin{exercise}\label{b5:ex:04-cyclic-quotient-extension}
设 \(I\) 为有限群，\(N\trianglelefteq I\)，且 \(I/N\) 为阶 \(d\) 的循环群。证明每个同构类被 \(I\) 固定的不可约复 \(N\)-表示都能延拓到 \(I\)。
\end{exercise}
\begin{solution}
选取 \(t\in I\)，使 \(tN\) 生成 \(I/N\)，并选可逆 \(T\) 满足 \(T\rho(n)T^{-1}=\rho(tnt^{-1})\)。\(T^d\) 与 \(\rho(t^d)\) 实现同一个共轭，所以 \(T^d=a\cdot\rho(t^d)\)，其中 \(a\in\mathbb C^\times\)。取 \(b\in\mathbb C^\times\) 满足 \(b^d=a^{-1}\)，令 \(S=bT\)，便有 \(S^d=\rho(t^d)\)。

每个群元素唯一写成 \(nt^j\)，其中 \(0\leq j<d\)。定义 \(\widetilde\rho(nt^j)=\rho(n)S^j\)。乘法中先用 \(S^j\rho(m)S^{-j}=\rho(t^jmt^{-j})\) 移动 \(N\) 中的因子；若指数越过 \(d\)，再以 \(S^d=\rho(t^d)\) 作相同的归并。于是 \(\widetilde\rho(gh)=\widetilde\rho(g)\widetilde\rho(h)\)，且它在 \(N\) 上等于 \(\rho\)。复数中 \(d\) 次根的存在是这一步调整所用的条件。
\end{solution}

\begin{exercise}\label{b5:ex:04-invariant-needed}
设 \(G=S_3\)，\(N=A_3\)，\(\lambda((123))=\zeta\)，其中 \(\zeta\neq1\) 为三次单位根。证明 \(\lambda\) 不能延拓到 \(G\)，并说明上一题的循环商群条件为何没有给出矛盾。
\end{exercise}
\begin{solution}
换位把 \((123)\) 共轭为其逆，因此把 \(\lambda\) 变成不同的 \(\lambda^{-1}\)。若 \(\lambda\) 有一维延拓，则共轭在一维标量矩阵上不起作用，迫使 \(\zeta=\zeta^{-1}\)，矛盾。虽然 \(G/N\cong C_2\) 为循环群，但上一题还要求不可约类型被整个 \(G\) 固定；这里惯性群只有 \(N\)。
\end{solution}

\begin{exercise}\label{b5:ex:04-twisted-klein-algebra}
令 \(Q=\mathbb F_2^2\)，并定义
\(\alpha\bigl((a,b),(c,d)\bigr)=(-1)^{bc}\)。证明 \(\alpha\) 是归一化 \(2\)-余圈但不是余边。再令
\[
X=\begin{pmatrix}1&0\\0&-1\end{pmatrix},\qquad
Y=\begin{pmatrix}0&1\\1&0\end{pmatrix},
\]
证明 \(u_{(a,b)}\mapsto X^aY^b\) 给出 \(\mathbb C^\alpha[Q]\cong M_2(\mathbb C)\)，并求其不可约射影表示的维数。
\end{exercise}
\begin{solution}
函数的指数 \(bc\) 对两个变量双线性，展开三个变量即得余圈等式，单位元处也为 \(1\)。但 \(\alpha\bigl((1,0),(0,1)\bigr)=1\)，反向为 \(-1\)，而交换群的余边对称，所以它不是余边。

\(X^2=Y^2=I_2\) 且 \(YX=-XY\)，故矩阵乘法恰好满足给定的因子集。四个矩阵 \(I_2,X,Y,XY\) 线性无关，分别张成对角和非对角的全部矩阵，因此映射是代数同构。为直接看出简单模的形式，设 \(E\) 是非零简单 \(M_2(\mathbb C)\)-模，选 \(0\neq v\in E_{11}E\)；若原先只有 \(E_{22}E\) 非零，用 \(E_{12}\) 即移到这里。\(v,E_{21}v\) 线性无关，其线性生成空间被所有矩阵单位保持，故由简单性等于 \(E\)。所以只有一个不可约类型，维数为 \(2\)。
\end{solution}

\begin{exercise}\label{b5:ex:04-quaternion-factor-change}
设 \(Q=\mathbb F_2^2\) 为加法群，\(\alpha((a,b),(c,d))=(-1)^{bc}\) 为复标量因子集。比较 \(\alpha\) 与 \(Q_8\) 延拓障碍产生的因子集
\(\beta\bigl((a,b),(c,d)\bigr)=(-1)^{ac+bd+bc}\)，各指数按模 \(2\) 计算。求一个显式函数 \(f:Q\to\mathbb C^\times\)，使 \(\beta=(\delta f)\alpha\)，其中 \((\delta f)(p,q)=f(p)f(q)/f(p+q)\)。解释为何两者描述同一个非零障碍类。
\end{exercise}
\begin{solution}
取 \(f(a,b)=\mathrm i^{a+b}\)，其中 \(a,b\) 用 \(0,1\) 表示。模 \(2\) 加法发生进位时，
\[
\frac{f(a,b)f(c,d)}{f(a+c,b+d)}
=(-1)^{ac+bd}.
\]
分母的参数按模 \(2\) 相加。因此 \(\beta=(\delta f)\alpha\)，它们相差余边，代表同一个类。这个类非零，因为上一题已由不对称性证明 \(\alpha\) 不是余边。相应扭曲群代数通过重标定基同构。
\end{solution}

\begin{exercise}\label{b5:ex:04-factor-determinant}
设有限群 \(Q\) 有一个 \(d\) 维射影复表示 \(R\)，因子集为 \(\alpha\)。证明 \(\alpha^d\) 是余边。如果 \([\alpha]\) 在余圈模余边的群中有有限阶 \(m\)，证明 \(m\mid d\)，并应用于前两题。
\end{exercise}
\begin{solution}
对 \(R(p)R(q)=\alpha(p,q)R(pq)\) 取行列式，得到
\[
\alpha(p,q)^d
=\frac{\det R(p)\det R(q)}{\det R(pq)}.
\]
右侧就是函数 \(q\mapsto\det R(q)\) 的余边，所以 \([\alpha]^d=1\)。阶的定义给出 \(m\mid d\)。前两题中 \(\alpha^2=1\) 且类非零，所以其类的阶为 \(2\)，任何具有该因子集的非零射影表示都具有偶数维数；构造出的二维表示达到最小可能维数。
\end{solution}

\begin{exercise}\label{b5:ex:04-central-extension-factor}
设有限群 \(Q\) 上的归一化余圈 \(\alpha\) 取值于 \(m\) 次单位根群 \(\mu_m\)。在 \(\mu_m\times Q\) 上定义
\[
(z,p)(w,q)=\bigl(zw\alpha(p,q),pq\bigr).
\]
证明它是一个群，且 \(\mu_m\times\{1\}\) 是中心子群。证明因子集为 \(\alpha\) 的射影表示 \(R\) 可以提升为这个群的普通表示。
\end{exercise}
\begin{solution}
三项乘法的标量部分相等，恰好等价于余圈等式；单位元为 \((1,1)\)。把余圈等式用于 \((p,p^{-1},p)\)，得到 \(\alpha(p,p^{-1})=\alpha(p^{-1},p)\)，因此逆元为
\[
(z,p)^{-1}
=\bigl(z^{-1}\cdot\alpha(p,p^{-1})^{-1},p^{-1}\bigr).
\]
归一化条件说明 \((z,1)\) 与所有元素对易；投影到 \(Q\) 的核正是这个中心子群。定义 \(\widetilde R(z,p)=z\cdot R(p)\)，则
\(\widetilde R(z,p)\widetilde R(w,q)=zw\alpha(p,q)R(pq)\)，与群乘法相容。这一表示把中心元素 \((z,1)\) 送到 \(zI\)。
\end{solution}

\begin{exercise}\label{b5:ex:04-twisted-center}
设 \(Q\) 为有限交换群，\(\alpha\) 为归一化复标量 \(2\)-余圈。证明扭曲群代数 \(\mathbb C^\alpha[Q]\) 的中心由满足
\[
\alpha(p,q)=\alpha(q,p)\qquad\hbox{对所有 }q\in Q
\]
的那些 \(u_p\) 张成，并求\cref{b5:ex:04-twisted-klein-algebra}中代数的中心。
\end{exercise}
\begin{solution}
写 \(z=\sum_p a_pu_p\)。对固定 \(q\)，因为 \(Q\) 交换，比较 \(zu_q\) 与 \(u_qz\) 的基向量 \(u_{pq}\) 的系数，得到
\[
a_p\bigl(\alpha(p,q)-\alpha(q,p)\bigr)=0.
\]
所以 \(a_p\neq0\) 的必要充分条件正是题中的全部等式。对 \(Q=\mathbb F_2^2\)、\(\alpha\bigl((a,b),(c,d)\bigr)=(-1)^{bc}\)，分别取 \((c,d)=(1,0),(0,1)\)，条件迫使 \(b=0,a=0\)。故中心只有 \(\mathbb C u_{(0,0)}\)，与 \(M_2(\mathbb C)\) 的标量中心一致。
\end{solution}

\begin{exercise}\label{b5:ex:04-direct-product-clifford}
设 \(N,Q\) 为有限群，\(G=N\times Q\)。利用本章的重数空间构造，证明 \(G\) 的每个不可约复表示都唯一写成 \(W\otimes E\)，其中 \(W,E\) 分别是 \(N,Q\) 的不可约复表示，作用为 \((n,q)(w\otimes e)=nw\otimes qe\)。
\end{exercise}
\begin{solution}
把 \(N\) 识别为第一因子。它的每个不可约类型均被 \(G\) 固定，因为第二因子与它对易，第一因子的共轭只改变同构表示。Clifford 限制定理使任何不可约 \(G\)-表示限制为 \(W^{\oplus e}\)，且 \(W\) 的同构类型唯一。用 \((n,q)\mapsto\rho_W(n)\) 延拓 \(W\)，可在商群 \(Q\) 上取所有 \(T_q=1\)，因子集为 \(1\)。\cref{b5:thm:clifford-projective-multiplicity}遂给出普通 \(Q\)-表示 \(E=\operatorname{Hom}_N(W,V)\)，并说明 \(V\) 不可约当且仅当 \(E\) 不可约。求值同构给出所述张量积，也保证 \(E\) 的类型唯一。反过来，每一对不可约 \(W,E\) 都由该定理给出不可约张量积。
\end{solution}

\begin{exercise}\label{b5:ex:04-clifford-square-sum}
设 \(G\) 为有限群，\(N\trianglelefteq G\)，\(W\) 为不可约复 \(N\)-表示，\(I=I_G(W)\)。让 \(U_1,\ldots,U_r\) 遍历限制到 \(N\) 时含有 \(W\) 的不可约 \(I\)-表示，写 \(U_j|_N\cong W^{\oplus e_j}\)。证明
\[
\sum_{j=1}^r e_j^2=[I:N].
\]
解释它在因子集非平凡时说明了什么。
\end{exercise}
\begin{solution}
把完全可约的 \(\operatorname{Ind}_N^I W\) 分解为不可约表示。由 Frobenius 互反性，\(U_j\) 的重数为
\(\dim\operatorname{Hom}_N(W,U_j|_N)=e_j\)，其他不可约表示的重数为零。因此
\[
[I:N]\dim W
=\dim\operatorname{Ind}_N^I W
=\sum_j e_j\dim U_j
=\sum_j e_j^2\dim W.
\]
消去非零的 \(\dim W\) 得到公式。由重数空间定理，\(e_j\) 是同一个反因子集的全部不可约射影表示的维数，因此即使障碍非零，仍有相应的平方和公式；不能把这些射影表示直接换成普通商群表示。
\end{solution}

\begin{exercise}\label{b5:ex:04-schurian-not-simple}
设 \(k\) 的特征为 \(3\)，让 \(S_3\) 在
\[
A=\bigl\{(x,y,z)\in k^3:x+y+z=0\bigr\}
\]
上置换坐标。证明 \(A\) 可约，但 \(\operatorname{End}_{k[S_3]}A=k\)。由此说明 Mackey 不可约性判据的证明为何使用完全可约性。
\end{exercise}
\begin{solution}
非零向量 \((1,1,1)\) 属于 \(A\)，其生成的直线被全群固定，所以二维空间 \(A\) 可约。取基 \(u=(1,-1,0)\)、\(v=(0,1,-1)\)。在此基下，\((123)\) 与 \((12)\) 的矩阵分别为
\[
C=\begin{pmatrix}0&-1\\1&-1\end{pmatrix},
\qquad
T=\begin{pmatrix}-1&1\\0&1\end{pmatrix}.
\]
若算子 \(F\) 与 \(C\) 对易，写 \(F(u)=au+bv\)，则 \(F(v)=F(Cu)=CF(u)=-bu+(a-b)v\)，故 \(F=aI+bC\)。它还与 \(T\) 对易，当且仅当 \(b(CT-TC)=0\)。矩阵 \(CT-TC\) 非零，例如在 \(u\) 上的值为 \(-u-2v\neq0\)，于是 \(b=0\)。所以所有群模自同态都是标量。自同态空间一维在缺少完全可约性时，并不能排除真子表示。
\end{solution}

\begin{exercise}\label{b5:ex:04-order-twenty-one}
设 \(G=\langle r,s:r^7=s^3=1,\ srs^{-1}=r^2\rangle\)，并令 \(N=\langle r\rangle\)。求 \(N\) 的非平凡一维复特征标在 \(G\) 共轭下的轨道，构造由它们得到的不可约三维表示，并证明这些表示与三个一维表示穷尽 \(G\) 的不可约复表示。
\end{exercise}
\begin{solution}
取本原七次单位根 \(\zeta\)，写 \(\lambda_j(r)=\zeta^j\)。由于 \(s^{-1}rs=r^4\)，共轭把 \(j\) 送到 \(4j\)。六个非零指数分成 \(\{1,2,4\}\) 与 \(\{3,5,6\}\) 两个轨道，各特征标的惯性群都是 \(N\)。Clifford 对应给出两个三维不可约表示 \(\operatorname{Ind}_N^G\lambda_1\)、\(\operatorname{Ind}_N^G\lambda_3\)；它们限制到 \(N\) 的类型不同，故不同构。

若一个不可约 \(G\)-表示限制到 \(N\) 含有平凡特征标，Clifford 限制定理说明整个 \(N\) 都平凡作用，因此该表示来自 \(G/N\cong C_3\)，共有三个一维类型。若限制中含有非平凡特征标，则它属于上述两个轨道之一，且惯性群等于 \(N\)，Clifford 对应说明它只能是相应的三维诱导表示。所有情形因而已经穷尽。
\end{solution}

% !TeX root = ../../Book5.tex
% 纲要标题：## 第5章　Artin 诱导、Brauer 诱导与群论应用
% 纲要位置：工作记录/计划纲要.md，第 3845 行。
% 纲要细目（写作范围）：
% > 本章顺读：表示环与整性 → Artin 有理诱导 → Brauer 整诱导的陈述与证明 → Burnside 可解性定理 → 群论应用。Burnside 证明只使用本章整性工具和此前有限群、特征标理论，不以前面整诱导的长证明为前提。

\chapter{Artin 诱导、Brauer 诱导与群论应用}
\label{b5:ch:05}
\BookChapterNumber{b5:ch:05}

\begin{chapterabstract}
诱导既能构造新的表示，也能将表示环中的元素写成来自较小子群的诱导特征标的线性组合。本章先建立特征标值与中心类和的整性，再证明 Artin 的循环子群有理诱导及 Brauer 的初等子群整诱导。整系数结论通过逐素数的评价检测获得，并进一步化为一维表示的诱导。Burnside 的 \(p^a q^b\) 阶群可解性定理则由类和整性和共轭类大小证明，不依赖前面的整诱导定理。这一应用说明特征标怎样约束有限群结构。
\end{chapterabstract}

\begin{learninggoals}
\begin{itemize}
\item 使用表示环、次数及不可约系数，区分实际特征标与虚特征标。
\item 证明特征标值与类和标量的整性，并说明次数整除的来源。
\item 区分有理系数诱导、有理值特征标及仅从平凡表示诱导的不同结论。
\item 掌握 \(p\)-初等子群的定义，逐步排除整诱导中的每个素数分母。
\item 用 Clifford 理论证明 \(p\)-群的单项式性，将整诱导化为来自一维表示的诱导。
\item 从素数幂大小的共轭类构造正规子群，证明 Burnside 可解性定理，并判断哪些群满足其条件。
\end{itemize}
\end{learninggoals}

本章使用前四章的复特征标、诱导公式、张量恒等式及 Clifford 对应，也使用有限群的 Sylow 定理。所需代数整数与有限交换环的引理在章内证明。\secref{b5:sec:0504}从特征标计算转入表示环的局部化与极大理想检测，给出 Brauer 整诱导的证明；它依赖\secref{b5:sec:0502}的有理诱导和\secref{b5:sec:0503}的准备。\secref{b5:sec:0505}的 Burnside 定理则只用\secref{b5:sec:0501}的整性工具及前四章的群与特征标知识，不依赖 Brauer 整诱导。整诱导的证明也不要求模表示理论。

\ProseParagraphBreak
参考阅读中，Etingof 等的《Introduction to Representation Theory》\cite[Sections 5.2--5.4; Section 5.26]{EtingofEtAl2011RepresentationTheory}分别讨论整性、Burnside 证明和 Artin 诱导；Serre 的《Linear Representations of Finite Groups》\cite[Chapter 10, pp. 74--80]{Serre1977FiniteRepresentations}集中讨论 Brauer 诱导。本章逐素数证明使用表示环的极大理想评价，并把相关环论步骤展开。

% !TeX root = ../../Book5.tex
% 纲要标题：### 5.1 表示环、代数整数性与类和
% 纲要位置：工作记录/计划纲要.md，第 3849 行。
% 纲要细目（写作范围）：
% * 表示环
% * 特征标值的代数整数性
% * 诱导法处理表示环问题
% * 分圆域与特征标的域
% * 从有限阶矩阵的单位根特征值证明特征标值的代数整数性，再建立群代数中心类和及其标量作用所需的整性引理
% * 所需代数整数的和、积与有理代数整数判别在本节就地准备，不把第三卷整扩张理论作为首次学习有限群表示的前提

\section{表示环、代数整数性与类和}
\label{b5:sec:0501}

特征标的加法来自直和，乘法来自张量积。允许相减以后，诱导问题便可以放进一个环中讨论。整数系数的诱导比有理系数的诱导要求更强，需要研究特征标格、单位根和有限生成的整系数模。本节介绍这些工具。

\subsection{表示环}
\label{b5:subsec:0501:01}

\begin{definition}\label{b5:def:complex-representation-ring}
设 \(G\) 为有限群，\(\operatorname{Irr}(G)\) 为其复不可约特征标的集合。所有有限和
\[
R_{\mathbb C}(G)\coloneqq
\bigoplus_{\chi\in\operatorname{Irr}(G)}\mathbb Z\chi
\]
在类函数的逐点加法、乘法下组成的环，称为 \(G\) 的复\term[biaoshihuan]{表示环}[representation ring]，以下简记为 \(R(G)\)。其元素称为\term[xutezhengbiao]{虚特征标}[virtual character]。乘法单位是平凡特征标 \(1_G\)，而不是正则特征标。
\end{definition}
\symidx{\(R(G)\)：有限群的复表示环}

乘积确实仍在上述整数格内：若 \(V,W\) 为复表示，则 \(\chi_{V\otimes W}=\chi_V\chi_W\)，将 \(V\otimes W\) 分解成不可约表示便得到非负整数系数。每个虚特征标都能写成 \(\chi_V-\chi_W\)；此记号只记录直和的形式差，不声称存在一个维数可能为负的向量空间。

\ProseParagraphBreak
次数映射 \(\dim:R(G)\rightarrow\mathbb Z\)、\(\alpha\mapsto\alpha(1)\) 是环同态。其核称为增广理想。次数非负的虚特征标未必是实际特征标：在 \(C_2\) 中，若非平凡一维特征标为 \(\varepsilon\)，则 \(2\cdot1-\varepsilon\) 的次数为 \(1\)，但不可约展开有负系数。由正交关系，\(\alpha\in R(G)\) 是实际特征标，当且仅当对所有不可约 \(\chi\)，\(\langle\chi,\alpha\rangle_G\ge0\)。本章内积对第一变量共轭线性。

\subsection{特征标值的代数整数性}
\label{b5:subsec:0501:02}

\begin{definition}\label{b5:def:algebraic-integer}
设 \(z\in\mathbb C\)。若存在 \(n\ge1\) 及整数 \(a_0,\ldots,a_{n-1}\)，使
\[
z^n+a_{n-1}z^{n-1}+\cdots+a_0=0,
\]
则称 \(z\) 为\term[daishuzhengshu]{代数整数}[algebraic integer]。
\end{definition}

下面的判据避免每遇到一个和或积便重新猜测其首一方程。
\begin{lemma}[有限生成模的整性判据]\label{b5:lem:finite-module-integrality}
设 \(A\subset B\) 为含同一单位的交换环，\(B\) 作为 \(A\)-模由有限多个元素生成。每个 \(b\in B\) 都满足某个首一 \(A\)-系数多项式。特别地，若 \(B\subset\mathbb C\) 且 \(B\) 是有限生成的 \(\mathbb Z\)-模，则 \(B\) 中每个元素都是代数整数。
\end{lemma}
\begin{proof}
选取模生成元 \(u_1,\ldots,u_r\)，写成 \(bu_j=\sum_i a_{ij}u_i\)，并记行向量 \(u=(u_1,\ldots,u_r)\)、矩阵 \(C=(a_{ij})\)。这些关系合在一起就是
\[
u(bI-C)=0.
\]
在交换环 \(B\) 上右乘伴随矩阵，得到
\[
u(bI-C)\operatorname{adj}(bI-C)
=u\det(bI-C)=0.
\]
因此首一多项式 \(F(T)=\det(TI-C)\in A[T]\) 的值 \(F(b)\) 零化每个 \(u_j\)。由于 \(1\) 是这些生成元的 \(A\)-线性组合，\(F(b)\cdot1=0\)，故 \(F(b)=0\)。这里使用的是生成关系与伴随矩阵恒等式，不要求生成元线性无关，也没有把这组生成元当作基。
\end{proof}

\begin{proposition}\label{b5:prop:character-values-integral}
设 \(G\) 为有限群，\(m\) 为各元素阶的最小公倍数，\(\zeta_m=e^{2\pi i/m}\)。每个复虚特征标 \(\alpha\in R(G)\) 的值均属于 \(\mathbb Z[\zeta_m]\)，并且都是代数整数。
\end{proposition}
\begin{proof}
任一表示矩阵 \(\rho(g)\) 满足 \(\rho(g)^m=I\)。由于 \(T^m-1\) 在 \(\mathbb C\) 中无重根，\(\rho(g)\) 可对角化，其特征值都是 \(\zeta_m\) 的幂。取迹及整数线性组合便得 \(\alpha(g)\in\mathbb Z[\zeta_m]\)。环 \(\mathbb Z[\zeta_m]\) 由 \(1,\zeta_m,\ldots,\zeta_m^{m-1}\) 生成，应用\cref{b5:lem:finite-module-integrality}即可。
\end{proof}

代数整数不必是实数，更不必是通常的整数。循环群的一维特征标可以取单位根为值；整性限制的是它满足的多项式，而非其在实轴上的位置。

\subsection{表示环中的诱导法}
\label{b5:subsec:0501:03}

对有限群的子群 \(H\le G\)，限制是环同态 \(\operatorname{Res}_H^G:R(G)\rightarrow R(H)\)，诱导首先是加法群同态 \(\operatorname{Ind}_H^G:R(H)\rightarrow R(G)\)。由\cref{b5:thm:tensor-identity}，它们满足
\begin{equation}\label{b5:eq:induction-ring-projection}
\alpha\operatorname{Ind}_H^G\beta
=\operatorname{Ind}_H^G\left(\operatorname{Res}_H^G\alpha\cdot\beta\right),
\qquad \alpha\in R(G),\quad\beta\in R(H).
\end{equation}
这条张量恒等式说明诱导像虽然未必含单位，却受到整个表示环乘法的控制。

\begin{proposition}\label{b5:prop:induction-ideal}
设 \(G\) 为有限群，\(\mathcal F\) 为它的一族子群，\(R(H)\) 表示有限群 \(H\) 的复表示环。令
\[
I_{\mathcal F}(G)=\sum_{H\in\mathcal F}\operatorname{Ind}_H^G R(H).
\]
则 \(I_{\mathcal F}(G)\) 是 \(R(G)\) 的理想。因此所有虚特征标均属于这个诱导像之和，当且仅当 \(1_G\in I_{\mathcal F}(G)\)。
\end{proposition}
\begin{proof}
定义已经给出加法子群。\eqref{b5:eq:induction-ring-projection}把一个诱导生成元与任意 \(\alpha\in R(G)\) 的乘积仍写成来自同一个 \(H\) 的诱导。若理想含 \(1_G\)，则含每个 \(\alpha=\alpha1_G\)；反向显然。
\end{proof}

诱导一般不保持乘法。若 \(H=1<G\)，则 \(\operatorname{Ind}_1^G1\) 是次数 \(|G|\) 的正则特征标。其平方的次数为 \(|G|^2\)，而 \(\operatorname{Ind}_1^G(1\cdot1)\) 的次数只有 \(|G|\)。

\subsection{分圆域与特征标的域}
\label{b5:subsec:0501:04}

\(\mathbb Q(\zeta_m)\) 称为相应的\term[fenyuanyu]{分圆域}[cyclotomic field]。对复特征标 \(\chi\)，由全部值生成的域 \(\mathbb Q\bigl(\chi(g):g\in G\bigr)\) 称为其特征标值域。\cref{b5:prop:character-values-integral}说明它包含于一个有限分圆扩张；这并不立即给出一组使全部表示矩阵也落在该值域中的基。

\begin{lemma}\label{b5:lem:cyclotomic-power-automorphisms}
设 \(m\ge1\) 为整数，\(\zeta_m=e^{2\pi\mathrm i/m}\)，\(a\) 为与 \(m\) 互素的整数。代换 \(\zeta_m\mapsto\zeta_m^a\) 给出 \(\mathbb Q(\zeta_m)\) 的自同构 \(\sigma_a\)。对有限群 \(G\) 的复虚特征标 \(\alpha\)，若 \(m\) 是 \(G\) 的指数，则对每个 \(g\in G\) 有
\[
\sigma_a\bigl(\alpha(g)\bigr)=\alpha(g^a).
\]
\end{lemma}
\begin{proof}
补充自同构存在性的理由。先说明将用到的有理首一因子判据。有限多个代数整数生成的环是有限生成的 \(\mathbb Z\)-模：依次用各自的首一方程约去高次幂即可。因此由\cref{b5:lem:finite-module-integrality}，这些数的和、积也都是代数整数。若既约分数 \(u/v\)、\(v>0\) 满足首一整方程，清除分母便得 \(v\mid u^n\)，故 \(v=1\)。一个首一整系数多项式的首一有理因子的系数，是原多项式部分根的整系数对称式，既有理又为代数整数，因而仍是整数。

分圆多项式 \(\Phi_m(T)\) 以所有本原 \(m\) 次单位根为根；递推分解 \(T^m-1=\prod_{d\mid m}\Phi_d(T)\) 及首一多项式的整系数除法给出 \(\Phi_m\in\mathbb Z[T]\)。取 \(\zeta_m\) 的首一不可约多项式 \(f\)，由刚证的因子判据有 \(f\in\mathbb Z[T]\)，且 \(f\mid\Phi_m\)。

若素数 \(\ell\nmid m\) 而 \(\zeta_m^\ell\) 不是 \(f\) 的根，令 \(h\ne f\) 为它在 \(\Phi_m\) 中的首一不可约因子。因为 \(h(\zeta_m^\ell)=0\)，有 \(f(T)\mid h(T^\ell)\)。这些多项式首一，整系数除法使这一整除关系可以模 \(\ell\) 约化，得到 \(\overline f\mid\overline h(T)^\ell\)。这里不要求 \(\overline f\) 不可约；它仍是正次数的首一多项式，故可选一个不可约因子 \(d\in\mathbb F_\ell[T]\)。于是 \(d\mid\overline h^\ell\)，不可约多项式在 \(\mathbb F_\ell[T]\) 中是素元，所以 \(d\mid\overline h\)。由于 \(f,h\) 是不同的不可约因子，有 \(fh\mid T^m-1\)，约化后便得
\[
d^2\mid\overline f\,\overline h\mid T^m-1.
\]
但 \(\ell\nmid m\) 时，\(T^m-1\) 与导数 \(mT^{m-1}\) 互素，不可能有这样的重因子。因此 \(\zeta_m^\ell\) 仍是 \(f\) 的根。

若已知本原根 \(\zeta_m^b\) 是 \(f\) 的根，由 \(f\) 首一且不可约，它同样是 \(\zeta_m^b\) 在 \(\mathbb Q\) 上的极小多项式。将上述论证中的 \(\zeta_m\) 换成 \(\zeta_m^b\)，便知对每个素数 \(\ell\nmid m\)，\(\zeta_m^{b\ell}\) 仍是 \(f\) 的根。每个与 \(m\) 互素的正整数都是这些素数的乘积，\(1\) 对应空积，故全部本原 \(m\) 次单位根都是 \(f\) 的根。它们互不相同，所以
\[
\deg f\ge\varphi(m)=\deg\Phi_m.
\]
另一方面，\(f\mid\Phi_m\) 给出反向不等式；两者次数相等且均首一，因而 \(f=\Phi_m\)。

任意与 \(m\) 互素的整数 \(a\) 都模 \(m\) 同余于一个与 \(m\) 互素的正整数，故 \(\zeta_m^a\) 也有极小多项式 \(f\)。因此代换 \(\zeta_m\mapsto\zeta_m^a\) 给出域同态，其逆由 \(a\) 模 \(m\) 的逆元给出。最后，将特征标值写成单位根的整数线性组合，逐项施行 \(\sigma_a\) 即得公式。
\end{proof}

上述论证由代数整数的封闭性得到首一有理因子的整系数性，再由模素数约化证明分圆多项式不可约。特征标值域与表示矩阵的定义域仍须分别考虑。

\subsection{单位根、中心类和与标量作用的整性}
\label{b5:subsec:0501:05}

设 \(K_1,\ldots,K_r\) 为 \(G\) 的共轭类，取 \(g_j\in K_j\)。在群代数中记 \(z_j=\sum_{g\in K_j}g\)。这些\term[leihe]{类和}[class sum]构成 \(Z\bigl(\mathbb C[G]\bigr)\) 的一组基：群代数元素位于中心，恰好意味着其系数在共轭类上不变。并且
\[
z_i z_j=\sum_k a_{ijk}z_k,\qquad a_{ijk}\in\mathbb Z_{\ge0},
\]
其中 \(a_{ijk}\) 等于把 \(K_k\) 中某个固定元素写成 \(xy\)、\(x\in K_i,y\in K_j\) 的方法数；同时共轭给出不同固定元素之间的双射。

\begin{proposition}[中心标量的整性]\label{b5:prop:central-character-integrality}
设 \(G\) 为有限群，\(\chi\) 为次数 \(d\) 的复不可约特征标，\(g\in G\)，\(K\) 为 \(g\) 的共轭类。类和 \(z_K=\sum_{x\in K}x\in\mathbb C[G]\) 在对应不可约表示上作用为标量
\[
\omega_\chi(K)=\frac{|K|\cdot\chi(g)}{d},
\]
且此标量是代数整数。
\end{proposition}
\begin{proof}
类和位于中心，由\cref{b5:cor:schur-algebraically-closed}，在复不可约表示上的作用为标量。取迹得到 \(d\cdot\omega_\chi(K)=|K|\cdot\chi(g)\)。整系数类和格 \(\bigoplus_j\mathbb Zz_j\) 是含单位的环，并且是有限生成的 \(\mathbb Z\)-模。\cref{b5:lem:finite-module-integrality}给出类和满足的首一整系数方程；施行表示后，标量 \(\omega_\chi(K)\) 满足同一方程。
\end{proof}

这比特征标值本身的整性更强，因为式中除以了次数 \(d\)。Burnside 证明将把这个分母与共轭类大小的互素条件结合。

\subsection{代数整数的基本运算与有理数判别}
\label{b5:subsec:0501:06}

\begin{lemma}\label{b5:lem:algebraic-integer-arithmetic}
代数整数的和、差、积仍为代数整数；有理数是代数整数，当且仅当它是整数。若代数整数 \(\alpha\) 的首一极小多项式为 \(f\in\mathbb Q[T]\)，则 \(f\in\mathbb Z[T]\)，因而全部共轭根的乘积是整数。
\end{lemma}
\begin{proof}
若 \(\alpha,\beta\) 分别满足次数 \(r,s\) 的首一整系数方程，则环 \(\mathbb Z[\alpha,\beta]\) 由 \(\alpha^i\beta^j\)、\(0\le i<r,0\le j<s\) 有限生成。\cref{b5:lem:finite-module-integrality}同时处理和、差、积。若既约分数 \(u/v\)、\(v>0\) 满足次数 \(n\) 的首一整系数方程，乘以 \(v^n\) 得 \(v\mid u^n\)，所以 \(v=1\)。

最后，\(f\) 的每个根也是 \(\alpha\) 所满足的那个首一整系数多项式的根，因而都是代数整数。\(f\) 的系数是这些根的整数多项式，又属于 \(\mathbb Q\)，由刚证的两项结论可知均为整数。根的乘积就是带符号的常数项。
\end{proof}

\begin{corollary}[Frobenius 次数整除定理]\label{b5:cor:character-degree-divisibility}
有限群 \(G\) 的每个复不可约特征标的次数 \(d\) 都整除 \(|G|\)。
\end{corollary}
\begin{proof}
设该特征标为 \(\chi\)。由\cref{b5:thm:character-hom-inner-product}与\cref{b5:cor:schur-algebraically-closed}，有 \(\langle\chi,\chi\rangle_G=1\)。按共轭类求和得
\[
\frac{|G|}{d}
=\sum_j\frac{|K_j|\cdot\chi(g_j)}{d}\,\overline{\chi(g_j)}.
\]
右侧各因子由\cref{b5:prop:central-character-integrality}及\cref{b5:prop:character-values-integral}均为代数整数。因此 \(|G|/d\) 是有理代数整数，由\cref{b5:lem:algebraic-integer-arithmetic}属于 \(\mathbb Z\)。
\end{proof}

整性与次数整除可参阅 Etingof 等\cite[Sections 5.2--5.3, pp. 93--98]{EtingofEtAl2011RepresentationTheory}。后文只需上述有限生成模的整性判据、类和与有理数判别，不预设一般整扩张理论。

% !TeX root = ../../Book5.tex
% 纲要标题：### 5.2 Artin 有理诱导
% 纲要位置：工作记录/计划纲要.md，第 3858 行。
% 纲要细目（写作范围）：
% * Artin 诱导定理
% * 有理系数与整数系数的区别
% * 明确所取 Artin 诱导版本：\(R_{\mathbb C}(G)\otimes_{\mathbb Z}\mathbb Q\) 由循环子群表示的诱导张成；有理系数的线性组合不意味着所有特征标值都在 \(\mathbb Q\) 中
% * 把从循环子群一般特征标诱导的表述与仅从平凡特征标诱导的有理值版本分开；用 Frobenius 互反和类函数检测组织证明

\section{Artin 有理诱导}
\label{b5:sec:0502}

每个群元素都在一个循环子群里。单靠这句话无法拼合实际表示，却足以检测类函数是否为零。借助 Frobenius 互反，把这种检测转置过来，便得到诱导的有理张成结论。

\subsection{Artin 诱导定理}
\label{b5:subsec:0502:01}

\begin{theorem}[Artin 有理诱导]\label{b5:thm:artin-induction}
设 \(G\) 为有限群。每个复虚特征标 \(\alpha\in R(G)\) 都能写成
\[
\alpha=\sum_{j=1}^s a_j\operatorname{Ind}_{C_j}^G\lambda_j,
\qquad a_j\in\mathbb Q,
\]
其中 \(C_j\le G\) 为循环子群，\(\lambda_j\) 为 \(C_j\) 的一维复特征标。等价地，这些诱导特征标张成 \(\mathbb Q\otimes_{\mathbb Z}R(G)\)。
\end{theorem}
\term[Artinyoudao]{Artin 诱导定理}[Artin's induction theorem]在这里是关于系数的结论；它没有断言 \(\alpha(g)\) 为有理数。

\begin{proof}
先在复类函数空间中证明张成。若类函数 \(f\) 与每个 \(\operatorname{Ind}_C^G\lambda\) 正交，则由\cref{b5:cor:frobenius-character}，\(\operatorname{Res}_C^Gf\) 与循环群 \(C\) 的每个不可约特征标正交。互反等式也适用于任意类函数：\cref{b5:thm:irreducible-character-basis}把类函数写成不可约特征标的复线性组合，等式可按半线性延拓。因此 \(f|_C=0\)。对每个 \(g\in G\) 取 \(C=\langle g\rangle\)，便得 \(f(g)=0\)。有限维内积空间中正交补为零，所以这些诱导特征标复线性张成全部类函数。

取不可约特征标为基，把全部循环子群的一维诱导特征标写为列向量，所得矩阵的元素是非负整数。刚证的结论说该矩阵满行秩，因而存在一个行数阶的非零整数子式。用对应方阵求解任一不可约特征标的坐标，逆矩阵的元素均为有理数。这就将复张成加强为有理张成，再对整数线性组合使用加法即可。
\end{proof}

最后一步必须使用整数矩阵；仅证明复线性张成，尚未自动说明一个指定向量的坐标是有理数。Etingof 等\cite[Theorem 5.26.1, Remark 5.26.2 and Corollary 5.26.3, pp. 141--142]{EtingofEtAl2011RepresentationTheory}给出了相应的子群检测表述。

\subsection{有理系数与整数系数}
\label{b5:subsec:0502:02}

\begin{corollary}\label{b5:cor:artin-finite-index}
设 \(G\) 为有限群，\(R(G)\) 为其复表示环，令 \(I_{\mathrm{cyc}}(G)=\sum_{C\le G,\ C\text{ 循环}}\operatorname{Ind}_C^G R(C)\)。加法商群
\[
R(G)/I_{\mathrm{cyc}}(G)
\]
有限；特别地，存在同一个正整数 \(N\)，使全部 \(\alpha\in R(G)\) 均满足 \(N\alpha\in I_{\mathrm{cyc}}(G)\)。
\end{corollary}
\begin{proof}
由\cref{b5:thm:artin-induction}，对有限多个不可约特征标分别清除有理系数的分母，再取这些分母的公倍数 \(N\)。于是 \(NR(G)\subset I_{\mathrm{cyc}}(G)\)，所述商群是有限群 \(R(G)/NR(G)\) 的商。
\end{proof}

分母确实可能无法只靠循环子群去掉。取四元交换群 \(G=C_2\times C_2\)。每个循环子群的阶为 \(1\) 或 \(2\)，从它诱导的一维特征标的次数为 \(4\) 或 \(2\)。因此循环诱导的任意整数线性组合都有偶数次数，不可能等于 \(1_G\)。若 \(H_1,H_2,H_3\) 是三个二阶子群，则
\begin{equation}\label{b5:eq:klein-artin-denominator}
1_G=\frac12\sum_{i=1}^3\operatorname{Ind}_{H_i}^G1_{H_i}
-\frac12\operatorname{Ind}_1^G1.
\end{equation}
在单位元处右侧等于 \((6-4)/2=1\)；每个非单位元恰属一个 \(H_i\)，相应置换特征标取值 \(2\)，其余项取值 \(0\)。这逐点核验了公式，也显示负系数与分母的不同职责。

\subsection{循环子群诱导与有理张成}
\label{b5:subsec:0502:03}

\begin{proposition}\label{b5:prop:subgroup-cover-induction}
设 \(G\) 为有限群，\(\mathcal F\) 为对 \(G\)-共轭封闭的一族子群，\(R(G)\) 为复表示环，\(\operatorname{Irr}(H)\) 为 \(H\) 的不可约复特征标集合。下列条件等价：每个 \(g\in G\) 属于某个 \(H\in\mathcal F\)；所有 \(\operatorname{Ind}_H^G\theta\)、\(H\in\mathcal F\)、\(\theta\in\operatorname{Irr}(H)\) 有理张成 \(\mathbb Q\otimes_{\mathbb Z}R(G)\)。
\end{proposition}
\begin{proof}
若子群覆盖 \(G\)，则与全部诱导正交的类函数在每个 \(H\) 上为零，因而恒为零。\cref{b5:thm:artin-induction}证明中的整数矩阵秩论证给出有理张成。

反之，若 \(g\) 不属于任何 \(H\in\mathcal F\)，共轭封闭性说明 \(g\) 也不能共轭到这些子群中。\cref{b5:thm:induced-character}的求和集合为空，所以全部诱导特征标在 \(g\) 处均为零。其有理线性组合便不可能得到在 \(g\) 处取值 \(1\) 的平凡特征标。
\end{proof}

循环子群不需要互不相交，也不需要每个元素只被覆盖一次。这里使用的是限制能检测零函数；重叠部分不会破坏正交补论证。但这种检测只给出有理系数结果，整系数结论还需要下一节的整性论证。

\subsection{一般特征标诱导与有理值版本}
\label{b5:subsec:0502:04}

若要求仅从 \(1_C\) 诱导，得到的都是置换特征标，因而必为整数值。它们的有理线性组合也只能是有理值。对有限群 \(G\) 中的元素 \(u,v\)，若存在与 \(|u|\) 互素的整数 \(a\)，使 \(v\) 共轭于 \(u^a\)，则称 \(u,v\) 有理共轭；这等价于它们生成的循环子群共轭。

\begin{proposition}[有理值版本]\label{b5:prop:rational-valued-artin}
设 \(G\) 为有限群。所有在 \(G\) 的有理共轭类上常值的函数 \(f:G\to\mathbb Q\)，由 \(\operatorname{Ind}_C^G1_C\) 在 \(\mathbb Q\) 上张成，其中 \(C\le G\) 遍历循环子群，\(1_C\) 为平凡复特征标。特别地，每个取值于 \(\mathbb Q\) 的复虚特征标都有这种表达。
\end{proposition}
\begin{proof}
从循环子群共轭类中选代表 \(C_1,\ldots,C_s\)，各取生成元 \(c_i\)，按 \(|C_i|\) 从大到小排列。置换特征标 \(\pi_j=\operatorname{Ind}_{C_j}^G1\) 在 \(c_i\) 处非零，恰好在某个共轭的 \(C_j\) 包含 \(C_i\) 时发生。因此矩阵 \(\bigl(\pi_j(c_i)\bigr)\) 为三角矩阵：同阶的不同代表之间该值为零，对角元是 \(\bigl[N_G(C_i):C_i\bigr]>0\)。矩阵为整系数且可逆，故每一组有理数值都能由这些 \(\pi_j\) 的有理线性组合得到。一个有理共轭类恰对应一个循环子群共轭类，所以这已覆盖所述函数空间。

若虚特征标 \(\alpha\) 为有理值，由\cref{b5:lem:cyclotomic-power-automorphisms}得 \(\alpha(g^a)=\alpha(g)\)，其中 \(a\) 与群指数互素。若只给定 \(a\) 与 \(|g|\) 互素，可在同一模 \(|g|\) 剩余类中选择与群指数互素的整数：对未整除 \(|g|\) 的素因子另规定非零剩余类，再用中国剩余定理。于是 \(\alpha\) 在有理共轭类上常值，应用前一结论即可。
\end{proof}

例如 \(C_3=\langle c\rangle\) 的特征标 \(\lambda(c)=e^{2\pi i/3}\) 不能由平凡特征标的循环诱导有理张成，因为它不是有理值；但它就是 \(\operatorname{Ind}_{C_3}^{C_3}\lambda\)，完全符合\cref{b5:thm:artin-induction}。两种表述的区别在于允许使用哪些子群特征标，不能把“有理系数”改说成“有理表示”。

% !TeX root = ../../Book5.tex
% 纲要标题：### 5.3 Brauer 诱导
% 纲要位置：工作记录/计划纲要.md，第 3865 行。
% 纲要细目（写作范围）：
% * 初等子群及一维特征标
% * Brauer 诱导定理的明确版本
% * 与原纲要未指明版本的“Brauer 定理”对接
% * 定义 \(p\)-初等群为 \(P\times C\)，其中 \(P\) 为有限 \(p\)-群、\(C\) 为阶与 \(p\) 互素的循环群；所有素数一起使用
% * 目标版本为复虚特征标是从初等子群诱导的特征标的整数线性组合，并细化到适当子群的一维特征标；完整整性证明在第5.4节
% * 本章不以前置第8章的模表示、Brauer 特征标或分解数为条件

\section{Brauer 诱导}
\label{b5:sec:0503}

四元交换群的例子说明，循环子群有时太小，不能整系数地生成表示环。需要加入的不是任意复杂的子群，而是把一个素数幂阶部分与一个互素的循环部分并列起来。本节确定目标及其准确范围，证明集中在下一节。

\subsection{初等子群与一维特征标}
\label{b5:subsec:0503:01}

\begin{definition}\label{b5:def:elementary-group}
设 \(p\) 为素数，\(E\) 为有限群。若 \(E\) 同构于 \(P\times C\)，其中 \(P\) 是有限 \(p\)-群，\(C\) 是阶与 \(p\) 互素的循环群，则称 \(E\) 为\term[pchudengqun]{\(p\)-初等群}[\(p\)-elementary group]。若存在某个素数 \(p\) 使此条件成立，则称 \(E\) 为\term[chudengqun]{初等群}[elementary group]。有限群 \(G\) 中具有上述结构的子群称为相应的初等子群。
\end{definition}

允许 \(P=1\) 或 \(C=1\)。因此全部有限 \(p\)-群都是 \(p\)-初等群；每个循环群对每一个素数 \(p\) 都是 \(p\)-初等群，只需把其阶分成 \(p\) 部分与互素部分。初等群不一定交换，例如 \(Q_8\times C_3\) 是 \(2\)-初等群。

\ProseParagraphBreak
从一维特征标 \(\lambda:H\rightarrow\mathbb C^\times\) 诱导得到的特征标称为\term[danxiangshitezhengbiao]{单项式特征标}[monomial character]。取左陪集代表后，诱导表示的每个群元素把每个基向量送到另一基向量的非零标量倍，所以矩阵每行、每列恰有一个非零元。名称来自这个具体矩阵形状，并不意味着诱导结果总不可约。

\subsection{Brauer 诱导定理的整系数版本}
\label{b5:subsec:0503:02}

\begin{theorem}[Brauer 整诱导]\label{b5:thm:brauer-induction}
设 \(G\) 为有限群。每个复虚特征标 \(\alpha\in R(G)\) 都是从初等子群的复特征标诱导而得的特征标的整数线性组合。进一步，存在初等子群 \(H_1,\ldots,H_t\le G\)、它们的一维复特征标 \(\lambda_i\) 及整数 \(n_i\)，使
\[
\alpha=\sum_{i=1}^t n_i\operatorname{Ind}_{H_i}^G\lambda_i.
\]
\end{theorem}
这就是本章所用的\term[Braueryoudao]{Brauer 诱导定理}[Brauer's induction theorem]。整系数结论将在\cref{b5:prop:brauer-integral-image}证明，再由\cref{b5:prop:elementary-monomial}化为一维特征标的诱导。

\ProseParagraphBreak
这是表示环中的加法恒等式。若 \(\alpha\) 是实际表示的特征标，右侧仍可能有负系数；把正负部分移到两边才得到实际表示的同构。定理也未声称每个不可约表示单独就是一个一维表示的诱导。

\subsection{Brauer 诱导定理的版本辨析}
\label{b5:subsec:0503:03}

Artin 定理允许循环子群、一般一维特征标和有理系数；Brauer 定理扩大子群族，以此把系数收紧为整数。如果只保留平凡特征标，则所有诱导都是整数值函数，不能产生 \(C_3\) 的非实一维特征标。故“只从平凡表示诱导”不能作为这里的 Brauer 定理。

\ProseParagraphBreak
初等群 \(P\times C\) 与初等交换 \(p\)-群 \(C_p^r\) 是不同概念：前者允许非交换 \(P\) 和互素循环因子。把直积放宽为某些半直积，会得到别的诱导定理中使用的子群族；本章始终使用上述直积定义。Serre 的书中将相应整诱导集中放在 Chapter 10，可作为进一步阅读\cite[Chapter 10, pp. 74--80]{Serre1977FiniteRepresentations}。

\subsection{\texorpdfstring{\(p\)-初等群的定义}{p-初等群的定义}}
\label{b5:subsec:0503:04}

定义中的直积具有可操作的内部判据：\(E\) 中存在子群 \(P,C\)，满足 \(E=PC\)、\(P\cap C=1\)、每个 \(P\) 元素与每个 \(C\) 元素交换，且阶与循环性满足\cref{b5:def:elementary-group}。此时 \((x,c)\mapsto xc\) 是群同构。

\begin{proposition}\label{b5:prop:elementary-subgroups}
设 \(p\) 为素数，\(P\) 为有限 \(p\)-群，\(C\) 为阶与 \(p\) 互素的有限循环群，\(E=P\times C\)。每个子群 \(H\le E\) 都是 \(p\)-初等群，具体地
\[
H=(H\cap P)\times(H\cap C),
\]
其中把 \(P,C\) 视作 \(E\) 的两个自然子群。
\end{proposition}
\begin{proof}
对 \(h=xc\in H\)，\(x\in P,c\in C\)，元素 \(x,c\) 的阶互素。中国剩余定理给出整数 \(a\)，使 \(a\equiv1\pmod{|x|}\)、\(a\equiv0\pmod{|c|}\)，所以 \(h^a=x\in H\)，继而 \(c=x^{-1}h\in H\)。因此 \(H=(H\cap P)(H\cap C)\)。两因子交换且交为 \(1\)，第二个因子是循环群的子群，阶仍与 \(p\) 互素。
\end{proof}

例如 \(S_3\) 不是初等群：其中心为 \(1\)，且它的阶不是素数幂；若它是 \(P\times C\) 且 \(C\ne1\)，则 \(C\) 必在中心中，矛盾。不过其二阶、三阶子群都是初等子群，已经足以给出整诱导。

\subsection{从初等子群到一维特征标的诱导}
\label{b5:subsec:0503:05}

设 \(E=P\times C\)。对复不可约 \(E\)-表示 \(V\)，由\cref{b5:cor:schur-algebraically-closed}，中心中的 \(C\) 作用为一个一维特征标 \(\lambda\)。任何 \(P\)-不变子空间自动也被这个标量作用保持，所以 \(V|_P\) 不可约。故 \(V\) 的特征标为 \((x,c)\mapsto\theta(x)\lambda(c)\)，其中 \(\theta\in\operatorname{Irr}(P)\)。

\ProseParagraphBreak
剩下的任务是证明 \(p\)-群的不可约特征标可从较小子群的一维特征标诱导。一旦写出 \(\theta=\operatorname{Ind}_Q^P\mu\)，就有
\[
\theta\otimes\lambda
=\operatorname{Ind}_{Q\times C}^{P\times C}(\mu\otimes\lambda).
\]
右侧子群仍是 \(p\)-初等群。这一步所需的 \(p\)-群结论将在\cref{b5:thm:p-group-monomial}证明；再由诱导传递性与从 \(E\) 到 \(G\) 的诱导复合。

\subsection{Brauer 诱导与模表示的关系}
\label{b5:subsec:0503:06}

本章始终讨论复表示。下一节会把特征标值模某个素理想约化到特征 \(p\) 的域中，但这只是对一个交换数环取商，并不把表示空间或群代数改为特征 \(p\) 上的对象。模表示在 \(p\,\bigm\vert\,|G|\) 时通常不完全可约。

\ProseParagraphBreak
第八章的 Brauer 特征标只在 \(p\)-正则元素上定义，并处理特征 \(p\) 的表示；这里的 Brauer 诱导则是全部群元素上的普通复特征标恒等式。名字相同，所研究的对象却不同。

% !TeX root = ../../Book5.tex
% 纲要标题：### 5.4 初等子群检测与 Brauer 整诱导的证明
% 纲要位置：工作记录/计划纲要.md，第 3874 行。
% 纲要细目（写作范围）：
% * 把诱导像组织为复表示环中的加法子群，利用限制、诱导与张量恒等式控制其环论行为
% * 建立初等子群上的限制检测与整性引理，分素数处理局部化后的整数系数问题，再由有限生成交换群的局部检测回到整数系数
% * 分别完成“存在某个整数倍的诱导表达”与“无需整体分母的整诱导表达”，指出第二步是超出 Artin 有理诱导的实质
% * 解释初等群上的表示如何进一步化为从适当子群一维表示的诱导；用完整权空间建立惯性群表示与诱导，不把权空间误当成一维；逐项说明所用子群仍处于允许的初等范围
% * 用小群的虚特征标表达说明整数系数可以为负，不把虚表示恒等式误解为实际表示的直接和分解

\section{初等子群检测与 Brauer 整诱导的证明}
\label{b5:sec:0504}

Artin 定理已经给出有限指数的子格。现在必须证明这个指数没有任何素因子。对每个素数 \(p\)，我们反演所有与 \(p\) 互素的整数，证明 \(p\)-初等子群的诱导在这个局部化以后生成整个表示环，从而排除指数中的素因子 \(p\)。论证中所需的交换环知识将在使用处给出证明。

\ProseParagraphBreak
张量恒等式说明诱导像是理想。为了证明它在 \(p\)-局部化后等于全环，我们把局部化表示环 \(R(G)_{(p)}\) 的极大理想写成模 \(p\) 的特征标值检测，并将检测元素约化为 \(p\)-正则部分。在相应中心化子中构造 \(p\)-初等子群，就能找到不落入该极大理想的诱导特征标。对每个素数重复这一论证，才得到整系数等式。

\subsection{诱导像与复表示环}
\label{b5:subsec:0504:01}

固定有限群 \(G\) 与素数 \(p\)，记 \(\mathcal E_p(G)\) 为全部 \(p\)-初等子群，并令
\[
I_p(G)=\sum_{E\in\mathcal E_p(G)}\operatorname{Ind}_E^G R(E).
\]
由\cref{b5:prop:induction-ideal}，这是 \(R(G)\) 的理想。每个循环子群都是 \(p\)-初等群，因此\cref{b5:cor:artin-finite-index}给出
\begin{equation}\label{b5:eq:artin-elementary-finite-index}
I_{\mathrm{cyc}}(G)\subset I_p(G)\subset R(G),
\qquad \bigl[R(G):I_p(G)\bigr]<\infty.
\end{equation}
还需排除这个有限指数中的素因子 \(p\)。

\ProseParagraphBreak
记 \(\mathbb Z_{(p)}=\{a/b\in\mathbb Q:p\nmid b\}\)。对加法群 \(M\)，记 \(M_{(p)}=\mathbb Z_{(p)}\otimes_{\mathbb Z}M\)，可将元素写成 \(x/s\)，其中 \(p\nmid s\)；判据 \(x/1=0\) 等价于存在这样的 \(s\) 使 \(sx=0\)。我们将证明
\[
I_p(G)_{(p)}=R(G)_{(p)}.
\]
这并不要求单独固定一个 \(p\) 就能在整数上生成 \(R(G)\)；它只说明最终的障碍不含 \(p\) 部分。

\subsection{限制检测、局部化与整性引理}
\label{b5:subsec:0504:02}

需要先把 \(p\)-局部化表示环 \(R(G)_{(p)}\) 的极大理想解释为模 \(p\) 的特征标值检测。为此，先证明有限交换环扩张中的极大理想提升引理。

\begin{lemma}[有限扩张中的极大理想提升]\label{b5:lem:finite-extension-maximal-lift}
设 \(A\subset B\) 为含同一单位的交换环，且 \(B\) 为有限生成 \(A\)-模。对每个极大理想 \(\mathfrak m\subset A\)，存在极大理想 \(\mathfrak n\subset B\) 使 \(\mathfrak n\cap A=\mathfrak m\)。
\end{lemma}
\begin{proof}
令 \(S=A\setminus\mathfrak m\)。分式环 \(A_{\mathfrak m}=S^{-1}A\) 的非单位恰组成理想 \(\mathfrak mA_{\mathfrak m}\)，且 \(A_{\mathfrak m}\subset S^{-1}B\) 仍为单射。因此 \(S^{-1}B\ne0\)，并为有限生成 \(A_{\mathfrak m}\)-模。

若 \(\mathfrak mB=B\)，则有限生成模 \(M=S^{-1}B\) 满足 \(\mathfrak mA_{\mathfrak m}M=M\)。取模生成元 \(u_1,\ldots,u_r\)，写成 \(u_j=\sum_i c_{ij}u_i\)，其中全部 \(c_{ij}\in\mathfrak mA_{\mathfrak m}\)。记 \(u=(u_1,\ldots,u_r)\)、\(C=(c_{ij})\)，这些关系给出 \(u(I-C)=0\)。右乘伴随矩阵便得
\[
u(I-C)\operatorname{adj}(I-C)=u\det(I-C)=0.
\]
因 \(u_j\) 生成 \(M\)，这个行列式零化整个 \(M\)，无须假定生成元线性无关。它模 \(\mathfrak mA_{\mathfrak m}\) 为 \(1\)，故为单位，迫使 \(M=0\)，矛盾。因此 \(\mathfrak mB\ne B\)。取包含它的极大理想 \(\mathfrak n\)。其收缩是包含 \(\mathfrak m\) 的真理想，只能等于 \(\mathfrak m\)。
\end{proof}

存在极大理想只需通常的链并与 Zorn 引理\footnote{佐恩（Max August Zorn，1906-1993），德国数学家，提出与选择公理等价的极大原理，并研究代数结构。}：真理想的链之并仍不含 \(1\)。在本节中 \(B/\mathfrak mB\) 还是域 \(A/\mathfrak m\) 上的有限维非零代数，也可直接从真理想的维数中取最大者。

\begin{lemma}[极大理想的求值描述]\label{b5:lem:representation-maximal-evaluation}
设 \(G\) 为有限群，\(m\) 为其指数，\(\zeta_m=e^{2\pi\mathrm i/m}\)，\(O=\mathbb Z[\zeta_m]\)，\(p\) 为素数，\(R(G)\) 为复表示环。下标 \((p)\) 表示反演所有与 \(p\) 互素整数的局部化。对 \(R(G)_{(p)}\) 的每个极大理想 \(\mathfrak m\)，存在 \(g\in G\) 及 \(O_{(p)}\) 的极大理想 \(\mathfrak q\)，使剩余域 \(O_{(p)}/\mathfrak q\) 的特征为 \(p\)，且
\[
\mathfrak m=\bigl\{\,\alpha\in R(G)_{(p)}:\alpha(g)\in\mathfrak q\,\bigr\}.
\]
\end{lemma}
\begin{proof}
设共轭类数为 \(r\)。在每个共轭类代表处取值给出单射环同态
\[
R(G)_{(p)}\longrightarrow B=O_{(p)}^r.
\]
单射性来自特征标作为类函数的相等判据；由\cref{b5:prop:character-values-integral}，各分量的特征标值都属于 \(O_{(p)}\)，故这确实是到 \(B\) 的求值映射。\(O\) 为有限生成 \(\mathbb Z\)-模，所以 \(B\) 为有限生成 \(\mathbb Z_{(p)}\)-模。在这个嵌入下，\(\mathbb Z_{(p)}\) 通过 \(a\mapsto a\,1_G\) 嵌入 \(R(G)_{(p)}\)，其在 \(B\) 中的像是对角标量 \((a,\ldots,a)\)。因此一组 \(\mathbb Z_{(p)}\)-模生成元也生成 \(B\) 作为 \(R(G)_{(p)}\)-模。\cref{b5:lem:finite-extension-maximal-lift}遂把 \(\mathfrak m\) 提升为 \(B\) 的极大理想。

有限直积的极大理想只在一个分量上取极大理想：各分量的单位幂等元在域商中只能取 \(0\) 或 \(1\)，它们的和是 \(1\)，且彼此乘积为零，因此恰有一个取 \(1\)。这便给出某个代表 \(g\) 及 \(O_{(p)}\) 的极大理想 \(\mathfrak q\)。

最后，域 \(O_{(p)}/\mathfrak q\) 对 \(\mathbb Z_{(p)}\) 的像是整的，整性由\cref{b5:lem:finite-module-integrality}给出。若域 \(L\) 整于子整环 \(D\)，则 \(D\) 本身是域：对 \(0\ne a\in D\)，把 \(a^{-1}\in L\) 的首一整方程乘以适当的 \(a\) 的幂，就得到 \(a^{-1}\in D\)。因此 \(\mathfrak q\cap\mathbb Z_{(p)}\) 为极大理想，只能是 \(p\mathbb Z_{(p)}\)。剩余域的特征遂为 \(p\)。
\end{proof}

群元素 \(g\) 有唯一分解 \(g=su=us\)，其中 \(s,u\) 均为 \(g\) 的幂，\(|s|\) 与 \(p\) 互素，\(|u|\) 为 \(p\) 的幂。这由循环群 \(\langle g\rangle\) 的阶的互素分解得到。这样的 \(s\) 称为 \(g\) 的 \(p\)-正则部分；阶与 \(p\) 互素的元素称为\term[pzhengzeyuansu]{\(p\)-正则元素}[\(p\)-regular element]。

\ProseParagraphBreak
表示环的求值像在整数上通常不是任意坐标的直积。以 \(C_2\) 为例，若非平凡一维特征标记为 \(t\)，则 \(t^2=1\)，而 \(1,t\) 是表示环的整数基，所以
\[
R(C_2)=\mathbb Z[t]/(t^2-1).
\]
在单位元与非单位元处取值，把 \(a+bt\) 送到 \((a+b,a-b)\)。像恰为同奇偶的整数对，因为反解是 \(a=(u+v)/2\)、\(b=(u-v)/2\)。局部化到奇素数 \(p\) 时，\(2\) 可逆，这个求值映射成为
\[
R(C_2)_{(p)}\cong\mathbb Z_{(p)}\times\mathbb Z_{(p)},
\]
两个极大理想分别由 \(t=1\) 与 \(t=-1\) 的模 \(p\) 求值给出。局部化到 \(p=2\) 时，模 \(2\) 后却有 \(t^2-1=(t-1)^2\)，两个求值都变成 \(t\mapsto1\)，得到同一个极大理想 \((2,t-1)\)。由刚证的求值描述，它也是此时唯一的极大理想。于是原本不同的两个共轭类在模 \(2\) 的检测中合并，不能任意指定它们的两个剩余值。

\begin{lemma}\label{b5:lem:character-p-regular-reduction}
设 \(G\) 为有限群，\(p\) 为素数，\(m\) 为 \(G\) 的指数，\(\zeta_m=e^{2\pi\mathrm i/m}\)，\(O=\mathbb Z[\zeta_m]\)，\(R(G)\) 为复表示环，下标 \((p)\) 表示反演所有与 \(p\) 互素整数的局部化。设 \(\mathfrak q\subset O_{(p)}\) 为剩余特征是 \(p\) 的极大理想。若 \(g\in G\) 有分解 \(g=su=us\)，其中 \(s,u\) 均为 \(g\) 的幂、\(|s|\) 与 \(p\) 互素、\(|u|\) 为 \(p\) 的幂，则每个 \(\alpha\in R(G)_{(p)}\) 均满足
\[
\alpha(g)\equiv\alpha(s)\pmod{\mathfrak q}.
\]
\end{lemma}
\begin{proof}
对实际表示，交换且可对角化的矩阵 \(\rho(s),\rho(u)\) 可同时对角化：先按 \(\rho(s)\) 的特征空间分解，再在每个被 \(\rho(u)\) 保持的空间内对角化。共同特征向量上的特征值分别记作 \(\xi,\eta\)，其中 \(\eta^{p^a}=1\)。在特征 \(p\) 的剩余域中，\(T^{p^a}-1=(T-1)^{p^a}\)，故 \(\eta\equiv1\)。于是 \(\xi\eta\equiv\xi\)，求迹即得结论。再取整数线性组合并除以与 \(p\) 互素的分母，便覆盖全部 \(\alpha\)。
\end{proof}

\subsection{从有理诱导到整系数诱导}
\label{b5:subsec:0504:03}

\begin{theorem}[逐素数的整诱导]\label{b5:thm:brauer-local-induction}
设 \(G\) 为有限群，\(p\) 为素数，\(R(G)\) 为复表示环，令 \(I_p(G)=\sum_{E\le G,\ E\text{ 为 }p\text{-初等}}\operatorname{Ind}_E^G R(E)\)。下标 \((p)\) 表示反演所有与 \(p\) 互素整数的局部化，则
\[
I_p(G)_{(p)}=R(G)_{(p)}.
\]
因此有限指数 \(\bigl[R(G):I_p(G)\bigr]\) 与 \(p\) 互素。
\end{theorem}
\begin{proof}
若理想 \(I_p(G)_{(p)}\) 为真理想，则它包含在某个极大理想 \(\mathfrak m\) 中。由\cref{b5:lem:representation-maximal-evaluation}及\cref{b5:lem:character-p-regular-reduction}，存在 \(p\)-正则元素 \(s\) 及剩余特征为 \(p\) 的 \(\mathfrak q\subset O_{(p)}\)，使
\[
\alpha\in\mathfrak m\quad\Longleftrightarrow\quad
\alpha(s)\equiv0\pmod{\mathfrak q}.
\]
我们构造一个 \(p\)-初等诱导，使其在 \(s\) 处不满足此同余。

令 \(C=\langle s\rangle\)，在中心化子 \(C_G(s)\) 中取 Sylow \(p\)-子群 \(P\)。由于 \(C\) 在该中心化子内居中，且 \(|C|\) 与 \(|P|\) 互素，
\[
E=CP=C\times P
\]
为 \(p\)-初等子群；整数 \(d=\bigl[C_G(s):E\bigr]\) 与 \(p\) 互素。对 \(C\) 的一维特征标 \(\lambda\)，将其延拓为 \(E\) 上在 \(P\) 处平凡的特征标 \(\widetilde\lambda\)。诱导公式\cref{b5:thm:induced-character}给出
\begin{equation}\label{b5:eq:elementary-induction-evaluation}
\left(\operatorname{Ind}_E^G\widetilde\lambda\right)(s)
=d\sum_{t\in s^G\cap C}\lambda(t).
\end{equation}
详细地说，\(E\) 中与 \(s\) 共轭的元素阶与 \(p\) 互素，故其 \(P\) 分量为 \(1\)，恰落在 \(C\) 内。对每个固定的 \(t\in s^G\cap C\)，满足 \(x^{-1}sx=t\) 的 \(x\in G\) 有 \(|C_G(s)|\) 个，除以 \(|E|\) 即得上式。

集合 \(T=s^G\cap C\) 非空。若对所有 \(\lambda\in\operatorname{Irr}(C)\)，和 \(S_\lambda=\sum_{t\in T}\lambda(t)\) 都属于 \(\mathfrak q\)，取 \(t_0\in T\)，利用循环群的有限几何级数可得
\[
\sum_{\lambda\in\operatorname{Irr}(C)}
S_\lambda\cdot\lambda(t_0)^{-1}
=\sum_{t\in T}\sum_{\lambda\in\operatorname{Irr}(C)}
\lambda(tt_0^{-1})
=|C|.
\]
左侧模 \(\mathfrak q\) 为零，右侧却因 \(p\mathrel{\vcenter{\hbox{\scalebox{1.25}{\(\nmid\)}}}}|C|\) 非零，矛盾。故至少一个 \(S_\lambda\notin\mathfrak q\)。整数 \(d\) 模 \(\mathfrak q\) 也非零，于是\eqref{b5:eq:elementary-induction-evaluation}不属于 \(\mathfrak q\)。这与该诱导属于 \(I_p(G)_{(p)}\subset\mathfrak m\) 矛盾，故理想等于全环。

由\eqref{b5:eq:artin-elementary-finite-index}，商 \(R(G)/I_p(G)\) 有限。若其阶被 \(p\) 整除，它含有阶为 \(p\) 的元素，而这样的元素不会在仅反演与 \(p\) 互素整数后消失。这与刚证的局部化商为零矛盾，故指数与 \(p\) 互素。
\end{proof}

\begin{proposition}\label{b5:prop:brauer-integral-image}
设 \(G\) 为有限群，\(I(G)\) 为全部初等子群的复表示环诱导像之和。则 \(I(G)=R(G)\)。
\end{proposition}
\begin{proof}
\(I(G)\) 包含循环诱导，由\cref{b5:cor:artin-finite-index}，加法商群 \(A=R(G)/I(G)\) 有限。对每个素数 \(p\)，它又是 \(R(G)/I_p(G)\) 的商，而\cref{b5:thm:brauer-local-induction}说明后者阶与 \(p\) 互素。故 \(|A|\) 与每个素数互素，只能等于 \(1\)，于是 \(A=0\)。
\end{proof}

这里先用 Artin 得到有限指数，再逐素数排除指数，正是从“某个整数倍可诱导”到“本身可整诱导”的差别。

\begin{corollary}[初等子群上的虚特征标判据]\label{b5:cor:brauer-character-criterion}
设 \(G\) 为有限群，\(f:G\rightarrow\mathbb C\) 为类函数。则 \(f\in R(G)\)，当且仅当对每个初等子群 \(E\le G\)，限制 \(f|_E\) 属于 \(R(E)\)。
\end{corollary}
\begin{proof}
正向来自实际表示的限制及可加性。反向中，\cref{b5:prop:brauer-integral-image}把每个不可约 \(\chi\) 写成 \(\sum_j n_j\operatorname{Ind}_{E_j}^G\theta_j\)，其中 \(n_j\in\mathbb Z\)、\(\theta_j\in R(E_j)\)。于是
\[
\langle\chi,f\rangle_G
=\sum_j n_j\langle\theta_j,f|_{E_j}\rangle_{E_j}\in\mathbb Z.
\]
这里使用 Frobenius 互反及两个虚特征标的内积为整数。以不可约特征标为正交规范基展开 \(f\)，所有系数均为整数，故 \(f\in R(G)\)。
\end{proof}

\subsection{初等群表示与一维表示的诱导}
\label{b5:subsec:0504:04}

下面用普通复表示论把定理化为一维特征标的诱导。
\begin{theorem}\label{b5:thm:p-group-monomial}
设 \(p\) 为素数，\(P\) 为有限 \(p\)-群。\(P\) 的每个复不可约特征标 \(\theta\) 都可写成 \(\operatorname{Ind}_Q^P\mu\)，其中 \(Q\le P\)，\(\mu\) 为 \(Q\) 的一维复特征标。
\end{theorem}
\begin{proof}
按 \(|P|\) 归纳。若表示有非平凡核 \(N\)，它来自较小 \(p\)-群 \(P/N\) 的不可约表示。由归纳假设，这个表示等于某个 \(\overline Q\le P/N\) 的一维特征标 \(\overline\mu\) 的诱导。取原像 \(Q\le P\)，令 \(\mu(q)=\overline\mu(qN)\)。映射
\[
\mathbb C[P]\otimes_{\mathbb C[Q]}\mathbb C_\mu
\longrightarrow
\mathbb C[P/N]\otimes_{\mathbb C[\overline Q]}\mathbb C_{\overline\mu},
\qquad g\otimes z\longmapsto(gN)\otimes z
\]
与张量平衡关系相容，因为 \(\mu(q)=\overline\mu(qN)\)，并与 \(P\) 作用相容。陪集双射 \(P/Q\simeq(P/N)/\overline Q\) 说明它把一组陪集基双射到另一组，故为同构。于是原表示为 \(\operatorname{Ind}_Q^P\mu\)。

以下设表示 \(\rho\) 忠实。若 \(P\) 交换，由\cref{b5:cor:schur-algebraically-closed}，全部 \(\rho(g)\) 为标量，不可约性迫使维数为 \(1\)；这也包含归纳起点 \(P=1\)。

设 \(P\) 非交换。非平凡有限 \(p\)-群中心非平凡，因为类方程中非中心共轭类的大小均被 \(p\) 整除，故中心的阶亦被 \(p\) 整除。将此结论用于非平凡 \(p\)-群 \(P/Z(P)\)，选取 \(xZ(P)\ne1\) 位于其中心。子群
\(A=\bigl\langle Z(P),x\bigr\rangle\)
交换且正规：对任意 \(g\in P\)，有 \(gxg^{-1}\in xZ(P)\)。又因 \(x\notin Z(P)\)，\(A\) 不包含于 \(Z(P)\)。

由于 \(A\) 交换，它的不可约复表示都是一维特征标。将 \(\rho\) 的表示空间 \(V\) 限制到 \(A\)，按这些特征标的完整等型分量分解：
\[
V=\bigoplus_\lambda V_\lambda,\qquad
V_\lambda=\{\,v\in V:\rho(a)v=\lambda(a)v\text{ 对所有 }a\in A\,\}.
\]
每个非零 \(V_\lambda\) 中可以有多个同类型的一维成分，因而 \(\dim V_\lambda\) 不必为一。若只出现一个特征标，则每个 \(a\in A\) 都作用为标量 \(\rho(a)=\lambda(a)I\)。对任意 \(g\in P\)，于是有
\[
\rho([a,g])
=\rho(a)\rho(g)\rho(a)^{-1}\rho(g)^{-1}=I.
\]
忠实性给出 \([a,g]=1\)，故 \(A\subseteq Z(P)\)，矛盾。因此至少出现两个权。

取其中一个 \(\lambda\)，其惯性群 \(J=I_P(\lambda)\) 为真子群；这是因为\cref{b5:thm:clifford-restriction}断言不可约表示中的所有权构成一个 \(P\)-轨道，而这里已有至少两个权。完整权空间 \(V_\lambda\) 被 \(J\) 保持：对 \(j\in J\)、\(a\in A\)、\(v\in V_\lambda\)，有
\[
\rho(a)\rho(j)v
=\rho(j)\rho(j^{-1}aj)v
=\lambda(j^{-1}aj)\rho(j)v
=\lambda(a)\rho(j)v.
\]
它作为 \(J\)-表示也不可约。事实上，若 \(0\ne Z\subseteq V_\lambda\) 是 \(J\)-子模，则对 \(P/J\) 的一组代表 \(g\)，各 \(\rho(g)Z\) 位于不同的权空间；它们的直和被 \(P\) 保持，故由 \(V\) 的不可约性等于 \(V\)。取其中的 \(\lambda\)-权分量便得 \(Z=V_\lambda\)。因此\cref{b5:thm:clifford-correspondence}在这里给出的 \(U\) 正是整个 \(V_\lambda\)，且 \(\rho\simeq\operatorname{Ind}_J^P U\)。对较小的 \(p\)-群 \(J\) 应用归纳假设，得 \(U\simeq\operatorname{Ind}_Q^J\mu\)。\cref{b5:thm:induction-transitivity}遂给出 \(\rho\simeq\operatorname{Ind}_Q^P\mu\)。
\end{proof}

\begin{proposition}\label{b5:prop:elementary-monomial}
有限初等群 \(E\) 的每个复不可约特征标都是其某个初等子群的一维复特征标的诱导。
\end{proposition}
\begin{proof}
写 \(E=P\times C\)。在不可约表示中，中心子群 \(C\) 作用为某个一维特征标 \(\lambda\)，而限制到 \(P\) 是不可约表示 \(\theta\)：若有真 \(P\)-不变子空间，标量的 \(C\)-作用也保持它，与 \(E\)-不可约性矛盾。由\cref{b5:thm:p-group-monomial}，\(\theta=\operatorname{Ind}_Q^P\mu\)。取 \(P/Q\) 的陪集基，便有
\[
\theta\otimes\lambda
=\operatorname{Ind}_{Q\times C}^{P\times C}(\mu\otimes\lambda).
\]
其中 \(Q\times C\) 仍是 \(p\)-初等群，且 \(\mu\otimes\lambda\) 是它的一维特征标，得到所需表达。
\end{proof}

\begin{proof}[\cref{b5:thm:brauer-induction}的证明]
由\cref{b5:prop:brauer-integral-image}，任意 \(\alpha\in R(G)\) 是初等子群虚特征标诱导的整数线性组合；把这些虚特征标展开为不可约特征标的整数线性组合，便得到定理的第一项结论。对每个不可约初等子群特征标应用\cref{b5:prop:elementary-monomial}，将其写成同一初等群内某个初等子群的一维诱导。最后用\cref{b5:thm:induction-transitivity}把两次诱导合成一次，即得全部系数仍为整数的一维表述。
\end{proof}

\subsection{负系数与虚特征标的例子}
\label{b5:subsec:0504:05}

取 \(G=S_3\)，其不可约特征标为 \(1,\varepsilon,\chi\)，次数分别为 \(1,1,2\)。设 \(C_2=\bigl\langle(12)\bigr\rangle\)，\(C_3=\bigl\langle(123)\bigr\rangle\)，取 \(C_3\) 的非平凡一维特征标 \(\lambda\)。由限制及 Frobenius 互反，
\[
\operatorname{Ind}_{C_2}^{S_3}1=1+\chi,\qquad
\operatorname{Ind}_{C_3}^{S_3}\lambda=\chi.
\]
第一式也可从三点置换表示中分离常数向量得到；第二式在三类上的值依次为 \(2,0,-1\)。所以
\begin{equation}\label{b5:eq:s3-brauer-negative}
1=\operatorname{Ind}_{C_2}^{S_3}1
-\operatorname{Ind}_{C_3}^{S_3}\lambda.
\end{equation}
两个子群都循环，因而都是允许的初等子群。此式移项后表示一个三维实际表示同构于一维平凡表示与一个二维表示的直和；负号没有制造负维数的表示空间。

\ProseParagraphBreak
\(S_3\) 自身不是初等群，所有初等子群都是真子群，故每个一维诱导的次数至少为 \(2\)。因此 \(1\) 不可能是这些诱导的非负整数线性组合。这也说明负系数不是上述公式偶然写得不巧，而是整诱导定理必须允许的现象。

% !TeX root = ../../Book5.tex
% 纲要标题：### 5.5 共轭类、特征标约束与 Burnside 可解性定理
% 纲要位置：工作记录/计划纲要.md，第 3882 行。
% 纲要细目（写作范围）：
% * 先准确陈述 Burnside 的 \(p^aq^b\) 阶群可解性定理，再完成下列共轭类论证；允许其中一个指数为零，并分别处理平凡群和素数幂阶群。
% * 从类和作用和代数整数性建立群论所需的特征标消失或标量作用引理；注明次数与共轭类大小互素的使用位置
% * 证明有限非交换单群不存在非平凡的素数幂大小共轭类，并处理类大小为1的中心情形
% * 对 \(|G|=p^aq^b\) 利用 Sylow 子群中心中的元素构造所需共轭类，再按群阶归纳并使用可解性的扩张稳定性完成证明
% * 以中心化子指数说明素数幂共轭类对有限非交换单群的限制

\section{共轭类、特征标约束与 Burnside 可解性定理}
\label{b5:sec:0505}

整性还能把“某个共轭类很小”变成正规子群的存在性。以下从特征标正交、\cref{b5:prop:central-character-integrality}与代数整数的基本运算出发。

\subsection{\texorpdfstring{Burnside 的 \(p^aq^b\) 阶群可解性定理}{Burnside 的 paqb 阶群可解性定理}}
\label{b5:subsec:0505:01}

对群 \(G\)，记 \(G'=\langle xyx^{-1}y^{-1}:x,y\in G\rangle\)，递归定义 \(G^{(0)}=G\)、\(G^{(n+1)}=\bigl(G^{(n)}\bigr)'\)。若某个 \(G^{(r)}=1\)，则称 \(G\) 为\term[kejiequn]{可解群}[solvable group]。这一定义等价于存在每个相邻商群均交换的有限次正规列
\[
G=G_0\supseteq G_1\supseteq\cdots\supseteq G_r=1,
\qquad G_{i+1}\trianglelefteq G_i.
\]
导出列本身给出一个方向；反过来，\(G_i/G_{i+1}\) 交换意味着 \(G_i'\subseteq G_{i+1}\)，由归纳得到 \(G^{(i)}\subseteq G_i\)，因而 \(G^{(r)}=1\)。

\begin{theorem}[Burnside 可解性定理]\label{b5:thm:burnside-solvability}
设 \(p,q\) 为素数，\(a,b\) 为非负整数。每个阶为 \(p^a q^b\) 的有限群 \(G\) 都是可解群。
\end{theorem}
此结论称为\term[Burnsidekejie]{Burnside 可解性定理}[Burnside's \(p^a q^b\) theorem]。允许 \(a=0\)、\(b=0\)；若 \(p=q\)，结论仍归入素数幂阶情形。证明在本节先建立共轭类工具，再于\subsecref{b5:subsec:0505:04}完成。

\ProseParagraphBreak
定理并不说 \(G\) 交换，也不说全部 Sylow 子群正规。例如 \(S_3\) 阶为 \(2\cdot3\)，其导出群为 \(A_3\)，故可解，但三阶以外的 Sylow 子群并不正规。这里的目标是在导出列中逐步消去非交换部分。

\subsection{类和作用与特征标消失引理}
\label{b5:subsec:0505:02}

\begin{lemma}[单位根平均值]\label{b5:lem:root-unity-average}
设 \(d\ge1\) 为整数，\(\xi_1,\ldots,\xi_d\) 为复单位根。若
\[
\alpha=\frac{\xi_1+\cdots+\xi_d}{d}
\]
是代数整数，则或者 \(\alpha=0\)，或者 \(\xi_1=\cdots=\xi_d\)。
\end{lemma}
\begin{proof}
若这些单位根不全相同，三角不等式的等号条件给出 \(|\alpha|<1\)。选取公倍数 \(m\)，使各 \(\xi_i\) 都属于 \(\mathbb Q(\zeta_m)\)。\(\alpha\) 的每个代数共轭仍是 \(d\) 个单位根的平均，故绝对值不超过 \(1\)。这里的共轭描述可这样核验：\(\mathbb Q(\alpha)\) 到 \(\mathbb C\) 的任一嵌入都可延拓到 \(\mathbb Q(\alpha,\zeta_m)\)，只需把 \(\zeta_m\) 在前一域上的极小多项式的系数施行嵌入，再选一个复根；在该根处求值给出商域上的延拓。延拓保持 \(\zeta_m^m=1\)，因而保持“是单位根”这一性质。

若 \(\alpha\ne0\)，其全部共轭根的乘积非零；由\cref{b5:lem:algebraic-integer-arithmetic}该乘积为整数。但其中一个因子的绝对值严格小于 \(1\)，其余不超过 \(1\)，故乘积的绝对值严格小于 \(1\)，矛盾。
\end{proof}

\begin{lemma}[消失或标量作用]\label{b5:lem:coprime-class-scalar}
设 \(G\) 为有限群，\(\rho:G\rightarrow\operatorname{GL}_{\mathbb C}(V)\) 为次数 \(d\) 的复不可约表示，特征标为 \(\chi\)，\(g\in G\)。设 \(K=g^G\) 且 \(\gcd(d,|K|)=1\)。则或者 \(\chi(g)=0\)，或者 \(\rho(g)\) 为标量矩阵。
\end{lemma}
\begin{proof}
取整数 \(a,b\)，使 \(a|K|+bd=1\)。由\cref{b5:prop:central-character-integrality}，\(|K|\cdot\chi(g)/d\) 是代数整数；\(\chi(g)\) 也由\cref{b5:prop:character-values-integral}为代数整数。因此
\[
\frac{\chi(g)}d
=a\cdot\frac{|K|\cdot\chi(g)}d+b\cdot\chi(g)
\]
是代数整数。这正是互素条件的使用处。\(\rho(g)\) 可对角化，其特征值 \(\xi_1,\ldots,\xi_d\) 为单位根；由\cref{b5:lem:root-unity-average}，或者它们的和为零，或者全部相同。后者连同可对角化性给出标量矩阵。
\end{proof}

不能删去互素条件。例如 \(S_3\) 的二维不可约表示在三循环处取值 \(-1\)，并非零；对应矩阵的两个特征值为互不相同的本原三次单位根，所以不是标量。该共轭类大小为 \(2\)，恰与表示次数有公因子。

\subsection{非交换单群的素数幂大小共轭类}
\label{b5:subsec:0505:03}

\begin{proposition}\label{b5:prop:prime-power-class-normal-subgroup}
设 \(G\) 为有限群，某个共轭类 \(K\) 的大小为 \(p^r\)，其中 \(p\) 为素数、\(r\ge1\)。则 \(G\) 有非平凡真正规子群。特别地，有限非交换单群没有这样的共轭类。
\end{proposition}
\begin{proof}
取 \(g\in K\)，则 \(g\ne1\)。由\cref{b5:thm:regular-decomposition}，正则特征标在非单位元处为零，故
\begin{equation}\label{b5:eq:burnside-regular-sum}
0=\sum_{\chi\in\operatorname{Irr}(G)}\chi(1)\chi(g).
\end{equation}
我们先证明存在非平凡不可约 \(\chi\)，满足 \(p\nmid\chi(1)\) 且 \(\chi(g)\ne0\)。否则，\eqref{b5:eq:burnside-regular-sum}除去平凡特征标的贡献 \(1\)，其余可能非零的项的次数均被 \(p\) 整除，从而
\[
-\frac1p
=\sum_{\substack{\chi\in\operatorname{Irr}(G)\\p\mid\chi(1)}}
\frac{\chi(1)}p\,\chi(g)
\]
为代数整数。这与\cref{b5:lem:algebraic-integer-arithmetic}的有理数判别矛盾。

取这样的 \(\chi\) 及表示 \(\rho\)。因为 \(\gcd\bigl(\chi(1),p^r\bigr)=1\) 且 \(\chi(g)\ne0\)，\cref{b5:lem:coprime-class-scalar}说明 \(\rho(g)\) 为标量。所有 \(K\) 中的元素共轭于 \(g\)，所以它们在 \(\rho\) 下给出同一个标量。令
\[
N=\langle xy^{-1}:x,y\in K\rangle.
\]
共轭保持 \(K\)，故 \(N\triangleleft G\)；又 \(N\subset\ker\rho\)。由于 \(\rho\) 非平凡，\(\ker\rho\ne G\)，所以 \(N\ne G\)。而 \(|K|>1\) 给出不同的 \(x,y\)，此时 \(xy^{-1}\ne1\)，所以 \(N\ne1\)。
\end{proof}

大小为 \(1\) 的类须另行处理，它对应中心元素。有限非交换单群的中心只能为 \(1\)：中心若非平凡，单性使整个群交换；有限交换单群则必为素数阶循环群。因此非交换单群的非单位元既不能有大小为 \(1\) 的类，也不能有更大的素数幂大小的类。

\subsection{Sylow 中心、群阶归纳与可解扩张}
\label{b5:subsec:0505:04}

先记录可解性的群论运算，避免把归纳的最后一步藏在一句“因此可解”中。
\begin{lemma}\label{b5:lem:solvable-extensions}
设 \(G\) 为群，\(N\triangleleft G\)。若 \(N\) 与 \(G/N\) 都可解，则 \(G\) 可解。可解群的子群和商群也可解。对每个素数 \(p\)，每个有限 \(p\)-群都可解。
\end{lemma}
\begin{proof}
若 \((G/N)^{(r)}=1\)，商映射保持交换子，故 \(G^{(r)}\subset N\)。若 \(N^{(s)}=1\)，则 \(G^{(r+s)}\subset N^{(s)}=1\)。子群 \(H\le G\) 满足 \(H^{(i)}\subset G^{(i)}\)，而满同态把导出列映为商群的导出列，由此得到另外两项稳定性。

对有限 \(p\)-群 \(P\) 按阶归纳。若 \(P\ne1\)，其类方程中非中心类的大小均被 \(p\) 整除，故 \(|Z(P)|\) 被 \(p\) 整除，中心非平凡。若 \(Z(P)=P\)，群交换；否则 \(P/Z(P)\) 阶较小，由归纳假设可解。中心自身交换，再由已经证明的扩张性质得 \(P\) 可解。
\end{proof}

\begin{proof}[\cref{b5:thm:burnside-solvability}的证明]
平凡群可解；若 \(a=0\)、\(b=0\) 或 \(p=q\)，由\cref{b5:lem:solvable-extensions}的有限 \(p\)-群结论可得。以下设 \(p\ne q\) 且 \(a,b>0\)，按 \(|G|\) 归纳。

若 \(G\) 有非平凡真正规子群 \(N\)，则 \(N\) 与 \(G/N\) 的阶仍只含素因子 \(p,q\)，且都小于 \(|G|\)。归纳假设使两者可解，再由\cref{b5:lem:solvable-extensions}得 \(G\) 可解。只需排除 \(G\) 为单群的情形。它不可能是交换单群，因为交换单群只有一个素因子。

取 Sylow \(p\)-子群 \(P\)。由 \(p\)-群的类方程，选 \(1\ne x\in Z(P)\)。于是 \(P\subset C_G(x)\)，故
\[
|x^G|=\bigl[G:C_G(x)\bigr]
\]
是 \(q\) 的幂。若这个指数为 \(1\)，则 \(x\in Z(G)\setminus\{1\}\)，与非交换单性矛盾；若指数大于 \(1\)，则\cref{b5:prop:prime-power-class-normal-subgroup}给出非平凡真正规子群，同样矛盾。因此 \(G\) 不可能为单群，归纳即得可解性。
\end{proof}

这个证明需要的是在 Sylow 子群的中心中取元素，不要求所取 Sylow 子群在 \(G\) 中正规。整性与共轭类步骤可对照 Etingof 等\cite[Theorems 5.4.3, 5.4.4, 5.4.6; Lemmas 5.4.5, 5.4.7, pp. 98--100]{EtingofEtAl2011RepresentationTheory}；这里的最后一步以 Sylow 中心显式构造所需共轭类。

\subsection{素数幂共轭类对单群的限制}
\label{b5:subsec:0505:05}

有限非交换单群的每个非单位元 \(g\) 都必须满足：中心化子指数 \([G:C_G(g)]\) 既不为 \(1\)，也不是任何素数的正整数次幂。因此每个非平凡共轭类的大小至少含两个不同素因子。这个限制施加在共轭类大小上，而不只是群阶上；对只含两个素因子的群，Sylow 子群中心提供的元素使中心化子指数只剩一个素因子，正是在这一点与单性矛盾。下一节将用可解性进一步确定合成因子与极小正规子群的形状。

% !TeX root = ../../Book5.tex
% 纲要标题：### 5.6 群论应用与三种 Burnside 定理的比较
% 纲要位置：工作记录/计划纲要.md，第 3890 行。
% 纲要细目（写作范围）：
% * Burnside 的 p^a q^b 阶群可解性定理
% * 与第二卷 Burnside 矩阵代数定理的区别
% * 结构结论如何由特征标约束推出
% * 本节运用第5.5节已经证明的 \(p^aq^b\) 阶群可解性定理，整理群论推论及适用范围；将其与轨道计数引理、矩阵代数定理逐项比较
% ---

\section{群论应用与三种 Burnside 定理的比较}
\label{b5:sec:0506}

可解性是一个结构结论，尚未给出具体群的分类。使用它时，通常再取极小或极大正规子群，把群拆成较简单的扩张。与此同时，特征标方法还可以研究置换作用和矩阵代数；应当分清每次得到了什么。

\subsection{\texorpdfstring{\(p^aq^b\) 阶群可解性定理的应用}{paqb 阶群可解性定理的应用}}
\label{b5:subsec:0506:01}

\begin{corollary}\label{b5:cor:two-prime-composition}
设 \(p,q\) 为素数，\(a,b\ge0\) 为整数，\(G\) 为阶大于 \(1\) 且 \(|G|=p^a q^b\) 的有限群。则 \(G\) 有素数阶循环商群，并且其每个单合成因子都是阶为 \(p\) 或 \(q\) 的循环群。特别地，有限非交换单群的阶至少含三个不同素因子。
\end{corollary}
\begin{proof}
由\cref{b5:thm:burnside-solvability}，\(G\) 可解。取极大真正规子群 \(N\)，则 \(G/N\) 为非平凡单群，并由\cref{b5:lem:solvable-extensions}仍可解。非平凡单可解群必须交换：其导出群若为自身，导出列不会终止；否则单性使导出群为 \(1\)。有限交换单群为素数阶循环群，阶又整除 \(|G|\)，故只能为 \(p\) 或 \(q\)。对子群的简单商同样应用此论证，即得全部合成因子结论。最后一项是\cref{b5:thm:burnside-solvability}的直接逆否应用。
\end{proof}

这并不使扩张自动分裂，也不能由合成因子的清单恢复乘法。例如 \(C_4\) 与 \(C_2\times C_2\) 都有两个 \(C_2\) 合成因子，却不同构。表示论给出的限制常要与具体正规子群及共轭作用一起使用。

\subsection{与 Burnside 矩阵代数定理的比较}
\label{b5:subsec:0506:02}

为便于对照，给出矩阵版本及一个有限维证明；它不参加上一节的可解性论证。
\begin{theorem}[Burnside 矩阵代数定理]\label{b5:thm:burnside-matrix-algebra}
设 \(k\) 为代数闭域，\(V\ne0\) 为有限维 \(k\)-向量空间，\(A\subset\operatorname{End}_k(V)\) 为含恒等映射的 \(k\)-子代数。若 \(V\) 没有非零真 \(A\)-不变子空间，则 \(A=\operatorname{End}_k(V)\)。
\end{theorem}
\term[Burnsidejuzhendaishu]{Burnside 矩阵代数定理}[Burnside's matrix algebra theorem]的假设中没有有限群，也没有两个素因子。
\begin{proof}
首先，每个与 \(A\) 交换的自同态 \(T\) 都是标量。由于 \(k\) 代数闭，\(T\) 有特征值 \(\lambda\)，非零子空间 \(\ker(T-\lambda I)\) 为 \(A\)-不变，故等于 \(V\)。

需要一个简单模的初等事实。把 \(V^r\) 写成坐标简单 \(A\)-模的直和。对任意子模 \(W\)，从这些坐标模中取一个最大的子集，使其直和 \(U\) 满足 \(U\cap W=0\)。若 \(U+W\ne V^r\)，便有一个坐标简单模 \(S\) 不包含于 \(U+W\)；由简单性，\(S\cap(U+W)=0\)，所以可以把 \(S\) 加入所选子集而仍与 \(W\) 零交，矛盾。故 \(V^r=W\oplus U\)。特别地，若 \(W\ne V^r\)，投影到 \(U\) 中一个坐标简单模给出非零 \(A\)-线性映射 \(V^r\rightarrow V\)，且零化 \(W\)。

现取 \(V\) 的基 \(v_1,\ldots,v_d\)，考察
\[
M=\bigl\{(av_1,\ldots,av_d):a\in A\bigr\}\subset V^d.
\]
它是 \(A\)-子模。若 \(M\ne V^d\)，前一事实给出零化 \(M\) 的非零 \(A\)-线性映射 \(\ell:V^d\rightarrow V\)。它在每个坐标上由第一段为标量，故
\(\ell(w_1,\ldots,w_d)=\sum_i c_iw_i\)，其中 \(c_i\) 不全为零。对 \(a=I\) 即得 \(\sum_i c_i v_i=0\)，与基的线性无关性矛盾。因此 \(M=V^d\)。任意指定各 \(v_i\) 的像都能由某个 \(a\in A\) 实现，故 \(A\) 是全部自同态代数。
\end{proof}

对复不可约有限群表示 \(\rho\)，代数 \(A=\operatorname{span}_{\mathbb C}\rho(G)\) 满足上述条件，所以群矩阵的线性生成空间就是全部矩阵代数。这并不声称群像 \(\rho(G)\) 等于全部可逆矩阵，更不推出 \(G\) 可解。若底域不是代数闭域，结论可能失败：\(\mathbb C\) 通过实线性乘法作用在 \(\mathbb R^2\) 上不可约，但所得二维实子代数不等于四维的 \(M_2(\mathbb R)\)。

\subsection{由特征标约束推出群结构}
\label{b5:subsec:0506:03}

\begin{proposition}\label{b5:prop:minimal-normal-solvable}
设 \(G\) 为有限可解群，\(N\triangleleft G\) 为非平凡极小正规子群。则存在素数 \(\ell\) 与正整数 \(r\)，使 \(N\simeq C_\ell^r\)。若另有素数 \(p,q\) 及非负整数 \(a,b\) 满足 \(|G|=p^a q^b\)，则 \(\ell\) 为 \(p\) 或 \(q\)。
\end{proposition}
\begin{proof}
由\cref{b5:lem:solvable-extensions}，\(N\) 可解。其导出群 \(N'\) 是 \(N\) 的特征子群，因而在 \(G\) 中正规。极小性说明 \(N'=1\) 或 \(N'=N\)；后一情形使非平凡的 \(N\) 的导出列永不终止，故 \(N\) 交换。选一个整除 \(|N|\) 的素数 \(\ell\)。有限交换群中所有 \(\ell\)-幂阶元素组成非平凡特征子群，极小性使其等于 \(N\)，所以 \(N\) 为交换 \(\ell\)-群。其子群 \(N[\ell]=\{x:x^\ell=1\}\) 同样非平凡且特征，故等于 \(N\)。用群乘法作加法、用整数次幂作 \(\mathbb F_\ell\) 的标量作用，\(N\) 成为有限维 \(\mathbb F_\ell\)-向量空间，因而 \(N\simeq C_\ell^r\)。最后由阶的整除性得到 \(\ell\in\{p,q\}\)。
\end{proof}

从小共轭类到正规子群的联系也有显式构造。\cref{b5:prop:prime-power-class-normal-subgroup}中的
\(N=\langle xy^{-1}:x,y\in K\rangle\)
只由共轭类 \(K\) 决定；它在找到的不可约表示下全被送到恒等矩阵。以 \(S_3\) 的三元换位类为例，其大小为 \(3\)，两个不同换位的商是三循环，所构造的 \(N\) 正是 \(A_3\)。因此证明不只是排除单性，有时还直接指出一个有用的正规子群。

\subsection{群论推论、适用范围与同名定理}
\label{b5:subsec:0506:04}

轨道计数公式也可完全放在特征标语言中。若有限群 \(G\) 作用在有限集合 \(X\) 上，置换特征标为 \(\pi_X(g)=|X^g|\)。不变向量在各轨道上系数恒定，因此其维数为轨道数；也可直接计算
\[
\dim\mathbb C[X]^G
=\langle1_G,\pi_X\rangle_G
=\frac1{|G|}\sum_{g\in G}|X^g|.
\]
从双重计数看，每个轨道的贡献是
\(\sum_{x\in\mathcal O}|G_x|=|\mathcal O|\cdot|G_x|=|G|\)，同样得到公式。

\ProseParagraphBreak
三种结论于是各有明确用途：轨道计数求不变向量的维数或轨道数；矩阵代数定理识别不可约作用所张成的全部代数；\(p^a q^b\) 定理由类和整性推出群的可解性。前两者对 \(A_5\) 一样适用，但 \(|A_5|=60=2^2\cdot3\cdot5\) 已超出第三者的两个素因子假设。对阶含三个以上素因子的群，仅凭本章定理既不能断言可解，也不能断言不可解。


\begin{chaptersummary}
表示环是以不可约特征标为基的整数格，张量积给出乘法。特征标值由单位根相加而成，类和标量也满足首一整系数方程；将这些整性与有理数判别结合，便得到不可约次数整除群阶。

\ProseParagraphBreak
Artin 诱导由循环子群上的限制检测零类函数，再利用整数矩阵的秩获得有理张成。它允许一般的一维特征标；若只从平凡特征标诱导，则必须相应限制为在有理共轭类上常值的函数。循环诱导格的有限指数可能大于 \(1\)，不能把清分母误当成整系数结论。

\ProseParagraphBreak
Brauer 诱导扩大到初等子群。对每个素数，表示环的极大理想都可通过一个 \(p\)-正则元素的模素理想评价检测；中心化子中的 Sylow 子群与该元素生成的循环群组成所需初等子群。有限 Fourier 求和排除了诱导理想落入极大理想的可能，再由逐素数检测回到整数。Clifford 对应给出的 \(p\)-群单项式性则将初等子群的特征标进一步写成一维特征标的诱导。

\ProseParagraphBreak
Burnside 可解性定理沿独立的论证使用整性：次数与类大小互素时，特征标或消失，或对应矩阵为标量；正则特征标之和随后从素数幂大小的类中检出正规子群。Sylow 中心和可解扩张完成按群阶的归纳。其结论、轨道计数公式和不可约矩阵代数定理分别处理群结构、群作用计数和线性算子代数，不宜只凭同名而互相替代。
\end{chaptersummary}

\BookExercises

\begin{exercise}\label{b5:ex:05-cyclic-ring}
设 \(C_n=\langle c\rangle\)，一维特征标 \(\lambda\) 满足 \(\lambda(c)=e^{2\pi i/n}\)。证明 \(R(C_n)\simeq\mathbb Z[T]/(T^n-1)\)，并在此描述次数映射及其核。
\end{exercise}
\begin{solution}
循环群的不可约特征标为 \(1,\lambda,\ldots,\lambda^{n-1}\)，乘法使指数模 \(n\) 相加。因此 \(T\mapsto\lambda\) 给出满同态。商环中的每个元素唯一写成次数小于 \(n\) 的整系数多项式，而这些特征标在加法群中线性无关，所以同态单射。次数把每个 \(\lambda^j\) 送到 \(1\)，对应 \(f(T)\mapsto f(1)\)。整系数除法给出 \(f(1)=0\) 当且仅当 \(f\) 被 \(T-1\) 整除，故核由 \(T-1\) 的剩余类生成。
\end{solution}

\begin{exercise}\label{b5:ex:05-s3-ring}
设 \(S_3\) 的不可约特征标为 \(1,\varepsilon,\chi\)，次数为 \(1,1,2\)。计算其表示环乘法，并判断 \(\chi-\varepsilon\) 是否为实际特征标。
\end{exercise}
\begin{solution}
在单位、换位、三循环三类上，\(\varepsilon=(1,-1,1)\)、\(\chi=(2,0,-1)\)。逐点乘法给出
\[
\varepsilon^2=1,\qquad\varepsilon\chi=\chi,\qquad
\chi^2=1+\varepsilon+\chi.
\]
三者为加法基，故这些关系确定整个环。\(\chi-\varepsilon\) 次数为 \(1\)，但 \(\langle\varepsilon,\chi-\varepsilon\rangle=-1\)，不可约系数为负，因此不是实际特征标。
\end{solution}

\begin{exercise}\label{b5:ex:05-cyclotomic-integer}
令 \(\zeta=e^{2\pi i/5}\)，\(a=\zeta+\zeta^{-1}\)。求 \(a\) 的首一整系数极小多项式，并说明代数整数不必是整数。
\end{exercise}
\begin{solution}
由 \(1+\zeta+\zeta^2+\zeta^3+\zeta^4=0\)，有
\(a^2=\zeta^2+\zeta^{-2}+2=1-a\)，故 \(a^2+a-1=0\)。多项式 \(T^2+T-1\) 的判别式为 \(5\)，不是有理数平方，因此在 \(\mathbb Q[T]\) 中不可约。于是 \(a\) 是次数为 \(2\) 的代数整数而非有理数，当然也非整数。
\end{solution}

\begin{exercise}\label{b5:ex:05-class-sums-s3}
在 \(\mathbb Z[S_3]\) 中，令 \(T\) 为三个换位的和，\(C\) 为两个三循环的和。计算 \(T^2,TC,C^2\)，并求二者在二维不可约表示上的标量作用。
\end{exercise}
\begin{solution}
同一换位相乘给出单位，不同换位的有序乘积各三次给出两种三循环；换位与三循环的六个乘积各两次给出三个换位；三循环相乘则给出两次单位及各一次三循环。因此
\[
T^2=3\cdot1+3C,\qquad TC=2T,\qquad C^2=2\cdot1+C.
\]
二维特征标在换位处为 \(0\)，在三循环处为 \(-1\)，由中心标量公式得 \(T\) 作用为 \(0\)、\(C\) 作用为 \(-1\)。这些标量也满足所算乘法关系。
\end{solution}

\begin{exercise}\label{b5:ex:05-p-group-degrees}
设 \(|P|=p^a\)。证明每个复不可约表示的次数都是 \(p\) 的幂；若 \(P\ne1\) 且存在次数 \(d>1\) 的不可约表示，说明为什么 \(d^2<|P|\)。
\end{exercise}
\begin{solution}
由\cref{b5:cor:character-degree-divisibility}，\(d\mid p^a\)，所以 \(d=p^b\)。由正则分解，全部不可约次数的平方和为 \(|P|\)。除了所取不可约表示，还至少有平凡表示，其平方贡献为 \(1\)，故 \(d^2+1\le|P|\)，从而严格小于 \(|P|\)。
\end{solution}

\begin{exercise}\label{b5:ex:05-rational-values-no-real-model}
四元数群 \(Q_8=\{\pm1,\pm i,\pm j,\pm k\}\) 的二维复不可约特征标在 \(1,-1\) 处分别为 \(2,-2\)，其余取值为 \(0\)。证明它不能由二维实表示实现，因而也不能由二维有理表示实现。
\end{exercise}
\begin{solution}
若有这样的实表示，\(-1\) 的矩阵为 \(-I\)：它阶为 \(2\)、可对角化且迹为 \(-2\)。因此 \(i,j\) 的矩阵 \(A,B\) 满足 \(A^2=B^2=-I\)、\(AB=-BA\)。选取非零 \(v\)，则 \(v,Av\) 实线性无关，可把 \(A\) 写成
\(\left(\begin{smallmatrix}0&-1\\1&0\end{smallmatrix}\right)\)。
反交换条件迫使
\(B=\left(\begin{smallmatrix}u&v\\v&-u\end{smallmatrix}\right)\)，
故 \(B^2=(u^2+v^2)I\)，不可能为 \(-I\)。有理模型扩标量后会给实模型，亦被排除。特征标值都是整数，并未保证矩阵可在其值域上实现。
\end{solution}

\begin{exercise}\label{b5:ex:05-klein-exact-index}
令 \(G=C_2\times C_2\)。证明循环子群诱导格 \(I_{\mathrm{cyc}}(G)\) 恰由偶数次数的虚特征标组成，故 \(R(G)/I_{\mathrm{cyc}}(G)\simeq\mathbb Z/2\mathbb Z\)。
\end{exercise}
\begin{solution}
用四个一维特征标为加法基 \(e_1,\ldots,e_4\)。每个二阶子群的一维诱导是两个限制相同的特征标之和；三个二阶子群给出四元素集合的三个配对，故全部 \(e_i+e_j\)、\(i\ne j\)，都是这样的诱导。它们次数均为偶数。反过来，若 \(i,j,k\) 不同，则
\[
e_i-e_j=(e_i+e_k)-(e_j+e_k),\qquad
2e_i=(e_i+e_j)+(e_i+e_k)-(e_j+e_k).
\]
因此任何坐标和为偶数的整数向量都在该格中。次数模 \(2\) 的映射便给出所述商群。
\end{solution}

\begin{exercise}\label{b5:ex:05-elementary-abelian-denominator}
设 \(G=C_p^r\)、\(r\ge2\)，\(\mathcal L\) 为其全部阶为 \(p\) 的子群。求使 \(N1_G\in I_{\mathrm{cyc}}(G)\) 成立的最小正整数 \(N\)，并给出一个表达。
\end{exercise}
\begin{solution}
每个非平凡循环子群的指数为 \(p^{r-1}\)，平凡子群的指数为 \(p^r\)，所以所有循环诱导的次数都被 \(p^{r-1}\) 整除，必要条件是 \(p^{r-1}\mid N\)。子群数 \(L=|\mathcal L|=(p^r-1)/(p-1)\)，每个非单位元恰属一个这样的子群。于是逐点可验
\[
p^{r-1}1_G
=\sum_{H\in\mathcal L}\operatorname{Ind}_H^G1_H
-\frac{L-1}{p}\operatorname{Ind}_1^G1.
\]
在非单位元处两侧均为 \(p^{r-1}\)，在单位元处右侧为
\(Lp^{r-1}-(L-1)p^{r-1}\)。又 \((L-1)/p\) 为整数，故最小值正是 \(p^{r-1}\)。
\end{solution}

\begin{exercise}\label{b5:ex:05-permutation-artin-s3}
只使用 \(S_3\) 的子群 \(1,C_2,C_3\) 的平凡特征标诱导，求平凡特征标 \(1\) 的有理表达。说明即使所有不可约特征标均为有理值，整数系数仍可能不足。
\end{exercise}
\begin{solution}
记三个诱导分别为 \(\rho,\pi_2,\pi_3\)，则
\(\rho=1+\varepsilon+2\chi\)、\(\pi_2=1+\chi\)、\(\pi_3=1+\varepsilon\)。故
\[
1=\pi_2+\tfrac12\pi_3-\tfrac12\rho.
\]
若 \(1=a\rho+b\pi_2+c\pi_3\) 且 \(a,b,c\in\mathbb Z\)，比较 \(\chi,\varepsilon\) 系数得 \(b=-2a,c=-a\)，再比较平凡系数得 \(-2a=1\)，矛盾。允许非平凡子群特征标后，\eqref{b5:eq:s3-brauer-negative}才给出整数表达。
\end{solution}

\begin{exercise}\label{b5:ex:05-rational-classes-c5}
求 \(C_5\) 上在有理共轭类中常值的有理函数空间，并用置换特征标写出任意这样的函数。比较它与全部复类函数空间的维数。
\end{exercise}
\begin{solution}
有理共轭类只有 \(\{1\}\) 与四个非单位元组成的集合，故该有理空间维数为 \(2\)。若函数在这两处的值分别为 \(a,b\in\mathbb Q\)，则
\[
f=b\,1_{C_5}+\frac{a-b}{5}\operatorname{Ind}_1^{C_5}1.
\]
而 \(C_5\) 交换，每个元素单独成共轭类，全部复类函数空间维数为 \(5\)。减少到两个参数来自有理共轭条件，不是来自循环群特征标不足。
\end{solution}

\begin{exercise}\label{b5:ex:05-sylow-cover}
分别对 \(S_3\) 与 \(C_6\)，判断只从全部 Sylow 子群的表示诱导，能否有理张成整个复表示环。
\end{exercise}
\begin{solution}
\(S_3\) 的元素阶只有 \(1,2,3\)，每个元素都属于某个 Sylow 子群，且全部 Sylow 子群对共轭封闭。由\cref{b5:prop:subgroup-cover-induction}，有理张成成立。\(C_6\) 的生成元不属于二阶或三阶 Sylow 子群，且群交换，不能通过共轭进入它们。全部所述诱导在这个生成元处为零，故不能张成平凡特征标。
\end{solution}

\begin{exercise}\label{b5:ex:05-elementary-examples}
对 \(C_6\)、\(C_2^2\times C_3\)、\(Q_8\)、\(S_3\)，分别判断它们是否为 \(2\)-初等群、\(3\)-初等群。
\end{exercise}
\begin{solution}
\(C_6=C_2\times C_3\) 两种都是。\(C_2^2\times C_3\) 为 \(2\)-初等群，但不是 \(3\)-初等群，因为其与 \(3\) 互素的部分为非循环的 \(C_2^2\)。\(Q_8\) 为 \(2\)-初等群；若为 \(3\)-初等群，其 \(3\) 部分只能平凡，从而全群必须循环，矛盾。\(S_3\) 两种都不是：若有非平凡互素循环直因子，则该因子位于中心，而 \(Z(S_3)=1\)；若此因子平凡，全群又须为素数幂阶群。
\end{solution}

\begin{exercise}\label{b5:ex:05-elementary-quotient}
证明有限 \(p\)-初等群的任意商群仍为 \(p\)-初等群；并说明同样的说法不能把“直积”随意改成“半直积”后仍当作定义。
\end{exercise}
\begin{solution}
设 \(E=P\times C\)，\(N\triangleleft E\)。由\cref{b5:prop:elementary-subgroups}，\(N=(N\cap P)\times(N\cap C)\)，两个交分别在对应因子中正规，故
\[
E/N\simeq P/(N\cap P)\times C/(N\cap C).
\]
第一因子为 \(p\)-群，第二因子循环且阶与 \(p\) 互素。另一方面，\(S_3=C_3\rtimes C_2\) 是非平凡半直积，却不是初等群；交换两个因子的要求在定义及证明中都有实质作用。
\end{solution}

\begin{exercise}\label{b5:ex:05-p-regular-s3}
在 \(S_3\) 中，对 \(p=2,3\) 分别列出各共轭类的 \(p\)-正则部分，并对三个不可约特征标验证模 \(p\) 约化公式。
\end{exercise}
\begin{solution}
按 \(1,(12),(123)\) 排列，三个特征标为
\(1=(1,1,1)\)、\(\varepsilon=(1,-1,1)\)、\(\chi=(2,0,-1)\)，均为整数值。\(p=2\) 时换位的正则部分为 \(1\)，而三循环保持不变；于是需比较换位与单位元两列，逐项差为 \(0,-2,-2\)，均被 \(2\) 整除。\(p=3\) 时三循环的正则部分为 \(1\)，换位不变；三循环列与单位元列之差为 \(0,0,-3\)，均被 \(3\) 整除。
\end{solution}

\begin{exercise}\label{b5:ex:05-local-indices-klein}
令 \(G=C_2^2\)，计算 \(I_p(G)\) 在 \(p=2\) 及奇素数 \(p\) 时的指数。这与逐素数整诱导定理如何相容？
\end{exercise}
\begin{solution}
\(p=2\) 时全群 \(G\) 自身为 \(2\)-初等群，恒等诱导即给 \(I_2(G)=R(G)\)。若 \(p\) 为奇素数，任何子群的 \(p\) 部分都平凡，所以 \(p\)-初等子群恰是循环子群。\cref{b5:ex:05-klein-exact-index}的直接格计算给出指数 \(2\)。因此 \(I_p(G)\) 未必在整数上等于 \(R(G)\)，但其指数确实与该奇素数 \(p\) 互素；反演 \(2\) 后等式成立。
\end{solution}

\begin{exercise}\label{b5:ex:05-fourier-choice}
设 \(C=\langle c\rangle\) 阶为 \(5\)，\(T=\{c,c^{-1}\}\)。在特征 \(2\) 的域中取一个五阶单位根 \(z\)。比较平凡特征标与 \(\lambda_j(c)=z^j\) 给出的和 \(\sum_{t\in T}\lambda_j(t)\)，说明整诱导证明中为什么不能固定只用平凡特征标。
\end{exercise}
\begin{solution}
平凡特征标给出 \(1+1=0\)。对 \(1\le j\le4\)，若 \(z^j+z^{-j}=0\)，则在特征 \(2\) 中得 \(z^{2j}=1\)，与 \(z\) 阶为 \(5\) 矛盾。因此四个非平凡特征标均给出非零和。Fourier 检测允许选择某个一维特征标是必要的；平凡特征标的求和可能被剩余特征消去。
\end{solution}

\begin{exercise}\label{b5:ex:05-finiteness-lifting}
用包含 \(\mathbb Z\subset\mathbb Q\) 说明：极大理想提升引理中的有限生成假设不能直接删去。这里是否还有元素满足的首一整方程？
\end{exercise}
\begin{solution}
\(\mathbb Q\) 是域，唯一极大理想为 \(0\)，在 \(\mathbb Z\) 中的收缩仍为 \(0\)，并非 \((p)\)。因此任意素数理想 \((p)\subset\mathbb Z\) 都不能提升为 \(\mathbb Q\) 的极大理想。环 \(\mathbb Q\) 不是有限生成的 \(\mathbb Z\)-模，也不整于 \(\mathbb Z\)：例如 \(1/p\) 不是代数整数。问题不是模 \((p)\) 后选错了理想，而是 \(p\mathbb Q=\mathbb Q\)，此时整个纤维已经为零。
\end{solution}

\begin{exercise}\label{b5:ex:05-local-detection-primes}
设 \(A\) 为有限交换群。证明 \(A_{(p)}\) 恰保留 \(A\) 的 \(p\)-主分量，并用一个例子说明只检测固定的一个素数不能推出 \(A=0\)。
\end{exercise}
\begin{solution}
若元素阶与 \(p\) 互素，该阶在 \(\mathbb Z_{(p)}\) 中可逆，所以元素局部化后为零。若元素阶为 \(p^r\)，则与 \(p\) 互素的整数在这个循环群上乘法可逆，故不会使它消失。由有限交换群的主分量分解，\(A_{(p)}\) 与 \(p\)-主分量同构。也可直接对 \(|A|=p^r s\)、\(p\nmid s\)，用 Bézout 等式把每个元素拆为分别被 \(p^r\) 和 \(s\) 零化的两部分。若 \(q\ne p\)，非零群 \(\mathbb Z/q\mathbb Z\) 在 \(p\) 处局部化为零，说明须逐素数检测。
\end{solution}

\begin{exercise}\label{b5:ex:05-positive-restrictions}
在 \(S_3\) 中令 \(f=-1+\varepsilon+\chi\)，其中 \(\varepsilon\) 为符号特征标、\(\chi\) 为二维不可约特征标。证明 \(f\) 限制到每个初等子群都是实际特征标，但 \(f\) 自身不是实际特征标。由此说明 Brauer 判据为何只断言“虚特征标”。
\end{exercise}
\begin{solution}
\(S_3\) 不是初等群，其初等子群只有 \(1\)、二阶子群、三阶子群及其共轭。在 \(C_2\) 上，\(\varepsilon\) 为非平凡一维特征标 \(\eta\)，而 \(\chi|_{C_2}=1+\eta\)，故 \(f|_{C_2}=2\eta\)。在 \(C_3\) 上，\(\varepsilon|_{C_3}=1\)、\(\chi|_{C_3}=\lambda+\lambda^2\)，故 \(f|_{C_3}=\lambda+\lambda^2\)。平凡子群上 \(f\) 次数为 \(2\)。这些均为实际特征标，但 \(f\) 在平凡不可约特征标上的系数是 \(-1\)，故全群上不是实际特征标。
\end{solution}

\begin{exercise}\label{b5:ex:05-quaternion-monomial}
分类四元数群 \(Q_8\) 的二维不可约特征标 \(\theta\) 的全部一维诱导来源，即求出所有 \((H,\lambda)\)，使 \(H\le Q_8\)、\(\lambda\) 一维且 \(\operatorname{Ind}_H^{Q_8}\lambda=\theta\)。说明其他一维特征标在这些子群上诱导后如何分解。
\end{exercise}
\begin{solution}
比较次数得 \([Q_8:H]=2\)，故 \(|H|=4\)。这样的子群恰为 \(\langle i\rangle,\langle j\rangle,\langle k\rangle\)，都循环。先取 \(H=\langle i\rangle\)，陪集代表为 \(1,j\)，且 \(jij^{-1}=i^{-1}\)。诱导特征标在 \(H\) 上为 \(\lambda+\lambda^{-1}\)，在其外为零。若 \(\lambda(i)=\pm\sqrt{-1}\)，则在 \(1,-1\) 处取值 \(2,-2\)，在其余处为零，恰为 \(\theta\)。两个另外的子群同理，各有两个忠实一维特征标可用，故共有六组来源。

若 \(\lambda(i)=1\) 或 \(-1\)，它不是忠实特征标，并且共轭固定它。它在 \(Q_8\) 上恰有两个一维延拓：令 \(i\) 的像为所给的 \(\pm1\)，令 \(j\) 的像分别为 \(1,-1\)，这些赋值均满足群的关系。由 Frobenius 互反，两种延拓各出现一次，且次数和已经为 \(2\)，所以诱导恰为这两个一维特征标的直和，不再不可约。对 \(\langle j\rangle,\langle k\rangle\) 的结论由相同计算得到。
\end{solution}

\begin{exercise}\label{b5:ex:05-product-monomial}
令 \(E=Q_8\times C_3\)，\(\theta\) 为 \(Q_8\) 的二维不可约特征标，\(\nu\) 为 \(C_3\) 的非平凡一维特征标。写出 \(\theta\otimes\nu\) 的一维诱导表达及全部特征标值。
\end{exercise}
\begin{solution}
取 \(H=\langle i\rangle\)、\(\lambda(i)=\sqrt{-1}\)，由陪集基或\cref{b5:prop:elementary-monomial}的构造，
\[
\theta\otimes\nu
=\operatorname{Ind}_{H\times C_3}^{Q_8\times C_3}(\lambda\otimes\nu).
\]
对子群 \(H\times C_3\)，\(H\) 为 \(2\)-群且 \(C_3\) 为互素循环因子，仍是允许的 \(2\)-初等群。若 \(c\) 为 \(C_3\) 生成元、\(\nu(c)=\zeta_3\)，则特征标在 \((1,c^j)\) 处为 \(2\zeta_3^j\)，在 \((-1,c^j)\) 处为 \(-2\zeta_3^j\)，在其余元素处为零。
\end{solution}

\begin{exercise}\label{b5:ex:05-sign-negative-induction}
用 \(S_3\) 的二阶、三阶子群的一维特征标整诱导表达符号特征标 \(\varepsilon\)，并证明任何这样的表达都不可能只含非负系数。
\end{exercise}
\begin{solution}
令 \(\eta\) 为 \(C_2\) 的非平凡一维特征标，\(\lambda\) 为 \(C_3\) 的非平凡一维特征标。Frobenius 互反给出
\(\operatorname{Ind}_{C_2}^{S_3}\eta=\varepsilon+\chi\)，而
\(\operatorname{Ind}_{C_3}^{S_3}\lambda=\chi\)。故
\[
\varepsilon=\operatorname{Ind}_{C_2}^{S_3}\eta
-\operatorname{Ind}_{C_3}^{S_3}\lambda.
\]
\(S_3\) 的初等子群都是真子群，一维诱导的次数至少为 \(2\)，所以任何非零非负整数线性组合的次数至少为 \(2\)，不能等于 \(\varepsilon(1)=1\)。
\end{solution}

\begin{exercise}\label{b5:ex:05-nonintegral-average}
证明两个单位根的平均 \((1+i)/2\) 不是代数整数，并计算它的极小多项式及共轭乘积。
\end{exercise}
\begin{solution}
设 \(z=(1+i)/2\)，则 \(z^2-z+1/2=0\)。判别式为 \(-1\)，该首一二次多项式在 \(\mathbb Q[T]\) 中不可约，所以是极小多项式。其常数项为 \(1/2\)，不是整数；由\cref{b5:lem:algebraic-integer-arithmetic}，\(z\) 不可能是代数整数。两个共轭根的乘积正是 \((1+i)(1-i)/4=1/2\)，展示了平均值证明中整数乘积条件不可缺少。
\end{solution}

\begin{exercise}\label{b5:ex:05-a4-coprime-failure}
从 \(A_4\) 在四个字母上的置换表示中除去平凡表示，得次数为 \(3\) 的特征标 \(\chi\)。证明它不可约；再用双换位说明消失或标量作用引理中的互素假设不能删除。
\end{exercise}
\begin{solution}
在单位、三个双换位、八个三循环处，\(\chi\) 的值分别为 \(3,-1,0\)。于是
\(\langle\chi,\chi\rangle=(9+3)/12=1\)，故不可约。双换位在四维置换表示上的特征值为 \(1,1,-1,-1\)，去掉一个平凡方向后为 \(1,-1,-1\)。所以其迹为 \(-1\ne0\)，矩阵亦非标量。双换位的共轭类大小为 \(3\)，与次数 \(3\) 不互素，正落在引理假设之外。
\end{solution}

\begin{exercise}\label{b5:ex:05-dihedral-class-normal}
设 \(D_{14}=\langle r,s:r^7=s^2=1,\ srs=r^{-1}\rangle\)，下标表示群阶。求反射的共轭类，并计算由这个类中元素之商生成的正规子群。
\end{exercise}
\begin{solution}
由 \(r^a s r^{-a}=r^{2a}s\)，且 \(2\) 在模 \(7\) 下可逆，七个反射 \(r^j s\) 构成同一个共轭类 \(K\)，大小为 \(7\)。两个反射之商为
\((r^i s)(r^j s)^{-1}=r^{i-j}\)，遍历全部旋转。因此
\(\langle xy^{-1}:x,y\in K\rangle=\langle r\rangle\)，是阶为 \(7\) 的非平凡真正规子群。把旋转送到 \(1\)、反射送到 \(-1\) 的一维表示正好零化这个子群，与正文中的标量作用构造一致。
\end{solution}

\begin{exercise}\label{b5:ex:05-perfect-two-primes}
设有限群 \(G\) 满足 \(G'=G\)，且 \(|G|\) 至多含两个不同素因子。证明 \(G=1\)。再证明每个非平凡有限可解群都有非平凡一维复特征标。
\end{exercise}
\begin{solution}
由\cref{b5:thm:burnside-solvability}，\(G\) 可解。但 \(G'=G\) 使全部导出项都等于 \(G\)，只有 \(G=1\) 才能终止于 \(1\)。对任意非平凡有限可解群，导出群必须是真子群，所以交换化 \(G/G'\) 非平凡；取其一个素数阶商并选非平凡一维复特征标，提升到 \(G\) 就得到非平凡一维特征标。因此非平凡有限可解群并非只具有次数大于一的非平凡特征标。
\end{solution}

\begin{exercise}\label{b5:ex:05-a4-derived}
求 \(A_4\) 的导出列，并说明 \(p^a q^b\) 阶群的可解性不保证全部 Sylow 子群正规。
\end{exercise}
\begin{solution}
四元交换子群 \(V_4\) 由单位与三个双换位组成，正规且商 \(A_4/V_4\simeq C_3\)，故 \(A_4'\subset V_4\)。\(A_4\) 不交换，所以 \(A_4'\ne1\)；三个双换位在 \(A_4\) 中共轭，\(V_4\) 内任何包含其中一个的正规子群必须包含全部三个。因此 \(A_4'=V_4\)，再有 \(V_4'=1\)。导出列为 \(A_4,V_4,1\)。八个三循环组成四个不同三阶子群，故 Sylow \(3\)-子群不正规，尽管 \(|A_4|=2^2\cdot3\) 且群可解。
\end{solution}

\begin{exercise}\label{b5:ex:05-minimal-normal-linear}
设有限可解群 \(G\) 的极小非平凡正规子群为 \(N\simeq C_p^r\)。证明共轭给出不可约的 \(\mathbb F_p\)-线性作用 \(G\rightarrow\operatorname{GL}_r(\mathbb F_p)\)。在 \(A_4\) 的 \(V_4\) 上解释这一作用。
\end{exercise}
\begin{solution}
把 \(N\) 的群运算视作加法，整数次幂给出 \(\mathbb F_p\) 标量作用。群自同构保持两者，所以共轭为线性作用。任意 \(G\)-不变向量子空间是 \(N\) 的子群，并被全部 \(G\)-共轭保持，因而是 \(G\) 的正规子群。极小性使它只能为 \(0\) 或 \(N\)，故作用不可约。对 \(A_4\)，\(V_4\) 中三个非单位元是 \(\mathbb F_2^2\) 的三个非零向量，三循环通过共轭循环置换它们，没有保持某一条一维子空间；于是这个二维作用确为不可约。
\end{solution}

\begin{exercise}\label{b5:ex:05-real-quaternion-density}
令实四元数除环 \(\mathbb H\) 通过左乘作用于实向量空间 \(\mathbb H\)。证明这是不可约实矩阵子代数，但不等于 \(\operatorname{End}_{\mathbb R}(\mathbb H)\)。
\end{exercise}
\begin{solution}
若非零实子空间 \(W\) 被全部左乘保持，取 \(0\ne w\in W\)。对任意 \(y\in\mathbb H\)，因四元数可逆，\(y=(yw^{-1})w\in W\)，所以 \(W=\mathbb H\)。故左乘代数作用不可约。该子代数的实维数为 \(4\)，而全部自同态代数实维数为 \(4^2=16\)，两者不相等。这是底域未代数闭时的失败情形；除环的非标量右乘也解释了标量交换子结论为何不成立。
\end{solution}

\begin{exercise}\label{b5:ex:05-explicit-matrix-generation}
设 \(k\) 为含 \(n\) 个不同元素 \(a_1,\ldots,a_n\) 的域。令 \(D=\operatorname{diag}(a_1,\ldots,a_n)\)，令 \(S\) 循环置换标准基。直接证明由 \(D,S\) 生成的含幺 \(k\)-代数为 \(M_n(k)\)。
\end{exercise}
\begin{solution}
Lagrange 插值多项式
\[
f_i(T)=\prod_{j\ne i}\frac{T-a_j}{a_i-a_j}
\]
满足 \(f_i(D)=E_{ii}\)，所以所有坐标投影均在代数中。对任意 \(i,j\)，选 \(0\le r<n\) 使 \(S^r e_j=e_i\)，则 \(E_{ii}S^rE_{jj}=E_{ij}\)。全部矩阵单位都在所生成代数中，而它们是 \(M_n(k)\) 的基，故结论成立。这给出具体的矩阵生成计算，无需先调用矩阵代数定理。
\end{solution}

\begin{exercise}\label{b5:ex:05-necklace-count}
设 \(\ell\) 为素数，\(C_\ell\) 通过循环旋转作用在长为 \(\ell\) 的二进制串上。求轨道数，并解释它是某个置换表示中平凡表示的重数。
\end{exercise}
\begin{solution}
单位旋转固定全部 \(2^\ell\) 个串。每个非单位旋转在位置集合上是一个 \(\ell\)-循环，所以固定串必须所有位相同，恰有两种。因此轨道数为
\[
\frac{2^\ell+2(\ell-1)}{\ell}.
\]
置换特征标在每个旋转处等于其固定串数，和 \(1_{C_\ell}\) 的内积便是上述平均，即平凡表示的重数。这里没有使用可解性或不可约矩阵代数结论。
\end{solution}

\begin{exercise}\label{b5:ex:05-scalar-subgroup}
设 \(\rho:G\rightarrow\operatorname{GL}(V)\) 是有限群的非零复不可约表示，令 \(Z_\rho=\bigl\{g\in G:\rho(g)\text{ 为标量}\bigr\}\)。证明 \(Z_\rho\) 为正规子群，且 \(Z_\rho/\ker\rho=Z\bigl(G/\ker\rho\bigr)\)。忠实情形有何结论？
\end{exercise}
\begin{solution}
非零标量矩阵组成 \(\operatorname{GL}(V)\) 的中心子群，其原像 \(Z_\rho\) 因而正规。若 \(\rho(g)\) 为标量，它与每个 \(\rho(h)\) 交换，故 \(g\ker\rho\) 位于商群中心。反过来，若这个陪集居中，则 \(\rho(g)\) 与全部表示矩阵交换，由 Schur 引理是标量。两边相等，且 \(\ker\rho\subset Z_\rho\)，故得到所述商群等式。若表示忠实，则 \(Z_\rho=Z(G)\)；在非交换单群的非平凡不可约表示中，标量元素只能是单位。
\end{solution}

\begin{exercise}\label{b5:ex:05-simple-character-constraint}
设 \(G\) 为有限非交换单群，\(g\ne1\)。证明存在非平凡不可约复特征标 \(\chi\)，使 \(\gcd\bigl(\chi(1),|g^G|\bigr)>1\)。
\end{exercise}
\begin{solution}
反设全部非平凡不可约次数均与 \(|g^G|\) 互素。每个非平凡不可约表示的核为正规真子群，故由单性为 \(1\)。由\cref{b5:lem:coprime-class-scalar}，若 \(\chi(g)\ne0\)，则 \(g\) 在一个忠实不可约表示中作用为标量，从而位于 \(Z(G)=1\)，矛盾。因此所有非平凡不可约特征标在 \(g\) 处均为零。正则特征标分解遂在 \(g\ne1\) 处给出 \(0=1\)，矛盾。这一结论约束类大小与次数的公共素因子，而不只是分别要求它们整除群阶。
\end{solution}

